%<paper-block id="appendix.h001" type="heading">
\section*{附录A\quad 支撑材料文件列表}
%</paper-block>
%<paper-block id="appendix.p001" type="paragraph">
以下路径均相对于支撑材料根目录。源程序保持仓库中的真实文件名与目录结构，未重命名、未合并。
%</paper-block>
%<paper-block id="appendix.t001" type="table">
\begin{longtable}{p{0.58\textwidth}p{0.34\textwidth}}
\toprule
文件名 & 功能描述 \\
\midrule
\endfirsthead
\toprule
文件名 & 功能描述 \\
\midrule
\endhead
\texttt{02\_\allowbreak{}最终代码/\allowbreak{}第一二问/\allowbreak{}src/\allowbreak{}q12/\allowbreak{}\_\allowbreak{}\_\allowbreak{}init\_\allowbreak{}\_\allowbreak{}.py} & 第一、二问求解包初始化 \\
\texttt{02\_\allowbreak{}最终代码/\allowbreak{}第一二问/\allowbreak{}src/\allowbreak{}q12/\allowbreak{}geometry.py} & 问题一测向半平面交、区域直径与最小包围圆 \\
\texttt{02\_\allowbreak{}最终代码/\allowbreak{}第一二问/\allowbreak{}src/\allowbreak{}q12/\allowbreak{}posterior.py} & 问题一物理可行区域及内外包络 \\
\texttt{02\_\allowbreak{}最终代码/\allowbreak{}第一二问/\allowbreak{}src/\allowbreak{}q12/\allowbreak{}policy.py} & 问题二接收安全域与第二检测点策略 \\
\texttt{02\_\allowbreak{}最终代码/\allowbreak{}第一二问/\allowbreak{}src/\allowbreak{}q12/\allowbreak{}interval.py} & 问题二区间运算与有向误差控制 \\
\texttt{02\_\allowbreak{}最终代码/\allowbreak{}第一二问/\allowbreak{}src/\allowbreak{}q12/\allowbreak{}proof.py} & 问题二极小极大策略的计算机辅助核验 \\
\texttt{02\_\allowbreak{}最终代码/\allowbreak{}第三问/\allowbreak{}strategy.py} & 问题三动作、状态及基础定位接口 \\
\texttt{02\_\allowbreak{}最终代码/\allowbreak{}第三问/\allowbreak{}problem1\_\allowbreak{}v4\_\allowbreak{}inline.py} & 问题三内联交会定位几何模块 \\
\texttt{02\_\allowbreak{}最终代码/\allowbreak{}第三问/\allowbreak{}strategy\_\allowbreak{}v8.py} & 问题三逐频道覆盖、定位、清除与失败恢复 \\
\texttt{02\_\allowbreak{}最终代码/\allowbreak{}第三问/\allowbreak{}strategy\_\allowbreak{}fast.py} & 问题三正式链路使用的搜索参数基类 \\
\texttt{02\_\allowbreak{}最终代码/\allowbreak{}第三问/\allowbreak{}strategy\_\allowbreak{}night.py} & 问题三连续区域覆盖证书与搜索基类 \\
\texttt{02\_\allowbreak{}最终代码/\allowbreak{}第三问/\allowbreak{}strategy\_\allowbreak{}transit.py} & 问题三剩余区域更新与排空判定 \\
\texttt{02\_\allowbreak{}最终代码/\allowbreak{}第三问/\allowbreak{}strategy\_\allowbreak{}joint.py} & 问题三补测点成本评估与选择 \\
\texttt{02\_\allowbreak{}最终代码/\allowbreak{}第三问/\allowbreak{}strategy\_\allowbreak{}task.py} & 问题三最终 OpticalTaskP3：数量上界与有限光学覆盖 \\
\texttt{02\_\allowbreak{}最终代码/\allowbreak{}第四问/\allowbreak{}strategy.py} & 问题四动作、状态及基础定位接口 \\
\texttt{02\_\allowbreak{}最终代码/\allowbreak{}第四问/\allowbreak{}problem1\_\allowbreak{}v4\_\allowbreak{}inline.py} & 问题四内联交会定位几何模块 \\
\texttt{02\_\allowbreak{}最终代码/\allowbreak{}第四问/\allowbreak{}strategy\_\allowbreak{}p4.py} & 问题四定向源观测更新与基础搜索逻辑 \\
\texttt{02\_\allowbreak{}最终代码/\allowbreak{}第四问/\allowbreak{}strategy\_\allowbreak{}fast.py} & 问题四搜索与定位调度基类 \\
\texttt{02\_\allowbreak{}最终代码/\allowbreak{}第四问/\allowbreak{}strategy\_\allowbreak{}cost.py} & 问题四有限光学覆盖函数及成本计算 \\
\texttt{02\_\allowbreak{}最终代码/\allowbreak{}第四问/\allowbreak{}strategy\_\allowbreak{}finish.py} & 问题四终局光学覆盖与清除控制 \\
\texttt{02\_\allowbreak{}最终代码/\allowbreak{}第四问/\allowbreak{}discovery\_\allowbreak{}mesh.py} & 问题四连续覆盖核验所需几何网格辅助 \\
\texttt{02\_\allowbreak{}最终代码/\allowbreak{}第四问/\allowbreak{}discovery\_\allowbreak{}rings.py} & 问题四搜索基类加载的环形站点构造模块 \\
\texttt{02\_\allowbreak{}最终代码/\allowbreak{}第四问/\allowbreak{}discovery\_\allowbreak{}coverage.py} & 问题四未知朝向下的连续发现覆盖核验 \\
\texttt{02\_\allowbreak{}最终代码/\allowbreak{}第四问/\allowbreak{}discovery\_\allowbreak{}compact.py} & 问题四21站紧凑发现构造 \\
\texttt{02\_\allowbreak{}最终代码/\allowbreak{}第四问/\allowbreak{}discovery\_\allowbreak{}prune.py} & 问题四带覆盖证书的后续扫描站删减 \\
\texttt{02\_\allowbreak{}最终代码/\allowbreak{}第四问/\allowbreak{}local\_\allowbreak{}hunt.py} & 问题四有限局部补测规划 \\
\texttt{02\_\allowbreak{}最终代码/\allowbreak{}第四问/\allowbreak{}strategy\_\allowbreak{}hunt.py} & 问题四发现、局部定位与清除协同 \\
\texttt{02\_\allowbreak{}最终代码/\allowbreak{}第四问/\allowbreak{}strategy\_\allowbreak{}route.py} & 问题四扫描站与清除任务联合路径规划 \\
\texttt{02\_\allowbreak{}最终代码/\allowbreak{}第四问/\allowbreak{}strategy\_\allowbreak{}route\_\allowbreak{}probe.py} & 问题四最终 RouteProbeP4：零绕路补测与调度 \\
\texttt{02\_\allowbreak{}最终代码/\allowbreak{}第一二问/\allowbreak{}requirements.txt} & 第一、二问运行依赖 \\
\texttt{02\_\allowbreak{}最终代码/\allowbreak{}第一二问/\allowbreak{}examples/\allowbreak{}q1\_\allowbreak{}cases.json} & 问题一典型算例输入 \\
\texttt{02\_\allowbreak{}最终代码/\allowbreak{}第一二问/\allowbreak{}examples/\allowbreak{}q1\_\allowbreak{}observations.json} & 问题一观测示例输入 \\
\texttt{02\_\allowbreak{}最终代码/\allowbreak{}第一二问/\allowbreak{}proofs/\allowbreak{}global\_\allowbreak{}certificate.json} & 问题二全局证明证书 \\
\texttt{02\_\allowbreak{}最终代码/\allowbreak{}第一二问/\allowbreak{}results/\allowbreak{}q1\_\allowbreak{}samples.csv} & 问题一数值实验逐场结果 \\
\texttt{02\_\allowbreak{}最终代码/\allowbreak{}第一二问/\allowbreak{}results/\allowbreak{}q1\_\allowbreak{}summary.json} & 问题一数值实验汇总 \\
\texttt{02\_\allowbreak{}最终代码/\allowbreak{}第一二问/\allowbreak{}results/\allowbreak{}q2\_\allowbreak{}paired\_\allowbreak{}samples.csv} & 问题二配对数值实验逐场结果 \\
\texttt{02\_\allowbreak{}最终代码/\allowbreak{}第一二问/\allowbreak{}results/\allowbreak{}q2\_\allowbreak{}summary.json} & 问题二数值实验汇总 \\
\texttt{02\_\allowbreak{}最终代码/\allowbreak{}第三问/\allowbreak{}requirements\_\allowbreak{}collab.txt} & 问题三运行依赖 \\
\texttt{02\_\allowbreak{}最终代码/\allowbreak{}第四问/\allowbreak{}requirements-route.txt} & 问题四最终路线策略运行依赖 \\
\texttt{04\_\allowbreak{}结果汇总/\allowbreak{}第三四问/\allowbreak{}第三问.md} & 问题三结果汇总说明 \\
\texttt{04\_\allowbreak{}结果汇总/\allowbreak{}第三四问/\allowbreak{}第三问.json} & 问题三脱敏逐场结果数据 \\
\texttt{04\_\allowbreak{}结果汇总/\allowbreak{}第三四问/\allowbreak{}第四问.md} & 问题四结果汇总说明 \\
\texttt{04\_\allowbreak{}结果汇总/\allowbreak{}第三四问/\allowbreak{}第四问.json} & 问题四脱敏逐场结果数据 \\
\texttt{04\_\allowbreak{}结果汇总/\allowbreak{}第三问正式/\allowbreak{}P3\_\allowbreak{}FORMAL\_\allowbreak{}RESULTS.md} & 问题三三次正式测试汇总 \\
\texttt{04\_\allowbreak{}结果汇总/\allowbreak{}第三问正式/\allowbreak{}evidence/\allowbreak{}p3\_\allowbreak{}formal\_\allowbreak{}summary.json} & 问题三正式测试机器可读证据 \\
\texttt{04\_\allowbreak{}结果汇总/\allowbreak{}正式运行说明.md} & 第三、四问正式运行与结果口径说明 \\
\texttt{figures/\allowbreak{}fig\_\allowbreak{}q1\_\allowbreak{}measurement\_\allowbreak{}geometry.png} & 单次误差楔形与测向交会图 \\
\texttt{figures/\allowbreak{}fig\_\allowbreak{}q1\_\allowbreak{}diameter\_\allowbreak{}cover\_\allowbreak{}jung.png} & 直径圆与Jung界对照图 \\
\texttt{figures/\allowbreak{}fig\_\allowbreak{}q1\_\allowbreak{}statistics.png} & 第一问共享场景统计图 \\
\texttt{figures/\allowbreak{}fig\_\allowbreak{}q2\_\allowbreak{}information\_\allowbreak{}safe\_\allowbreak{}candidate.png} & 第二问首次信息、安全域与候选区图 \\
\texttt{figures/\allowbreak{}fig\_\allowbreak{}q2\_\allowbreak{}global\_\allowbreak{}optimality.png} & 第二问下界、临界根与最优性图 \\
\texttt{figures/\allowbreak{}fig\_\allowbreak{}q2\_\allowbreak{}strategy\_\allowbreak{}benchmark.png} & 第二问配对策略比较图 \\
\texttt{figures/\allowbreak{}fig\_\allowbreak{}app\_\allowbreak{}q1\_\allowbreak{}two\_\allowbreak{}station\_\allowbreak{}detail.png} & 两次测向反例局部放大图 \\
\texttt{figures/\allowbreak{}fig\_\allowbreak{}app\_\allowbreak{}q1\_\allowbreak{}disk\_\allowbreak{}bracket.png} & 圆盘近似与直径夹逼补充图 \\
\texttt{figures/\allowbreak{}fig\_\allowbreak{}app\_\allowbreak{}q1\_\allowbreak{}error\_\allowbreak{}sensitivity.png} & 固定配置的误差半宽灵敏度图 \\
\texttt{figures/\allowbreak{}fig\_\allowbreak{}app\_\allowbreak{}q2\_\allowbreak{}continuous\_\allowbreak{}certificate.png} & 非活动读数区间的连续核验图 \\
\texttt{paper\_figures.pdf} & 本文全部图件的高分辨率与矢量合并版 \\
\bottomrule
\end{longtable}
%</paper-block>
\clearpage
\linespread{1.3}\selectfont
%<paper-block id="appendix.h002" type="heading">
\section*{附录B\quad 完整源程序代码}
%</paper-block>
%<paper-block id="appendix.p002" type="paragraph">
本附录仅收录最终模型实际依赖的自研源码闭包。代码按仓库原文完整收录，不改文件名、不合并文件、不删改代码。
Python 标准库以及 NumPy、SciPy、Shapely 等第三方软件包不属于本队源程序，故不附其源码。
第三问以 \texttt{OpticalTaskP3} 为最终策略入口，第四问以 \texttt{RouteProbeP4} 为最终策略入口。
%</paper-block>
%<paper-block id="appendix.h003" type="heading">
\subsection*{B.1 第一、二问模型求解与证明程序}
%</paper-block>
%<paper-block id="appendix.p003" type="paragraph">
\noindent\textbf{文件：}\texttt{02\_\allowbreak{}最终代码/\allowbreak{}第一二问/\allowbreak{}src/\allowbreak{}q12/\allowbreak{}\_\allowbreak{}\_\allowbreak{}init\_\allowbreak{}\_\allowbreak{}.py}\\
\textbf{功能：}第一、二问求解包初始化
%</paper-block>
%<paper-block id="appendix.code001" type="code">
\begin{lstlisting}[style=appendixcode]
"""Reproducible Q1/Q2 algorithms and complete minimax proof verifier."""
\end{lstlisting}
%</paper-block>

%<paper-block id="appendix.p004" type="paragraph">
\noindent\textbf{文件：}\texttt{02\_\allowbreak{}最终代码/\allowbreak{}第一二问/\allowbreak{}src/\allowbreak{}q12/\allowbreak{}geometry.py}\\
\textbf{功能：}问题一测向半平面交、区域直径与最小包围圆
%</paper-block>
%<paper-block id="appendix.code002" type="code">
\begin{lstlisting}[style=appendixcode]
"""Validated planar set-membership geometry for CUMCM 2026 B, Q1/Q2.

P is the pure bearing intersection (possibly unbounded). K is a separately
labelled convex relaxation using the target disk and reception upper bound.
An ordinary direction excludes a 5 m disk; that nonconvex exclusion is NOT
silently claimed to be part of K. Empty, point, segment, bounded and unbounded
states are represented separately. All public angles are in degrees.
"""
from __future__ import annotations
from dataclasses import dataclass
from functools import lru_cache
import math
from typing import Sequence
import numpy as np
from scipy.optimize import linprog

try:
    from numba import njit
except ImportError:  # Correct, slower fallback; no network or runtime installation.
    def njit(*args, **kwargs):
        return lambda f: f

ABS_TOL = 2e-9

@dataclass(frozen=True)
class Observation:
    x: float
    y: float
    bearing_deg: float
    def __post_init__(self):
        if not np.isfinite([self.x, self.y, self.bearing_deg]).all():
            raise ValueError('Observation values must be finite.')
        if max(abs(self.x), abs(self.y)) > 2_000_000:
            raise ValueError('Position exceeds the interface coordinate limit.')
    @property
    def position(self):
        return np.array([self.x, self.y], dtype=float)

@dataclass
class Region:
    status: str
    vertices: np.ndarray
    description: str
    def metrics(self) -> dict:
        if self.status == 'empty':
            return dict(status='empty', D=None, radius=None, q=None,
                        diameter_circle_covers=None, vertices=[])
        if self.status == 'unbounded':
            return dict(status='unbounded', D=None, radius=None, q=None,
                        diameter_circle_covers=False, vertices=self.vertices.tolist())
        p = self.vertices
        d, ia, ib = diameter(p)
        c, r = minimum_enclosing_circle(p)
        mid = (p[ia] + p[ib]) / 2
        off = float(np.max(np.linalg.norm(p-mid, axis=1)) - d/2)
        return dict(status=self.status, D=d, radius=r, q=(2*r/d if d > 0 else None),
                    mec_center=c.tolist(), A=p[ia].tolist(), B=p[ib].tolist(),
                    diameter_circle_center=mid.tolist(), off=off,
                    diameter_circle_covers=off <= 1e-7 * max(1., d/1000),
                    area=polygon_area(p), vertices=p.tolist())


def wrap_deg(a):
    return (np.asarray(a) + 180.) % 360. - 180.


def validate_eps(eps):
    if not np.isfinite(eps) or not 0 < eps < 90:
        raise ValueError('eps must lie strictly between 0 and 90 degrees.')


def bearing_halfplanes(obs: Observation, eps_deg=1.):
    validate_eps(eps_deg)
    lo, hi = np.deg2rad([obs.bearing_deg-eps_deg, obs.bearing_deg+eps_deg])
    a = np.array([[np.sin(lo), -np.cos(lo)], [-np.sin(hi), np.cos(hi)]])
    return a, a @ obs.position


@njit(cache=True)
def clip_raw(poly, nx, ny, b):
    """Clip a convex point sequence, retaining point/segment degeneracies.

    The half-plane has a tiny outward floating-point safety displacement;
    there is no geometric vertex simplification that could cut off a source.
    """
    m = len(poly)
    if m == 0:
        return np.empty((0, 2), dtype=np.float64)
    bound = b + ABS_TOL
    out = np.empty((m+2, 2), dtype=np.float64)
    count = 0
    for i in range(m):
        j = (i-1) % m
        va = nx*poly[j,0]+ny*poly[j,1]-bound
        vb = nx*poly[i,0]+ny*poly[i,1]-bound
        ai, bi = va <= 0., vb <= 0.
        if ai != bi:
            t = va/(va-vb)
            out[count,0] = poly[j,0]+t*(poly[i,0]-poly[j,0])
            out[count,1] = poly[j,1]+t*(poly[i,1]-poly[j,1])
            count += 1
        if bi:
            out[count] = poly[i]
            count += 1
    return out[:count].copy()


def clip(poly, a, b):
    p = np.asarray(poly, dtype=float).reshape(-1,2)
    for n, d in zip(a,b):
        p = clip_raw(p, float(n[0]), float(n[1]), float(d))
        if len(p)==0: break
    return p


def convex_hull(points):
    p = np.asarray(points, dtype=float).reshape(-1,2)
    if len(p) == 0: return p
    pts = sorted(set(map(tuple, p)))
    if len(pts) <= 2: return np.asarray(pts)
    def cross(o,a,b):
        return (a[0]-o[0])*(b[1]-o[1])-(a[1]-o[1])*(b[0]-o[0])
    lower=[]
    for x in pts:
        while len(lower)>=2 and cross(lower[-2], lower[-1], x)<=0: lower.pop()
        lower.append(x)
    upper=[]
    for x in reversed(pts):
        while len(upper)>=2 and cross(upper[-2], upper[-1], x)<=0: upper.pop()
        upper.append(x)
    return np.asarray(lower[:-1]+upper[:-1], dtype=float)


def pure_bearing_region(observations: Sequence[Observation], eps_deg=1.) -> Region:
    """Exact half-plane model; vectorized O(h^3) vertex feasibility, h=2n.

    No bounding box is used to turn an unbounded wedge into a finite polygon.
    A feasibility LP handles empty/degenerate inputs; a recession-cone test
    based on normal-angle gaps distinguishes bounded from unbounded sets.
    """
    validate_eps(eps_deg)
    if not observations: raise ValueError('At least one direction observation is required.')
    aa,bb=zip(*(bearing_halfplanes(o,eps_deg) for o in observations))
    a,b=np.vstack(aa),np.concatenate(bb)
    # Translation improves conditioning for off-origin coordinates.
    origin=np.mean([o.position for o in observations],axis=0)
    bs=b-a@origin
    feas=linprog([0.,0.],A_ub=a,b_ub=bs,bounds=[(None,None)]*2,method='highs')
    if feas.status==2:
        return Region('empty',np.empty((0,2)),'pure bearing intersection')
    if not feas.success:
        raise ArithmeticError(f'Feasibility LP failed: {feas.message}')
    ang=np.sort(np.mod(np.arctan2(a[:,1],a[:,0]),2*np.pi))
    gap=np.max(np.diff(np.r_[ang,ang[0]+2*np.pi]))
    if gap >= np.pi-1e-12:
        return Region('unbounded',np.empty((0,2)),'pure bearing intersection')
    i,j=np.triu_indices(len(b),k=1)
    det=a[i,0]*a[j,1]-a[i,1]*a[j,0]
    ok=np.abs(det)>1e-13
    i,j,det=i[ok],j[ok],det[ok]
    pts=np.column_stack(((bs[i]*a[j,1]-a[i,1]*bs[j])/det,
                         (a[i,0]*bs[j]-bs[i]*a[j,0])/det))
    pts=pts[np.all(pts@a.T <= bs+2e-7,axis=1)]
    if len(pts)==0:
        # Bounded feasible set without reliable vertices signals conditioning,
        # not perfect localization. Do not fabricate a zero diameter.
        raise ArithmeticError('Numerically unresolved bounded intersection.')
    p=convex_hull(pts+origin)
    return Region('point' if len(p)==1 else 'segment' if len(p)==2 else 'bounded',p,
                  'pure bearing intersection')


@lru_cache(maxsize=64)
def normals(n: int, phase: float=0.):
    if n<8: raise ValueError('Use at least 8 polygon sides.')
    a=phase+2*np.pi*np.arange(n)/n
    v=np.column_stack((np.cos(a),np.sin(a)))
    v.setflags(write=False)
    return v


def disk_polygon(center, radius, n=512, outer=True, phase=0.):
    if radius<=0 or not np.isfinite(radius): raise ValueError('radius must be positive.')
    r=radius/np.cos(np.pi/n) if outer else radius
    # Both inner and outer polygons use the same outward normal grid.
    return np.asarray(center)+r*normals(n,phase+np.pi/n)


def clip_disk(poly, center, radius, n=512, outer=True, phase=0.):
    p=np.asarray(poly,dtype=float).reshape(-1,2)
    if len(p)==0: return p
    ns=normals(n,phase)
    rr=radius if outer else radius*np.cos(np.pi/n)
    b=ns@np.asarray(center)+rr
    active=np.max(p@ns.T-b,axis=0)>0.
    # Omitting already satisfied planes is exact for a shrinking convex set.
    return clip(p,ns[active],b[active])


def relaxed_region(observations: Sequence[Observation], eps_deg=1., n=512,
                   outer=True, target_radius=1800., reception_upper=1500.) -> Region:
    """K: convex relaxation of direction data + known radius upper bounds.

    K does not enforce distance>5. For inner=True (outer=False), floating point
    slack is offset inward so the intended inner geometry remains conservative.
    """
    validate_eps(eps_deg)
    if not observations: raise ValueError('Need direction observations.')
    if n<8: raise ValueError('n must be >=8.')
    # Align polygon orientation with the first observation for equivariance.
    phase=math.radians(observations[0].bearing_deg)
    safety=1e-7 if not outer else 0.
    p=disk_polygon((0.,0.),target_radius-safety,n,outer,phase)
    for o in observations:
        a,b=bearing_halfplanes(o,eps_deg)
        p=clip(p,a,b-(1e-7 if not outer else 0.))
        if len(p)==0: return Region('empty',p,'convex relaxation K')
    for o in observations:
        p=clip_disk(p,o.position,reception_upper-safety,n,outer,phase)
        if len(p)==0: return Region('empty',p,'convex relaxation K')
    # Remove only exact duplicate/collinear interior points by convex hull.
    p=convex_hull(p)
    return Region('point' if len(p)==1 else 'segment' if len(p)==2 else 'bounded',p,
                  ('outer' if outer else 'inner')+' polygon of convex relaxation K')


@njit(cache=True)
def diameter_raw(poly):
    # Quadratic reference, also faster than calipers for tiny posteriors.
    best=0.; bi=0; bj=0
    for i in range(len(poly)):
        for j in range(i+1,len(poly)):
            dx=poly[i,0]-poly[j,0]; dy=poly[i,1]-poly[j,1]
            ds=dx*dx+dy*dy
            if ds>best: best=ds; bi=i; bj=j
    return math.sqrt(best),bi,bj


def diameter(points):
    p=convex_hull(points)
    m=len(p)
    if m==0: raise ValueError('The empty set is not a zero-diameter localization.')
    if m<24:
        # Return indices into the original input, not the sorted hull.
        d,i,j=diameter_raw(p)
        original=np.asarray(points)
        return float(d),int(np.argmin(np.sum((original-p[i])**2,axis=1))),int(np.argmin(np.sum((original-p[j])**2,axis=1)))
    j=1; best=-1.; pair=(0,1)
    def area(i,ni,k):
        a=p[ni]-p[i];b=p[k]-p[i]
        return abs(a[0]*b[1]-a[1]*b[0])
    for i in range(m):
        ni=(i+1)%m
        for _ in range(m):
            nj=(j+1)%m
            if area(i,ni,nj)>area(i,ni,j)+1e-12: j=nj
            else: break
        for a in (i,ni):
            for b in (j,(j+1)%m):  # Include ties at parallel supporting edges.
                ds=float(np.sum((p[a]-p[b])**2))
                if ds>best: best=ds; pair=(a,b)
    original=np.asarray(points)
    return math.sqrt(best),*(int(np.argmin(np.sum((original-p[k])**2,axis=1))) for k in pair)


def polygon_area(p):
    if len(p)<3:return 0.
    q=np.asarray(p)-p[0]
    return abs(float(np.sum(q[:,0]*np.roll(q[:,1],-1)-q[:,1]*np.roll(q[:,0],-1))))/2


def minimum_enclosing_circle(points,seed=17):
    p=np.asarray(points,dtype=float).reshape(-1,2)
    if len(p)==0:raise ValueError('No enclosing circle requested for empty input.')
    offset=p.mean(axis=0); scale=max(1.,float(np.ptp(p,axis=0).max()))
    p=(p-offset)/scale
    p=p[np.random.default_rng(seed).permutation(len(p))]
    def pair(a,b):
        c=(a+b)/2;return c,float(np.linalg.norm(a-c))
    def inside(x,c):return np.linalg.norm(x-c[0])<=c[1]+2e-12
    def cross(a,b):return a[0]*b[1]-a[1]*b[0]
    def circum(a,b,c):
        u=b-a;v=c-a;det=2*cross(u,v)
        if abs(det)<1e-15:return None
        ux=(np.dot(u,u)*v[1]-np.dot(v,v)*u[1])/det
        uy=(u[0]*np.dot(v,v)-v[0]*np.dot(u,u))/det
        center=a+np.array([ux,uy]);return center,float(np.linalg.norm(center-a))
    def two_boundary(prefix,a,b):
        base=pair(a,b);left=None;right=None
        for x in prefix:
            if inside(x,base):continue
            side=cross(b-a,x-a);c=circum(a,b,x)
            if c is None:continue
            cc=cross(b-a,c[0]-a)
            if side>0 and (left is None or cc>cross(b-a,left[0]-a)):left=c
            if side<0 and (right is None or cc<cross(b-a,right[0]-a)):right=c
        if left is None:return right or base
        if right is None:return left
        return left if left[1]<=right[1] else right
    circle=None
    for i,a in enumerate(p):
        if circle is not None and inside(a,circle):continue
        circle=(a.copy(),0.)
        for j,b in enumerate(p[:i]):
            if inside(b,circle):continue
            circle=pair(a,b) if circle[1]==0 else two_boundary(p[:j+1],a,b)
    center=circle[0]*scale+offset
    # Repair only floating-point last bits; this enforces containment explicitly.
    radius=float(np.max(np.linalg.norm(np.asarray(points)-center,axis=1)))
    return center,radius


def contains(poly,point,tol=1e-6):
    p=np.asarray(poly);x=np.asarray(point)
    if len(p)==0:return False
    if len(p)==1:return np.linalg.norm(p[0]-x)<=tol
    if len(p)==2:
        d=p[1]-p[0];t=np.clip(np.dot(x-p[0],d)/max(np.dot(d,d),1e-30),0,1)
        return np.linalg.norm(p[0]+t*d-x)<=tol
    edges=np.roll(p,-1,axis=0)-p
    v=x-p
    return bool(np.all(edges[:,0]*v[:,1]-edges[:,1]*v[:,0]>=-tol*np.maximum(1.,np.linalg.norm(edges,axis=1))))
\end{lstlisting}
%</paper-block>

%<paper-block id="appendix.p005" type="paragraph">
\noindent\textbf{文件：}\texttt{02\_\allowbreak{}最终代码/\allowbreak{}第一二问/\allowbreak{}src/\allowbreak{}q12/\allowbreak{}posterior.py}\\
\textbf{功能：}问题一物理可行区域及内外包络
%</paper-block>
%<paper-block id="appendix.code003" type="code">
\begin{lstlisting}[style=appendixcode]
"""Full physical posterior: target/reception disks AND exclusions of near disks.

Polygonal lower/upper sets are used only to bracket the SAME physical set.
They are not silently substituted as a different optimization objective.
"""
from __future__ import annotations
import math
import numpy as np
from shapely.geometry import Polygon,LineString,Point
from .geometry import Observation,relaxed_region,disk_polygon,convex_hull,Region

def _shape(vertices):
    if len(vertices)==0:return Polygon()
    if len(vertices)==1:return Point(vertices[0])
    if len(vertices)==2:return LineString(vertices)
    return Polygon(vertices)

def _coords(shape):
    if shape.is_empty:return []
    if shape.geom_type=='Polygon':return list(shape.exterior.coords)[:-1]
    if shape.geom_type in ('LineString','LinearRing','Point'):return list(shape.coords)
    return [p for g in shape.geoms for p in _coords(g)]

def physical_posterior(observations,n=512,outer=True):
    observations=list(observations)
    if not observations:raise ValueError('At least one observed direction is required')
    base=relaxed_region(observations,n=n,outer=outer)
    shape=_shape(base.vertices)
    for obs in observations:
        # Upper F removes INSCRIBED near polygons; lower F removes CIRCUMSCRIBED ones.
        near=_shape(disk_polygon(obs.position,5.,n=n,outer=not outer))
        shape=shape.difference(near)
    points=_coords(shape)
    hull=convex_hull(points) if points else np.empty((0,2))
    status='empty' if len(hull)==0 else 'point' if len(hull)==1 else 'segment' if len(hull)==2 else 'bounded'
    return shape,Region(status,hull,('outer' if outer else 'inner')+' physical posterior (diameter computed on its hull)')

def posterior_bracket(observations,n=512):
    _,a=physical_posterior(observations,n,False);shape,b=physical_posterior(observations,n,True)
    ma,mb=a.metrics(),b.metrics()
    return dict(lower_m=ma['D'],upper_m=mb['D'],outer_metrics=mb,inner_metrics=ma),shape
\end{lstlisting}
%</paper-block>

%<paper-block id="appendix.p006" type="paragraph">
\noindent\textbf{文件：}\texttt{02\_\allowbreak{}最终代码/\allowbreak{}第一二问/\allowbreak{}src/\allowbreak{}q12/\allowbreak{}policy.py}\\
\textbf{功能：}问题二接收安全域与第二检测点策略
%</paper-block>
%<paper-block id="appendix.code004" type="code">
\begin{lstlisting}[style=appendixcode]
"""Constant-time universal minimax second-station policy.

For an untruncated first sector the two ideal stations are the EXACT global
minimizers over the entire reception-safe plane. For any valid first observation,
including outside the target disk, their worst residual diameter is <=D_star.
The policy is therefore minimax-optimal over all first-observation states.
Target-boundary truncation may permit an even better state-specific policy;
that different claim is not made by this module.
"""
from __future__ import annotations
from functools import lru_cache
from dataclasses import dataclass
import math
import numpy as np
from .geometry import Observation,relaxed_region,minimum_enclosing_circle

@dataclass(frozen=True)
class PolicyConstants:
    t_star:float
    diameter_star:float
    lower_x:float
    lower_y:float
    local_x:float
    local_y:float

@lru_cache(maxsize=1)
def policy_constants()->PolicyConstants:
    # These constants implement the analytic quadratic, not a stored grid result.
    a=5.;R=1000.;b=1500.;eps=math.pi/180;d=2*eps
    sd,cd,s2,c2=math.sin(d),math.cos(d),math.sin(2*d),math.cos(2*d)
    K=sd*(R*R-a*(b-a));H=(b-a)*s2
    qa=R*R-H*H
    qb=2*R*R*(sd*(b-2*a)*s2-b*cd*c2)-2*K*H
    qc=R*R*((sd*(b-2*a))**2+(b*cd)**2)-K*K
    disc=qb*qb-4*qa*qc
    if disc<=0:raise ArithmeticError('Unexpected quadratic discriminant')
    t=(-qb-math.sqrt(disc))/(2*qa)
    A=sd*(b-2*a)+t*s2;B=b*cd-t*c2;den=math.hypot(A,B)
    x=a+R*A/den;y=R*B/den
    lx=x*math.cos(eps)+y*math.sin(eps);ly=-x*math.sin(eps)+y*math.cos(eps)
    D=math.hypot(t-b*cd,b*sd)
    return PolicyConstants(t,D,x,y,lx,ly)

def safe_centers_local()->np.ndarray:
    e=math.pi/180
    return np.array([[r*math.cos(a),r*math.sin(a)] for r in (5.,1000.) for a in (-e,e)])

def reception_safe_local(points,tol=0.):
    """Return the exact four-circle reception-safe predicate for standard F0.

    For the standard first sector F0 that is not truncated by the target
    region, this four-circle intersection is the necessary-and-sufficient
    complete reception-safe set.  For an actual target-truncated F1 subset of
    F0, it remains a general sufficient safe domain but need not be the
    complete state-specific safe set.  This is an analytic four-circle
    predicate, not an empirical check of four sampled points.
    """
    pts=np.atleast_2d(np.asarray(points,dtype=float))
    if pts.shape[1]!=2 or not np.isfinite(pts).all():raise ValueError('Need finite 2D points')
    return np.all(np.sum((pts[:,None,:]-safe_centers_local()[None,:,:])**2,axis=2)<=(1000.+tol)**2,axis=1)

def transform(points,s1,bearing_deg):
    a=math.radians(float(bearing_deg));c,s=math.cos(a),math.sin(a)
    return np.asarray(points)@np.array([[c,s],[-s,c]])+np.asarray(s1)

def first_direction_possible(s1,bearing_deg):
    """Exact ray/disk feasibility criterion; only three possible angular extrema."""
    p=np.asarray(s1,float);angles=[bearing_deg-1.,bearing_deg+1.]
    if np.linalg.norm(p)>0:
        inward=math.degrees(math.atan2(-p[1],-p[0]));difference=(inward-bearing_deg+180)%360-180
        if abs(difference)<=1:angles.append(inward)
    for a in angles:
        u=np.array([math.cos(math.radians(a)),math.sin(math.radians(a))]);z=float(p@u)
        disc=z*z-float(p@p)+1800.**2
        if disc<0:continue
        root=math.sqrt(max(0.,disc));lo=max(0.,-z-root);hi=min(1500.,-z+root)
        if hi>5. and hi>=lo:return True
    return False

def choose_second_point(s1,bearing_deg,side=1,operational_margin=True)->dict:
    """Only observable inputs. side chooses one of the reflection-equivalent optima.

    operational_margin adds a 0.000056 m-scale inward displacement to prevent
    floating rounding of an active reception circle from causing a boundary miss.
    The ideal point and executed point are both returned and never conflated.
    """
    obs=Observation(float(s1[0]),float(s1[1]),float(bearing_deg))
    if side not in (-1,1):raise ValueError('side must be +1 or -1')
    # A direction cannot originate beyond the 1500m reception limit from Omega.
    if not first_direction_possible(obs.position,obs.bearing_deg):raise ValueError('No source can produce the stated first direction')
    k=policy_constants();ideal=np.array([k.local_x,side*k.local_y])
    shrink=1e-7 if operational_margin else 0.
    local=(1-shrink)*ideal+shrink*np.array([750.,0.])
    actual=transform(local,obs.position,obs.bearing_deg)
    return dict(S1=obs.position.tolist(),bearing_deg=obs.bearing_deg,side=side,
                theoretical_local=ideal.tolist(),operational_local=local.tolist(),S2=actual.tolist(),
                theoretical_S2=transform(ideal,obs.position,obs.bearing_deg).tolist(),
                theoretical_worst_diameter_m=k.diameter_star,travel_m=float(np.linalg.norm(local)),
                guard_displacement_m=float(np.linalg.norm(local-ideal)),
                minimum_reception_disk_margin_m=float(1000-np.max(np.linalg.norm(local-safe_centers_local(),axis=1))),
                guarantee='universal minimax; exact theorem for ideal coordinates; inward floating-point guard')

def good_candidate_mask(points):
    pts=np.atleast_2d(np.asarray(points,float));k=policy_constants()
    in_box=((np.abs(pts[:,0]-k.local_x)<=4)&(np.abs(np.abs(pts[:,1])-k.local_y)<=4))
    return reception_safe_local(pts)&in_box
\end{lstlisting}
%</paper-block>

%<paper-block id="appendix.p007" type="paragraph">
\noindent\textbf{文件：}\texttt{02\_\allowbreak{}最终代码/\allowbreak{}第一二问/\allowbreak{}src/\allowbreak{}q12/\allowbreak{}interval.py}\\
\textbf{功能：}问题二区间运算与有向误差控制
%</paper-block>
%<paper-block id="appendix.code005" type="code">
\begin{lstlisting}[style=appendixcode]
"""Directed, fixed-point rational interval arithmetic (Python integers only).

Endpoints are integers / 10**40. Every multiplication, division and square
root rounds outwards. Trigonometric constants use Machin's formula and
Taylor polynomials with explicit remainder bounds. No binary floating point
enters the proof arithmetic.
"""
from __future__ import annotations
from fractions import Fraction
from dataclasses import dataclass
from math import isqrt, factorial
from functools import lru_cache
SCALE=10**40

def ceil_div(a: int,b: int)->int:
    if b<0:a,b=-a,-b
    return -((-a)//b)

@dataclass(frozen=True)
class IV:
    lo:int
    hi:int
    def __post_init__(self):
        if self.lo>self.hi:raise ValueError('Reversed interval')
    @staticmethod
    def of(x):
        if isinstance(x,IV):return x
        f=Fraction(str(x)) if isinstance(x,float) else Fraction(x)
        return IV(f.numerator*SCALE//f.denominator,ceil_div(f.numerator*SCALE,f.denominator))
    @staticmethod
    def hull(a,b):
        a,b=IV.of(a),IV.of(b);return IV(min(a.lo,b.lo),max(a.hi,b.hi))
    def __add__(self,x):
        x=IV.of(x);return IV(self.lo+x.lo,self.hi+x.hi)
    __radd__=__add__
    def __neg__(self):return IV(-self.hi,-self.lo)
    def __sub__(self,x):return self+-IV.of(x)
    def __rsub__(self,x):return IV.of(x)+-self
    def __mul__(self,x):
        x=IV.of(x);p=[self.lo*x.lo,self.lo*x.hi,self.hi*x.lo,self.hi*x.hi]
        return IV(min(p)//SCALE,ceil_div(max(p),SCALE))
    __rmul__=__mul__
    def __truediv__(self,x):
        x=IV.of(x)
        if x.lo<=0<=x.hi:raise ZeroDivisionError('Interval divisor contains zero')
        vals=[Fraction(a*SCALE,b) for a in (self.lo,self.hi) for b in (x.lo,x.hi)]
        l,h=min(vals),max(vals)
        return IV(l.numerator//l.denominator,ceil_div(h.numerator,h.denominator))
    def __rtruediv__(self,x):return IV.of(x)/self
    def sq(self):
        p=[self.lo*self.lo,self.hi*self.hi]
        return IV(0 if self.lo<=0<=self.hi else min(p)//SCALE,ceil_div(max(p),SCALE))
    def sqrt(self):
        if self.lo<0:raise ValueError('Negative square-root interval')
        a,b=isqrt(self.lo*SCALE),isqrt(self.hi*SCALE)
        return IV(a,b+(b*b<self.hi*SCALE))
    def __pow__(self,n):
        if n<0:return 1/(self**(-n))
        r=IV.of(1)
        for _ in range(n):r=r*self
        return r
    def maxabs(self):return Fraction(max(abs(self.lo),abs(self.hi)),SCALE)
    def endpoints(self):return Fraction(self.lo,SCALE),Fraction(self.hi,SCALE)
    def json(self):
        def fmt(v):
            sign='-' if v<0 else '';v=abs(v)
            return sign+str(v//SCALE)+'.'+str(v%SCALE).zfill(40)
        return [fmt(self.lo),fmt(self.hi)]
    def midpoint(self):return Fraction(self.lo+self.hi,2*SCALE)

def _atan_recip(n,terms=70):
    z=Fraction(1,n);s=Fraction(0)
    for k in range(terms):s+=(-1)**k*z**(2*k+1)/(2*k+1)
    term=(-1)**terms*z**(2*terms+1)/(2*terms+1)
    return IV.hull(s,s+term)

@lru_cache(maxsize=1)
def pi():
    # Alternating arctangent series and Machin: pi=16 atan(1/5)-4 atan(1/239).
    return 16*_atan_recip(5)-4*_atan_recip(239)

def sincos_rad(x:IV):
    """Valid for |x|<=2; proof callers use a reduced angle <=pi/2.
    Horner polynomial, then Lagrange remainder |x|**(2n)/(2n)!.
    """
    x=IV.of(x)
    if x.maxabs()>2:raise ValueError('Reduce the angle before Taylor evaluation')
    z=x.sq();n=35
    c=IV.of(Fraction((-1)**(n-1),factorial(2*n-2)))
    t=IV.of(Fraction((-1)**(n-1),factorial(2*n-1)))
    for k in range(n-2,-1,-1):
        c=c*z+Fraction((-1)**k,factorial(2*k))
        t=t*z+Fraction((-1)**k,factorial(2*k+1))
    sn=x*t
    # Same conservative order-68 Taylor remainder bounds both polynomials.
    rem=IV.of(x.maxabs()**(2*n-1)/factorial(2*n-1))
    return IV(sn.lo-rem.hi,sn.hi+rem.hi),IV(c.lo-rem.hi,c.hi+rem.hi)

@lru_cache(maxsize=10000)
def sincos_deg(deg:Fraction):
    d=Fraction(deg)%360;q=int(d//90);r=d-90*q
    sn,cs=sincos_rad(pi()*r/180)
    return [(sn,cs),(cs,-sn),(-sn,-cs),(-cs,sn)][q]

def sincos_range(lo,hi):
    lo,hi=Fraction(lo),Fraction(hi)
    if hi<lo:raise ValueError('Reversed angular interval')
    if hi-lo>=360:return IV.of(-1)+IV(0,2*SCALE),IV(-SCALE,SCALE)
    vals=[sincos_deg(lo),sincos_deg(hi)]
    for k in range(int(lo//90)-1,int(hi//90)+2):
        if lo<=90*k<=hi:vals.append(sincos_deg(Fraction(90*k)))
    return IV(min(v[0].lo for v in vals),max(v[0].hi for v in vals)),IV(min(v[1].lo for v in vals),max(v[1].hi for v in vals))

def dot(a,b):return a[0]*b[0]+a[1]*b[1]
def cross(a,b):return a[0]*b[1]-a[1]*b[0]
def vecsub(a,b):return [a[i]-b[i] for i in range(2)]
def norm2(a):return a[0].sq()+a[1].sq()
\end{lstlisting}
%</paper-block>

%<paper-block id="appendix.p008" type="paragraph">
\noindent\textbf{文件：}\texttt{02\_\allowbreak{}最终代码/\allowbreak{}第一二问/\allowbreak{}src/\allowbreak{}q12/\allowbreak{}proof.py}\\
\textbf{功能：}问题二极小极大策略的计算机辅助核验
%</paper-block>
%<paper-block id="appendix.code006" type="code">
\begin{lstlisting}[style=appendixcode]
"""Computer-assisted, whole-plane minimax certificate for the universal policy.

No sampled location is used to establish the global lower bound. It follows
from a quadratic-over-disk maximization and the exact reception domain.
The upper proof covers ALL headings, with one analytic active interval.
All verifier arithmetic is directed rational arithmetic (interval.py).
"""
from __future__ import annotations
import json,time
from fractions import Fraction as F
from pathlib import Path
from itertools import combinations
from .interval import IV,SCALE,sincos_deg,sincos_range,pi,dot,cross,norm2,vecsub

R0=5;RM=1000;RU=1500

def constants():
    se,ce=sincos_deg(F(1));sd,cd=sincos_deg(F(2));s2,c2=sincos_deg(F(4))
    K=sd*(RM**2-R0*(RU-R0));H=s2*(RU-R0)
    qa=RM**2-H.sq()
    qb=2*RM**2*(sd*(RU-2*R0)*s2-RU*cd*c2)-2*K*H
    qc=RM**2*(sd.sq()*(RU-2*R0)**2+RU**2*cd.sq())-K.sq()
    discriminant=qb.sq()-4*qa*qc
    T=(-qb-discriminant.sqrt())/(2*qa)
    A=sd*(RU-2*R0)+T*s2;B=RU*cd-T*c2
    den=(A.sq()+B.sq()).sqrt()
    sx=R0+RM*A/den;sy=RM*B/den
    # Rotate from lower-ray coordinates (first lower ray has global angle -1deg).
    s=[sx*ce+sy*se,-sx*se+sy*ce]
    far=[RU*ce,-RU*se]
    d2=RU**2+T.sq()-2*RU*T*cd
    D=d2.sqrt()
    return dict(se=se,ce=ce,sd=sd,cd=cd,s2=s2,c2=c2,K=K,H=H,qa=qa,qb=qb,qc=qc,
                T=T,A=A,B=B,den=den,sx=sx,sy=sy,s=s,far=far,D2=d2,D=D)

def assert_interval(cond,msg):
    if not cond:raise AssertionError(msg)

def lower_checks(k):
    # Algebraic identity: the chosen root gives sup_disk(N-T H)=0.
    # Check branch, positive discriminant, stationary point outside the disk,
    # exact four-disk feasibility, and the physical lower witness.
    T=k['T'];sd=k['sd'];cd=k['cd'];sn=k['s2'];cs=k['c2'];sx=k['sx'];sy=k['sy']
    assert_interval(T.lo>1401*SCALE and T.hi<1402*SCALE,'Wrong quadratic root')
    assert_interval((k['K']+k['H']*T).lo>0,'Extraneous squared-equation root')
    # v-c = (A,B)/(2 sin delta); outside the radius-RM disk.
    assert_interval((k['den']/(2*sd)).lo>RM*SCALE,'Disk maximum not on boundary')
    assert_interval((k['A']).lo>0 and k['B'].lo>0,'Wrong normal quadrant')
    # Circle centred at (5,0) is active exactly by construction. Other three are strict.
    s=[sx,sy]
    centers=[[IV.of(RM),IV.of(0)],[R0*cd,R0*sd],[RM*cd,RM*sd]]
    distances=[]
    for c in centers:
        d=norm2(vecsub(s,c));assert_interval(d.hi<RM**2*SCALE,'Other reception disk violated');distances.append(d.sqrt().json())
    # Exterior-above-wedge regime and physically valid active pair.
    assert_interval((sy*cd-sx*sd).lo>0,'Station is not above the upper ray')
    q=[T*cd,T*sd];p=[IV.of(RU),IV.of(0)]
    assert_interval(norm2(vecsub(p,s)).lo>25*SCALE,'Far witness is near station')
    assert_interval(norm2(vecsub(q,s)).lo>25*SCALE,'Upper-ray witness is near station')
    assert_interval(T.lo>5*SCALE and T.hi<RU*SCALE,'Witness outside first annular sector')
    assert_interval(T.hi<(RU*cd).lo,'Wrong branch of distance monotonicity')
    assert_interval(k['D'].hi<111*SCALE,'Fallback lower-bound comparisons need D<111')
    return dict(other_disk_distances_m=distances,root_branch='smaller real root; unsquared sign positive',
                active_circle='center 5*u_minus, radius 1000; equality by construction',pass_=True)

def line_intersection(k,gamma,lo,hi):
    """Range of G on first ray gamma intersecting second ray beta in [lo,hi]."""
    sn,cs=sincos_range(lo,hi);sg,cg=sincos_deg(F(gamma));u=[cg,sg];v=[cs,sn]
    den=cg*sn-sg*cs
    r=cross(k['s'],v)/den
    rp=(k['s'][1]*cg-k['s'][0]*sg)/(den.sq())
    return [r*cg,r*sg],r,rp

def active_checks(k):
    a,b=F('-42.1'),F('-42.0')
    A,ra,da=line_intersection(k,1,a-1,b-1)
    B,rb,db=line_intersection(k,-1,a-1,b-1)
    C,rc,dc=line_intersection(k,1,a+1,b+1)
    V,rv,dv=line_intersection(k,-1,a+1,b+1)
    p=k['far'];sh,ch=sincos_deg(F('0.5'));tangent=[IV.of(RU),-RU*sh/ch]
    # Event occurs strictly within active interval: far-lower intersection crosses R.
    _,rleft,_=line_intersection(k,-1,a+1,a+1)
    _,rright,_=line_intersection(k,-1,b+1,b+1)
    assert_interval(rleft.hi<RU*SCALE<rright.lo,'Active event not bracketed')
    for name,r in [('A',ra),('B',rb),('C',rc)]:
        assert_interval(r.lo>5*SCALE and r.hi<RU*SCALE,f'{name} crosses radial cap')
    assert_interval(ra.hi<rc.lo and rb.hi<rv.lo,'Active quadrilateral radial ordering failed')
    assert_interval(all(z.lo>0 for z in (da,db,dc,dv)),'Ray intersection not increasing')
    forward={}
    for name,gamma,l,h in [('A',1,a-1,b-1),('B',-1,a-1,b-1),('C',1,a+1,b+1),('V',-1,a+1,b+1)]:
        sn,cs=sincos_range(l,h);sg,cg=sincos_deg(F(gamma))
        coefficient=cross(k['s'],[cg,sg])/(cg*sn-sg*cs)
        assert_interval(coefficient.lo>0,f'{name} lies on a backward second ray')
        forward[name]=coefficient.json()
    # 1/2 derivative of |A-V|^2 in radians, valid across the WHOLE active interval.
    deriv=da*(ra-rv*k['cd'])+dv*(rv-ra*k['cd'])
    assert_interval(deriv.lo>0,'Left active diameter is not increasing')
    assert_interval((ra-RU*k['cd']).hi<0,'Right active diameter is not decreasing')
    # The upper second ray hits the first outer tangent BEFORE its x=RU corner.
    sn,cs=sincos_range(a+1,b+1);v=[cs,sn]
    outside=cross(v,vecsub(tangent,k['s']))
    assert_interval(outside.lo>0,'Right clipping may pass beyond first tangent edge')
    # Every other possible vertex pair is uniformly separated from D*.
    other=[]
    for label,pts,skip in [('left',[A,B,C,V],(0,3)),('right',[A,B,C,p,tangent],(0,3))]:
        for i,j in combinations(range(len(pts)),2):
            if (i,j)==skip:continue
            d=norm2(vecsub(pts[i],pts[j]))
            assert_interval(d.hi<k['D2'].lo,f'Nonactive pair {label,i,j} is not excluded')
            other.append(dict(side=label,pair=[i,j],upper_m=d.sqrt().json()[1]))
    return dict(interval_degrees=[str(a),str(b)],left_squared_distance_derivative_half=deriv.json(),
                other_pairs=other,forward_ray_coefficients=forward,pass_=True)

# Exact rational convex polygon tools: the proof does not depend on scipy/NumPy.
def cross_q(a,b,c):return (b[0]-a[0])*(c[1]-a[1])-(b[1]-a[1])*(c[0]-a[0])
def hull(points):
    points=sorted(set(points))
    if len(points)<=2:return points
    l=[];u=[]
    for p in points:
        while len(l)>1 and cross_q(l[-2],l[-1],p)<=0:l.pop()
        l.append(p)
    for p in reversed(points):
        while len(u)>1 and cross_q(u[-2],u[-1],p)<=0:u.pop()
        u.append(p)
    return l[:-1]+u[:-1]

def rational_outer(k):
    sh,ch=sincos_deg(F('0.5'))
    vertices=[[IV.of(0),IV.of(0)],k['far'],[IV.of(RU),-RU*sh/ch],
              [IV.of(RU),RU*sh/ch],[RU*k['ce'],RU*k['se']]]
    pts=[]
    for v in vertices:
        for x in v[0].endpoints():
            for y in v[1].endpoints():pts.append((x,y))
    return hull(pts)

def clip(poly,nx,ny,b):
    if not poly:return []
    out=[]
    for i,cur in enumerate(poly):
        prev=poly[i-1];va=nx*prev[0]+ny*prev[1]-b;vb=nx*cur[0]+ny*cur[1]-b
        if (va<=0)!=(vb<=0):
            t=va/(va-vb);out.append((prev[0]+t*(cur[0]-prev[0]),prev[1]+t*(cur[1]-prev[1])))
        if vb<=0:out.append(cur)
    return out

def plane_relax(n,k):
    # Select rational coefficient midpoints and bound the residual on P,
    # where |x|<=1501, |y|<=27. This makes the rational plane an OUTER plane.
    mids=[q.midpoint() for q in n]
    dev=[max(abs(q.endpoints()[0]-m),abs(q.endpoints()[1]-m)) for q,m in zip(n,mids)]
    rhs=dot(n,k['s']).endpoints()[1]+1501*dev[0]+27*dev[1]
    return mids[0],mids[1],rhs

def angle_box_upper(k,poly,a,b,halfwidth=F(1)):
    # Union of all W(phi,eps), phi in [a,b], is contained in W([a-eps,b+eps]).
    if b-a+2*halfwidth>=180:return None
    sn,cs=sincos_deg(a-halfwidth);p=clip(poly,*plane_relax([sn,-cs],k))
    sn,cs=sincos_deg(b+halfwidth);p=clip(p,*plane_relax([-sn,cs],k))
    d=F(0)
    for x,y in combinations(p,2):d=max(d,(x[0]-y[0])**2+(x[1]-y[1])**2)
    return d

def heading_certificate(k,max_nodes=30000):
    poly=rational_outer(k);Dlo=k['D2'].endpoints()[0]
    stack=[(F(-180),F('-42.1')),(F('-42.0'),F(180))];leaves=[];n=0;max_other=F(0)
    while stack:
        a,b=stack.pop();n+=1
        if n>max_nodes:raise RuntimeError('Certificate node limit reached; no proof exported')
        ub=angle_box_upper(k,poly,a,b)
        if ub is not None and ub<Dlo:
            max_other=max(max_other,ub)
            leaves.append(dict(lo_deg=str(a),hi_deg=str(b),upper_squared_m=str(ub)))
        else:
            m=(a+b)/2
            if b-a<F(1,10**9):raise RuntimeError('Unresolved angle interval')
            stack.extend([(a,m),(m,b)])
    leaves.sort(key=lambda z:F(z['lo_deg']))
    # Independently check the partition, not just that the loop ended.
    intervals=[(F(x['lo_deg']),F(x['hi_deg'])) for x in leaves]+[(F('-42.1'),F('-42.0'))]
    intervals.sort();assert_interval(intervals[0][0]==-180 and intervals[-1][1]==180,'Missing ends')
    assert_interval(all(x[1]==y[0] for x,y in zip(intervals,intervals[1:])),'Coverage gap or overlap')
    return dict(nodes=n,leaf_count=len(leaves),outside_active_upper_m=IV.of(max_other).sqrt().json(),
                partition=leaves,pass_=True)


def candidate_certificate(k):
    """All reception-safe points within a +/-4 m square of either exact optimum.
    Position uncertainty increases the angular half-width by at most 0.625deg.
    """
    dmin=k['s'][1]-RU*k['se'];sb,cb=sincos_deg(F('0.625'))
    assert_interval((dmin*sb).sq().lo>32*SCALE,'Position-to-angle perturbation bound failed')
    poly=rational_outer(k);limit=F(36,25)*k['D2'].endpoints()[0]
    stack=[(F(-180),F(180))];parts=[];nodes=0;ubmax=F(0)
    while stack:
        a,b=stack.pop();nodes+=1
        if nodes>30000:raise RuntimeError('Candidate proof failed to terminate')
        u=angle_box_upper(k,poly,a,b,F('1.625'))
        if u is not None and u<limit:
            ubmax=max(ubmax,u);parts.append([str(a),str(b),str(u)])
        else:
            m=(a+b)/2;stack.extend([(a,m),(m,b)])
    parts.sort(key=lambda x:F(x[0]));assert F(parts[0][0])==-180 and F(parts[-1][1])==180
    assert all(F(a[1])==F(b[0]) for a,b in zip(parts,parts[1:]))
    return dict(half_side_m=4,extra_angle_degrees='0.625',position_to_angle_proved=True,
                bound_m=IV.of(ubmax).sqrt().json(),threshold_m=(k['D']*F(6,5)).json(),
                partition=parts,leaf_count=len(parts),pass_=True)

def operational_certificate(k):
    from .policy import choose_second_point
    proposed=choose_second_point((0.,0.),0.)['operational_local']
    # Each proposed finite decimal is treated exactly; the extra 1e-9m box
    # also covers ordinary coordinate conversion roundoff in this bounded task.
    point=[IV.of(str(x))+IV.of('0') for x in proposed]
    box=[IV(q.lo-10**31,q.hi+10**31) for q in point]
    se,ce=k['se'],k['ce'];lower=[box[0]*ce-box[1]*se,box[0]*se+box[1]*ce]
    for r in [5,1000]:
        for sign in [-1,1]:
            cc=[r*ce,sign*r*se]
            assert_interval(norm2(vecsub(box,cc)).hi<1000000*SCALE,'Operational guard is not strictly reception safe')
    A,B=lower;N=(RU*A-A.sq()-B.sq())*k['sd']+RU*B*k['cd']
    H=(RU-A)*k['s2']+B*k['c2'];T=N/H
    v=dict(k);v['s']=box;v['T']=T;v['D2']=RU**2+T.sq()-2*RU*T*k['cd'];v['D']=v['D2'].sqrt()
    active_checks(v);rest=heading_certificate(v)
    gap=v['D']-k['D']
    assert_interval(gap.hi<IV.of('0.0001').lo,'Operational deviation exceeds 0.1 mm objective budget')
    return dict(local_decimal_coordinates=[str(x) for x in proposed],
                coordinate_uncertainty_m='0.000000001',bound_m=v['D'].json(),
                optimality_loss_upper_m=gap.json()[1],nonactive_heading_leaves=rest['leaf_count'],pass_=True)

def verify_all(out:Path|None=None):
    t=time.perf_counter();k=constants()
    res=dict(arithmetic='integer fixed-point intervals, scale 10^40; exact rational clipping',
             pi=pi().json(),t_star_m=k['T'].json(),diameter_star_m=k['D'].json(),
             station_lower_coordinates_m=[k['sx'].json(),k['sy'].json()],
             station_original_coordinates_m=[x.json() for x in k['s']],
             lower_proof=lower_checks(k),active_upper_proof=active_checks(k))
    res['all_other_headings']=heading_certificate(k)
    res['candidate_region']=candidate_certificate(k)
    res['operational_guard']=operational_certificate(k)
    res['elapsed_s']=time.perf_counter()-t;res['pass']=True
    if out:
        out.parent.mkdir(parents=True,exist_ok=True);out.write_text(json.dumps(res,ensure_ascii=False,indent=2),encoding='utf-8')
    return res
\end{lstlisting}
%</paper-block>

%<paper-block id="appendix.h004" type="heading">
\subsection*{B.2 第三问最终策略程序}
%</paper-block>
%<paper-block id="appendix.p009" type="paragraph">
\noindent\textbf{文件：}\texttt{02\_\allowbreak{}最终代码/\allowbreak{}第三问/\allowbreak{}strategy.py}\\
\textbf{功能：}问题三动作、状态及基础定位接口
%</paper-block>
%<paper-block id="appendix.code007" type="code">
\begin{lstlisting}[style=appendixcode]
# -*- coding: utf-8 -*-
# strategy.py  (v2 重写版)
# 第三问：4 套定位清除策略（统一接口）
#
# 设计原则：
#   - 每个频道 ch 有状态：NEW / M1 / READY / CLEARED / SKIP
#   - 策略 step() 在外层 cursor 上推进，对当前 ch 调用 _next_action_for_ch(ch)
#   - 状态机不会回退：M1 → 一定要补第 2 个 bearing，然后切到 READY（再试定位）
#   - 失败的 clear 不会无限重试：超过 N 次失败 → SKIP
#
# 4 套策略的区别 = 第二个测向点的位置选取 + 清除前是否再移动补测：
#   1) StaticScan：固定在 (1000, 0)，第 2 点 = 沿 θ 走 800 m
#   2) Patrol：4 个站点轮换
#   3) BlinkHop：第 1 点随机在 (0,0) 起步，每测到 direction 沿 θ 跳 800 m
#   4) Hybrid：早期闪烁跳跃广撒网，后期 StaticScan 精定位
#
# 复用 Q1 v4 的 solve_problem_1 / localization_polygon。

import math
import os
import importlib.util as _ilu
from dataclasses import dataclass, field
from typing import List, Tuple, Optional, Dict, Set

# 复用 Q1 v4 算法 —— 优先用内联版本 problem1_v4_inline.py（自包含），
# 没有内联文件时再用外部多级 fallback 找 problem1_v4.py。
# 这样 v5 文件夹可以单独复制到任何机器，不需要 Q1 v4 源码。
def _locate_problem1_v4():
    """优先内联；找不到再 fallback 找外部 problem1_v4.py。
    返回 (module, source_label)。"""
    here = os.path.dirname(os.path.abspath(__file__))
    # 1) 内联版本（v5 自带）
    inline = os.path.join(here, 'problem1_v4_inline.py')
    if os.path.isfile(inline):
        return inline, 'inline'

    # 2) 外部多级 fallback
    fname = 'problem1_v4.py'
    candidates = []
    # 2a) 与本文件相对
    candidates.append(os.path.normpath(os.path.join(
        here, '..', '..', '第一问', 'v4', fname)))
    # 2b) 从 cwd 向上找
    cwd = os.getcwd()
    cur = cwd
    for _ in range(8):
        cand = os.path.join(cur, '第一问', 'v4', fname)
        if os.path.isfile(cand):
            candidates.append(os.path.abspath(cand))
            break
        parent = os.path.dirname(cur)
        if parent == cur:
            break
        cur = parent
    # 2c) 环境变量
    env = os.environ.get('CUMCM2026_ROOT')
    if env:
        candidates.append(os.path.normpath(os.path.join(env, '第一问', 'v4', fname)))
    # 2d) 常见路径
    common = [
        r'C:\Users\LENOVO\Desktop\CUMCM2026Problems',
        r'D:\CUMCM2026Problems',
        r'/mnt/c/Users/LENOVO/Desktop/CUMCM2026Problems',
        os.path.expanduser('~/Desktop/CUMCM2026Problems'),
        os.path.expanduser('~/CUMCM2026Problems'),
    ]
    for c in common:
        candidates.append(os.path.join(c, '第一问', 'v4', fname))

    for c in candidates:
        if os.path.isfile(c):
            return os.path.abspath(c), 'external'

    raise FileNotFoundError(
        '找不到 Q1 v4 算法源。请确认 v5 文件夹里有 problem1_v4_inline.py，'
        '或设置 CUMCM2026_ROOT 指向 CUMCM2026Problems 根目录。\n'
        '尝试过的路径：\n  ' + '\n  '.join(candidates))


_FIRST_FILE, _SOURCE = _locate_problem1_v4()
_spec = _ilu.spec_from_file_location('problem1_v4', _FIRST_FILE)
p1 = _ilu.module_from_spec(_spec)
_spec.loader.exec_module(p1)
print(f'[strategy] Q1 v4 算法源 = {_SOURCE}: {os.path.basename(_FIRST_FILE)}')

Vec = Tuple[float, float]
CHANNELS = list(range(1, 21))
ARENA_R = 1800.0
EPS_DEG = 1.0
CLEAR_R = 20.0
SPEED = 5.0
HOP_DIST = 800.0           # 闪烁跳跃：沿 direction 走多远
STATIC_POS: Vec = (1000.0, 0.0)
PATROL_POS: List[Vec] = [
    (1000.0, 0.0), (0.0, 1000.0), (-1000.0, 0.0), (0.0, -1000.0)
]
MAX_CLEAR_FAILS = 3        # 同一频道清除失败 ≥ N 次 → SKIP

# 每个频道的内部状态
CH_NEW = 'new'
CH_M1 = 'm1'
CH_READY = 'ready'
CH_CLEARED = 'cleared'
CH_SKIP = 'skip'


# ============================================================
# Action
# ============================================================
@dataclass
class Action:
    kind: str            # 'measure' | 'clear' | 'done'
    pos: Vec = (0.0, 0.0)
    ch: int = 1

    def __repr__(self):
        if self.kind == 'done':
            return "Action(done)"
        return f"Action({self.kind}, pos={self.pos}, ch={self.ch})"


@dataclass
class State:
    pos: Vec
    ch: int
    cleared: Set[int] = field(default_factory=set)
    skipped: Set[int] = field(default_factory=set)
    bearing_records: Dict[int, List[Tuple[Vec, float]]] = field(default_factory=dict)
    virtual_time: float = 0.0


# ============================================================
# 工具：定位 + 清除
# ============================================================
def localize(records: List[Tuple[Vec, float]],
             eps: float = EPS_DEG) -> Optional[Tuple[Vec, float]]:
    """对一组 (S, θ) 用 Q1 算法定位，返回 MEC 圆心和半径。"""
    if len(records) < 2:
        return None
    dets = [r[0] for r in records]
    thetas = [r[1] for r in records]
    res = p1.solve_problem_1(dets, thetas, eps_deg=eps,
                             R_eff=1500.0, R_target=ARENA_R, N_disk=32)
    if res['n_vertices'] < 3 or res['D'] < 1e-3:
        return None
    return res['mec_center'], res['mec_radius']


def second_pos_from_bearing(S: Vec, theta: float, dist: float = HOP_DIST) -> Vec:
    """从 S 沿 θ 方向走 dist 米。"""
    rad = math.radians(theta)
    return (S[0] + dist * math.cos(rad),
            S[1] + dist * math.sin(rad))


# ============================================================
# 基类
# ============================================================
class BaseStrategy:
    name: str = "base"

    def __init__(self):
        self.ch_state: Dict[int, str] = {ch: CH_NEW for ch in CHANNELS}
        self.bearings: Dict[int, List[Tuple[Vec, float]]] = {ch: [] for ch in CHANNELS}
        self.clear_fails: Dict[int, int] = {ch: 0 for ch in CHANNELS}
        self.cursor = 0
        self.done = False
        self._attempts = 0  # 总步数上限

    # 子类重写：返回某个频道的第 1 / 第 2 个 measure 位置
    def first_pos(self, ch: int, state: State) -> Vec:
        return STATIC_POS

    def second_pos(self, ch: int, state: State,
                   first_record: Tuple[Vec, float]) -> Vec:
        S, theta = first_record
        return second_pos_from_bearing(S, theta, HOP_DIST)

    def next_action_for_ch(self, ch: int, state: State) -> Optional[Action]:
        """对某个频道返回下一个动作；None 表示切到下个频道。"""
        st = self.ch_state[ch]
        if st == CH_CLEARED or st == CH_SKIP:
            return None
        recs = self.bearings[ch]
        if st == CH_NEW:
            return Action('measure', self.first_pos(ch, state), ch)
        if st == CH_M1:
            assert len(recs) == 1
            return Action('measure',
                          self.second_pos(ch, state, recs[0]), ch)
        if st == CH_READY:
            # 已经有 ≥ 2 个 bearing；尝试定位清除
            loc = localize(recs)
            if loc is None:
                # 几何退化（两扇形反向等），跳过
                self.ch_state[ch] = CH_SKIP
                state.skipped.add(ch)
                return None
            c, r = loc
            if r <= CLEAR_R:
                return Action('clear', c, ch)
            # 定位还不够准，再补一个 bearing
            # 简单策略：在上次某个 bearing 附近 + 沿另一方向 600m
            extra = self._extra_pos(ch, recs)
            return Action('measure', extra, ch)
        return None

    def _extra_pos(self, ch: int, recs: List[Tuple[Vec, float]]) -> Vec:
        """第 3+ 个 bearing 的位置：取已有 bearing 中心 + 沿某方向 600m。"""
        cx = sum(r[0][0] for r in recs) / len(recs)
        cy = sum(r[0][1] for r in recs) / len(recs)
        # 沿"已收 bearing 角度差最大"的方向
        thetas = [r[1] for r in recs]
        a = math.radians(sum(thetas) / len(thetas) + 90)  # 垂直方向
        return (cx + 600 * math.cos(a), cy + 600 * math.sin(a))

    def step(self, state: State) -> Action:
        if self.done:
            return Action('done')
        self._attempts += 1
        if self._attempts > 8000:
            self.done = True
            return Action('done')

        # 扫一圈找下一个需要处理的频道
        for _ in range(len(CHANNELS)):
            ch = CHANNELS[self.cursor % len(CHANNELS)]
            self.cursor += 1
            if ch in state.cleared or ch in state.skipped:
                continue
            a = self.next_action_for_ch(ch, state)
            if a is not None:
                return a
        # 全部清完或无新动作
        self.done = True
        return Action('done')

    # --------- 事件回调 ---------
    def on_measure(self, state: State, ch: int, result: str,
                   svd: Optional[float]):
        if result == 'direction' and svd is not None:
            recs = self.bearings[ch]
            # 避免重复（同一 pos 的 bearing 只记一次）
            pos = state.pos
            if not any((abs(S[0] - pos[0]) < 1e-3 and
                        abs(S[1] - pos[1]) < 1e-3) for S, _ in recs):
                recs.append((pos, svd))
            # 更新状态
            if self.ch_state[ch] == CH_NEW:
                self.ch_state[ch] = CH_M1
            elif self.ch_state[ch] in (CH_M1, CH_READY):
                # M1 → READY（如果 ≥ 2 个）；READY 再多收也无所谓
                if len(recs) >= 2:
                    self.ch_state[ch] = CH_READY
        elif result == 'no_signal':
            self.ch_state[ch] = CH_SKIP
            state.skipped.add(ch)
        elif result == 'near':
            # 已经在源附近 5m：直接进入"尝试清除"状态（用当前位置）
            self.ch_state[ch] = CH_READY
            # 不需要再测，把 state.pos 作为一个"假 bearing"也加入
            self.bearings[ch].append((state.pos, 0.0))
            self.bearings[ch].append((state.pos, 90.0))

    def on_clear(self, state: State, ch: int, success: bool):
        if success:
            self.ch_state[ch] = CH_CLEARED
            state.cleared.add(ch)
        else:
            self.clear_fails[ch] += 1
            if self.clear_fails[ch] >= MAX_CLEAR_FAILS:
                self.ch_state[ch] = CH_SKIP
                state.skipped.add(ch)


# ============================================================
# 具体策略
# ============================================================
class StaticScan(BaseStrategy):
    """固定在 (1000, 0)，对每个频道：
       1) measure(ch) @ (1000, 0)
       2) 若 direction：沿 θ 走 800m，再 measure(ch)
       3) 定位 → clear
    """
    name = "StaticScan"

    def first_pos(self, ch: int, state: State) -> Vec:
        return STATIC_POS


class Patrol(BaseStrategy):
    """在 4 个站位之间轮换：每站位扫所有频道。"""
    name = "Patrol"

    def __init__(self):
        super().__init__()
        self.site_idx = 0

    def first_pos(self, ch: int, state: State) -> Vec:
        return PATROL_POS[self.site_idx]

    def step(self, state: State) -> Action:
        # 在每站位扫一轮（len(CHANNELS) 个动作）
        a = super().step(state)
        # 当 cursor 回到 0，说明本轮结束 → 换站位
        if self.cursor % len(CHANNELS) == 0:
            self.site_idx = (self.site_idx + 1) % len(PATROL_POS)
        return a


class BlinkHop(BaseStrategy):
    """在初始 (0,0) 起步，每收 1 个 direction 就沿 θ 跳 800m 当作下个测点。"""
    name = "BlinkHop"

    def first_pos(self, ch: int, state: State) -> Vec:
        return state.pos  # 沿用当前位置

    def second_pos(self, ch: int, state: State,
                   first_record: Tuple[Vec, float]) -> Vec:
        S, theta = first_record
        return second_pos_from_bearing(S, theta, HOP_DIST)


class Hybrid(BaseStrategy):
    """早期：闪烁跳跃（pos 跟着 direction 走）；后期：固定到 STATIC_POS 精定位。"""
    name = "Hybrid"

    def __init__(self):
        super().__init__()
        self._phase = 'hop'   # 'hop' | 'static'
        self._switch_at = 5   # 已清除 ≥ 5 切换到 static

    def first_pos(self, ch: int, state: State) -> Vec:
        if self._phase == 'hop':
            return state.pos
        return STATIC_POS

    def second_pos(self, ch: int, state: State,
                   first_record: Tuple[Vec, float]) -> Vec:
        if self._phase == 'hop':
            S, theta = first_record
            return second_pos_from_bearing(S, theta, HOP_DIST)
        # static：固定基线，沿 θ 反方向 800m 作第 2 点
        S, theta = first_record
        rad = math.radians(theta)
        # 取与 first_pos 反向 + 800m
        return (STATIC_POS[0] - 800 * math.cos(rad),
                STATIC_POS[1] - 800 * math.sin(rad))

    def step(self, state: State) -> Action:
        # 阶段切换
        if self._phase == 'hop' and len(state.cleared) >= self._switch_at:
            self._phase = 'static'
        return super().step(state)


# ============================================================
# 工厂
# ============================================================
def make_strategy(name: str) -> BaseStrategy:
    return {
        'BlinkHop': BlinkHop,
        'StaticScan': StaticScan,
        'Patrol': Patrol,
        'Hybrid': Hybrid,
    }[name]()


if __name__ == "__main__":
    s = make_strategy('Hybrid')
    st = State(pos=(0.0, 0.0), ch=1)
    for i in range(8):
        a = s.step(st)
        print(i, a)
        if a.kind == 'done':
            break
\end{lstlisting}
%</paper-block>

%<paper-block id="appendix.p010" type="paragraph">
\noindent\textbf{文件：}\texttt{02\_\allowbreak{}最终代码/\allowbreak{}第三问/\allowbreak{}problem1\_\allowbreak{}v4\_\allowbreak{}inline.py}\\
\textbf{功能：}问题三内联交会定位几何模块
%</paper-block>
%<paper-block id="appendix.code008" type="code">
\begin{lstlisting}[style=appendixcode]
# -*- coding: utf-8 -*-
# problem1_v4_inline.py  （v5 自包含的 Q1 v4 算法副本）
#
# 这是 Q1 v4（problem1_v4.py）的完整内联版，src 在 解题库/第一问/v4/。
# v5 不依赖外部 Q1 v4 源码；本文件已经包含：
#   - HalfPlane、make_sector_halfplanes（楔形张角）
#   - hpi_intersect、halfplane_intersect、clip_polygon_by_disk（Sutherland-Hodgman）
#   - polygon_diameter、mec_welzl、diameter_circle_covers（直径 + Welzl MEC）
#   - localization_polygon、solve_problem_1（定位 + 直径圆覆盖判据）

import math
import random
from typing import List, Tuple, Optional, Sequence

Vec = Tuple[float, float]
Poly = List[Vec]


class HalfPlane:
    """半平面 n·P <= d（外向法向 n，闭合侧为内侧）。"""
    __slots__ = ('n', 'd')

    def __init__(self, n: Vec, d: float):
        self.n = n
        self.d = d

    def contains(self, P: Vec) -> bool:
        return self.n[0] * P[0] + self.n[1] * P[1] <= self.d + 1e-9


def make_sector_halfplanes(S: Vec, theta_deg: float, eps_deg: float
                           ) -> Tuple[HalfPlane, HalfPlane]:
    """扇形构造：点 P 与 S 的连线方向 ∈ [θ−ε, θ+ε]。"""
    lo = math.radians(theta_deg - eps_deg)
    hi = math.radians(theta_deg + eps_deg)
    u_lo = (math.cos(lo), math.sin(lo))
    u_hi = (math.cos(hi), math.sin(hi))
    n_lo = (u_lo[1], -u_lo[0])
    n_hi = (-u_hi[1], u_hi[0])
    d_lo = n_lo[0] * S[0] + n_lo[1] * S[1]
    d_hi = n_hi[0] * S[0] + n_hi[1] * S[1]
    return HalfPlane(n_lo, d_lo), HalfPlane(n_hi, d_hi)


def _seg_line_intersect(p1: Vec, p2: Vec, n: Vec, d: float) -> Vec:
    dx = p2[0] - p1[0]
    dy = p2[1] - p1[1]
    denom = n[0] * dx + n[1] * dy
    if abs(denom) < 1e-12:
        return p1
    t = (d - n[0] * p1[0] - n[1] * p1[1]) / denom
    return (p1[0] + t * dx, p1[1] + t * dy)


def hpi_intersect(poly: Poly, halfplanes: Sequence[HalfPlane]
                  ) -> Optional[Poly]:
    """Sutherland-Hodgman：给定初始凸多边形 poly，依次用各半平面裁剪。"""
    for hp in halfplanes:
        if poly is None or len(poly) == 0:
            return None
        new_poly: List[Vec] = []
        n, d = hp.n, hp.d
        m = len(poly)
        for i in range(m):
            cur = poly[i]
            prv = poly[i - 1]
            cur_in = (n[0] * cur[0] + n[1] * cur[1]) <= d + 1e-9
            prv_in = (n[0] * prv[0] + n[1] * prv[1]) <= d + 1e-9
            if cur_in:
                if not prv_in:
                    new_poly.append(_seg_line_intersect(prv, cur, n, d))
                new_poly.append(cur)
            elif prv_in:
                new_poly.append(_seg_line_intersect(prv, cur, n, d))
        poly = new_poly
    if poly is None or len(poly) < 3:
        return None
    return poly


def _bbox_polygon(bbox: float) -> Poly:
    b = float(bbox)
    return [(-b, -b), (b, -b), (b, b), (-b, b)]


def halfplane_intersect(halfplanes: Sequence[HalfPlane], bbox: float = 5000.0
                        ) -> Optional[Poly]:
    return hpi_intersect(_bbox_polygon(bbox), halfplanes)


def clip_polygon_by_disk(poly: Poly, c: Vec, r: float, N: int = 64
                         ) -> Optional[Poly]:
    """用 N 条线段近似圆盘 |P−c| ≤ r，截多边形 poly。"""
    if poly is None or len(poly) < 3:
        return None
    hps = []
    for k in range(N):
        a = 2 * math.pi * k / N
        n = (math.cos(a), math.sin(a))
        d = r + n[0] * c[0] + n[1] * c[1]
        hps.append(HalfPlane(n, d))
    return hpi_intersect(poly, hps)


def polygon_diameter(poly: Poly) -> Tuple[float, Vec, Vec]:
    """返回 (D, A, B)。O(n²) 暴力。"""
    n = len(poly)
    if n < 2:
        p = poly[0] if poly else (0.0, 0.0)
        return 0.0, p, p
    best_d2 = -1.0
    best_i, best_j = 0, 1
    for i in range(n):
        xi, yi = poly[i]
        for j in range(i + 1, n):
            dx = xi - poly[j][0]
            dy = yi - poly[j][1]
            d2 = dx * dx + dy * dy
            if d2 > best_d2:
                best_d2, best_i, best_j = d2, i, j
    return math.sqrt(best_d2), poly[best_i], poly[best_j]


def _circle_from_2(a: Vec, b: Vec) -> Tuple[Vec, float]:
    c = ((a[0] + b[0]) * 0.5, (a[1] + b[1]) * 0.5)
    r2 = (a[0] - c[0]) ** 2 + (a[1] - c[1]) ** 2
    return c, r2


def _circle_from_3(a: Vec, b: Vec, c: Vec) -> Optional[Tuple[Vec, float]]:
    ax, ay = a
    bx, by = b
    cx, cy = c
    d = 2 * (ax * (by - cy) + bx * (cy - ay) + cx * (ay - by))
    if abs(d) < 1e-12:
        return None
    ux = ((ax * ax + ay * ay) * (by - cy) +
          (bx * bx + by * by) * (cy - ay) +
          (cx * cx + cy * cy) * (ay - by)) / d
    uy = ((ax * ax + ay * ay) * (cx - bx) +
          (bx * bx + by * by) * (ax - cx) +
          (cx * cx + cy * cy) * (bx - ax)) / d
    ctr = (ux, uy)
    r2 = (ax - ux) ** 2 + (ay - uy) ** 2
    return ctr, r2


def mec_welzl(points: Sequence[Vec]) -> Tuple[Vec, float]:
    """Welzl 最小包围圆（迭代版，O(n) 期望时间）。"""
    import sys
    sys.setrecursionlimit(10000)
    P = list(points)
    random.shuffle(P)

    def circle_trivial(R: Sequence[Vec]) -> Tuple[Vec, float]:
        if len(R) == 0:
            return (0.0, 0.0), 0.0
        if len(R) == 1:
            return R[0], 0.0
        if len(R) == 2:
            return _circle_from_2(R[0], R[1])
        c, r2 = _circle_from_3(R[0], R[1], R[2])
        if c is None:
            return _circle_from_2(R[0], R[1])[0], max(
                (R[0][0] - R[1][0]) ** 2 + (R[0][1] - R[1][1]) ** 2,
                (R[0][0] - R[2][0]) ** 2 + (R[0][1] - R[2][1]) ** 2,
                (R[1][0] - R[2][0]) ** 2 + (R[1][1] - R[2][1]) ** 2,
            ) * 0.25
        return c, r2

    def mec(P: List[Vec], R: List[Vec], n: int) -> Tuple[Vec, float]:
        if n == 0 or len(R) == 3:
            return circle_trivial(R)
        p = P[n - 1]
        c, r2 = mec(P, R, n - 1)
        if (p[0] - c[0]) ** 2 + (p[1] - c[1]) ** 2 <= r2 + 1e-9:
            return c, r2
        return mec(P, R + [p], n - 1)

    c, r2 = mec(P, [], len(P))
    return c, math.sqrt(r2)


def diameter_circle_covers(poly: Poly, A: Vec, B: Vec
                           ) -> Tuple[bool, float, Tuple[Vec, float]]:
    """判断以 AB 为直径的圆是否能覆盖多边形。"""
    mid = ((A[0] + B[0]) * 0.5, (A[1] + B[1]) * 0.5)
    D = math.hypot(A[0] - B[0], A[1] - B[1])
    R = D / 2
    off = -1e9
    for p in poly:
        d = math.hypot(p[0] - mid[0], p[1] - mid[1])
        off = max(off, d - R)
    covered = off <= 1e-9
    c_mec, r_mec = mec_welzl(poly)
    return covered, off, (c_mec, r_mec)


def localization_polygon(dets: Sequence[Vec], thetas: Sequence[float],
                         eps_deg: float = 1.0,
                         R_eff: float = 1500.0,
                         R_target: float = 1800.0,
                         N_disk: int = 64) -> Optional[Poly]:
    """构造定位多边形 P = ∩扇形_i ∩ D(S_i, R_eff) ∩ D(O, R_target)。"""
    hps = []
    for S, th in zip(dets, thetas):
        hp_lo, hp_hi = make_sector_halfplanes(S, th, eps_deg)
        hps.extend([hp_lo, hp_hi])

    poly = halfplane_intersect(hps, bbox=5000.0)
    if poly is None:
        return None

    poly = clip_polygon_by_disk(poly, (0.0, 0.0), R_target, N_disk)
    if poly is None:
        return None

    for S in dets:
        poly = clip_polygon_by_disk(poly, S, R_eff, N_disk)
        if poly is None:
            return None

    return poly


def solve_problem_1(dets: Sequence[Vec], thetas: Sequence[float],
                    eps_deg: float = 1.0,
                    R_eff: float = 1500.0,
                    R_target: float = 1800.0,
                    N_disk: int = 64) -> dict:
    """完整求解第一问。返回含 mec_center, mec_radius, D 等的字典。"""
    poly = localization_polygon(dets, thetas, eps_deg, R_eff, R_target, N_disk)
    if poly is None:
        return {
            'D': 0.0, 'R_star': 0.0, 'ratio_2R_D': 0.0,
            'covered': False, 'max_offset': 0.0,
            'A': (0, 0), 'B': (0, 0), 'mec_center': (0, 0), 'mec_radius': 0.0,
            'n_vertices': 0, 'poly': [],
        }
    D, A, B = polygon_diameter(poly)
    covered, off, (c_mec, r_mec) = diameter_circle_covers(poly, A, B)
    return {
        'D': D,
        'R_star': r_mec,
        'ratio_2R_D': 2 * r_mec / D if D > 0 else 0.0,
        'covered': covered,
        'max_offset': off,
        'A': A,
        'B': B,
        'mec_center': c_mec,
        'mec_radius': r_mec,
        'n_vertices': len(poly),
        'poly': poly,
    }
\end{lstlisting}
%</paper-block>

%<paper-block id="appendix.p011" type="paragraph">
\noindent\textbf{文件：}\texttt{02\_\allowbreak{}最终代码/\allowbreak{}第三问/\allowbreak{}strategy\_\allowbreak{}v8.py}\\
\textbf{功能：}问题三逐频道覆盖、定位、清除与失败恢复
%</paper-block>
%<paper-block id="appendix.code009" type="code">
\begin{lstlisting}[style=appendixcode]
# -*- coding: utf-8 -*-
"""问题三 v8：就近定位、沿途测向、按频道证明搜索覆盖。

策略仅接收 State 和动作响应，不读取模拟器的源位置、数量或随机种子。
全向信号的 no_signal 排除半径 1000 m 的圆盘；网格按整格覆盖判断，
不是将未覆盖格点简单当成没有目标。虚拟时间限制由 /enter 的驱动器处理。
"""

import math
from dataclasses import dataclass, field

import numpy as np

from strategy import Action, State, p1

Vec = tuple[float, float]
CHANNELS = tuple(range(1, 21))
ARENA_R = 1800.0
RECEIVE_MIN = 1000.0
RECEIVE_MAX = 1500.0
# 官方 svd_deg 保留两位小数；量化最多再引入 0.005 度。
EPS_DEG = 1.0050001
CLEAR_R = 20.0


def distance(a, b):
    return math.hypot(a[0] - b[0], a[1] - b[1])


def centroid(poly):
    area = cx = cy = 0.0
    for a, b in zip(poly, poly[1:] + poly[:1]):
        cross = a[0] * b[1] - b[0] * a[1]
        area += cross
        cx += (a[0] + b[0]) * cross
        cy += (a[1] + b[1]) * cross
    if abs(area) < 1e-9:
        return tuple(sum(p[i] for p in poly) / len(poly) for i in (0, 1))
    return cx / (3 * area), cy / (3 * area)


def convex_hull(points):
    points = sorted(set(points))
    if len(points) < 3:
        return points
    def cross(a, b, c):
        return (b[0] - a[0]) * (c[1] - a[1]) - (b[1] - a[1]) * (c[0] - a[0])
    lower, upper = [], []
    for p in points:
        while len(lower) >= 2 and cross(lower[-2], lower[-1], p) <= 1e-9:
            lower.pop()
        lower.append(p)
    for p in reversed(points):
        while len(upper) >= 2 and cross(upper[-2], upper[-1], p) <= 1e-9:
            upper.pop()
        upper.append(p)
    return lower[:-1] + upper[:-1]


def exclude_disk(poly, pos, radius=RECEIVE_MIN):
    """从凸定位区减去无信号圆盘的内接多边形，再取保守凸包。

    只删除能确认不可能的位置；凸包可能保留额外区域，不会删除真目标。
    """
    center, bound = enclosing_circle(poly)
    if distance(center, pos) - bound >= radius:
        return poly
    pieces = []
    rest = poly
    n = 64
    support = radius * math.cos(math.pi / n) - 1e-6
    for k in range(n):
        angle = 2 * math.pi * k / n
        normal = (math.cos(angle), math.sin(angle))
        d = support + normal[0] * pos[0] + normal[1] * pos[1]
        outer = p1.hpi_intersect(rest, [p1.HalfPlane((-normal[0], -normal[1]), -d)])
        if outer:
            pieces.extend(outer)
        rest = p1.hpi_intersect(rest, [p1.HalfPlane(normal, d)])
        if not rest:
            break
    return convex_hull(pieces)


def reception_halfplanes(positive, negative):
    """同一源的接收半径不变：d(X,S) <= R < d(X,N) 给出线性约束。"""
    return [p1.HalfPlane((2 * (n[0] - s[0]), 2 * (n[1] - s[1])),
                        n[0]**2 + n[1]**2 - s[0]**2 - s[1]**2)
            for s in positive for n in negative]


def enclosing_circle(poly):
    """小凸多边形的确定性包围圆；优先检验直径圆，余下枚举三点。

    最后重新取最大顶点距离，避免浮点误差低估清除半径。
    """
    if len(poly) == 1:
        return poly[0], 0.0
    diameter, a, b = p1.polygon_diameter(poly)
    center = ((a[0] + b[0]) / 2, (a[1] + b[1]) / 2)
    radius = max(distance(center, v) for v in poly)
    if radius <= diameter / 2 + 1e-7:
        return center, radius + 1e-7
    best = (center, radius)
    for i, a in enumerate(poly):
        for j in range(i):
            b = poly[j]
            for k in range(j):
                candidate = p1._circle_from_3(a, b, poly[k])
                if candidate is None:
                    continue
                c, r2 = candidate
                if r2 >= best[1] ** 2:
                    continue
                r = max(distance(c, v) for v in poly)
                if r * r <= r2 + 1e-5:
                    best = c, r + 1e-7
    return best


@dataclass
class Track:
    records: list = field(default_factory=list)
    poly: list = field(default_factory=list)
    center: Vec = (0.0, 0.0)
    radius: float = math.inf
    near: Vec | None = None
    failures: int = 0
    estimate: Vec = (0.0, 0.0)
    probe_after_fail: bool = False


class Coverage:
    """覆盖证明：每格外接圆全部落入一次无信号检测的 1000 m 圆内。

    保留所有与目标圆相交的方格，包含边缘格；已排除的格子不会遗漏
    目标圆边缘。此条件比圆盘并集覆盖更保守，但不会虚报完整覆盖。
    """

    def __init__(self, cell_size=60.0):
        n = math.ceil(2 * ARENA_R / cell_size)
        step = 2 * ARENA_R / n
        a = -ARENA_R + step * (np.arange(n) + 0.5)
        x, y = np.meshgrid(a, a)
        nearest_x = np.maximum(np.abs(x) - step / 2, 0)
        nearest_y = np.maximum(np.abs(y) - step / 2, 0)
        inside = nearest_x**2 + nearest_y**2 <= ARENA_R**2
        self.points = np.column_stack((x[inside], y[inside]))
        self.cover_radius = RECEIVE_MIN - step / math.sqrt(2) - 1e-6
        self.remaining = {ch: np.ones(len(self.points), dtype=bool)
                          for ch in CHANNELS}

    def covered_at(self, pos):
        delta = self.points - pos
        return np.sum(delta * delta, axis=1) <= self.cover_radius**2

    def gain(self, ch, mask):
        return int(np.count_nonzero(self.remaining[ch] & mask))

    def exclude(self, ch, pos):
        self.remaining[ch] &= ~self.covered_at(pos)

    def absent(self, ch):
        return not np.any(self.remaining[ch])


class AdaptiveV8:
    name = 'AdaptiveV8'

    def __init__(self, lateral=150.0, scan_fraction=0.10,
                 cell_size=30.0, opportunistic=True, first_hop=300.0,
                 probe_radius=80.0):
        self.lateral = lateral
        self.scan_fraction = scan_fraction
        self.opportunistic = opportunistic
        self.coverage = Coverage(cell_size)
        self.tracks = {}
        self.negative_positions = {ch: [] for ch in CHANNELS}
        self.cleared = set()
        self.absent = set()
        self.done = False
        self.completion_reason = None
        self.scan_queue = []
        self.active_ch = None
        self.scan_pending = True
        self.exploring = False
        self.geometry_conflicts = 0
        self.first_hop = first_hop
        self.probe_radius = probe_radius
        self.station_positions = []
        for radius in (800., 1150., 1350., 1550.):
            for k in range(24):
                a = k * math.pi / 12
                self.station_positions.append((radius * math.cos(a), radius * math.sin(a)))
        self.station_positions.extend((float(x), float(y))
                                      for x in range(-1600, 1601, 200)
                                      for y in range(-1600, 1601, 200)
                                      if x * x + y * y <= 1800**2)
        self.station_masks = np.asarray([self.coverage.covered_at(p)
                                        for p in self.station_positions])

    def _unknown(self):
        return [ch for ch in CHANNELS if ch not in self.tracks
                and ch not in self.cleared and ch not in self.absent]

    def _start_scan(self, state, force=False):
        mask = self.coverage.covered_at(state.pos)
        unknown = [ch for ch in self._unknown()
                   if self.coverage.gain(ch, mask) >= (1 if force else max(1, int(
                       min(len(self.coverage.points) * self.scan_fraction,
                           np.count_nonzero(self.coverage.remaining[ch]) * 0.8))))]
        useful = []
        if self.opportunistic:
            for ch, track in self.tracks.items():
                if ch in self.cleared or track.radius <= max(CLEAR_R, self.probe_radius):
                    continue
                if not track.records:
                    continue
                nearest = min(distance(state.pos, p) for p, _ in track.records)
                if nearest < 100 or distance(state.pos, track.center) > 1250:
                    continue
                # 只有新视点能提供足够夹角时才补测，避免原地重复检测。
                old = min(track.records, key=lambda r: distance(state.pos, r[0]))[0]
                a = math.atan2(track.center[1] - old[1], track.center[0] - old[0])
                b = math.atan2(track.center[1] - state.pos[1],
                               track.center[0] - state.pos[0])
                if abs(math.sin(a - b)) > 0.15:
                    useful.append(ch)
        self.scan_queue = sorted(set(unknown + useful),
                                 key=lambda ch: (ch != state.ch, ch))
        self.scan_pending = False

    def _local_action(self, ch, state):
        track = self.tracks[ch]
        if track.probe_after_fail:
            if any(distance(state.pos, p) < 1 for p, _ in track.records):
                return Action('measure', (state.pos[0] + 25, state.pos[1] + 25), ch)
            return Action('measure', state.pos, ch)
        if track.near is not None:
            return Action('clear', track.near, ch)
        if track.radius <= CLEAR_R - 1e-5:
            # 利用 20 m 清除半径，向当前位置平移仍能覆盖整个定位区域。
            d = distance(state.pos, track.center)
            shift = min(d, CLEAR_R - track.radius - 1e-5)
            if d > 1e-9:
                pos = (track.center[0] + (state.pos[0] - track.center[0]) * shift / d,
                       track.center[1] + (state.pos[1] - track.center[1]) * shift / d)
            else:
                pos = track.center
            return Action('clear', pos, ch)

        if (track.radius <= self.probe_radius and len(track.records) >= 2
                and track.failures == 0):
            return Action('clear', track.estimate, ch)

        pos = track.estimate
        # 单条方向只约束一条长带：增加横向基线，避免沿射线走造成退化。
        if len(track.records) == 1:
            old, angle = track.records[-1]
            theta = math.radians(angle)
            if self.first_hop is not None:
                if (distance(state.pos, old) > 150 and distance(state.pos, pos) < 1250
                        and not any(distance(state.pos, p) < 1 for p in self.negative_positions[ch])):
                    return Action('measure', state.pos, ch)
                length = min(distance(old, pos), self.first_hop)
                pos = (old[0] + length * math.cos(theta),
                       old[1] + length * math.sin(theta))
            options = [(pos[0] + sign * self.lateral * -math.sin(theta),
                        pos[1] + sign * self.lateral * math.cos(theta))
                       for sign in (-1, 1)]
            pos = min(options, key=lambda p: distance(state.pos, p))
        elif any(distance(pos, p) < 10 for p, _ in track.records):
            _, a, b = p1.polygon_diameter(track.poly)
            theta = math.atan2(b[1] - a[1], b[0] - a[0]) + math.pi / 2
            length = max(25.0, min(self.lateral, track.radius * 0.5))
            options = [(pos[0] + sign * length * math.cos(theta),
                        pos[1] + sign * length * math.sin(theta))
                       for sign in (-1, 1)]
            pos = min(options, key=lambda p: distance(state.pos, p))
        return Action('measure', pos, ch)

    def _explore_action(self, state):
        unknown = self._unknown()
        if not unknown:
            return None
        weights = np.sum([self.coverage.remaining[ch] for ch in unknown], axis=0)
        indices = np.flatnonzero(weights)
        if len(indices) == 0:
            return None
        points = self.coverage.points
        # 候选仅由未覆盖区域生成，不读取源位置。稀疏取点并加入最近未覆盖点。
        stride = max(1, len(indices) // 100)
        choices = points[indices[::stride]]
        closest = points[indices[np.argmin(np.sum((points[indices] - state.pos)**2,
                                                axis=1))]]
        choices = np.vstack((choices, closest))
        best = None
        for candidate in choices:
            # 朝未覆盖区域走到足以覆盖其附近的站位，可比走到格心少走一段。
            d = distance(state.pos, candidate)
            if d > 600:
                candidate = np.asarray(state.pos) + (
                    candidate - state.pos) * max(0.4, (d - 500) / d)
            mask = self.coverage.covered_at(candidate)
            gain = int(np.sum(weights[mask]))
            if not gain:
                continue
            nscan = sum(self.coverage.gain(ch, mask) > 0 for ch in unknown)
            cost = distance(state.pos, candidate) / 5 + 6 * nscan + 25
            score = gain / cost
            if best is None or score > best[0]:
                best = (score, tuple(candidate), mask)
        if best is None:
            raise RuntimeError('搜索区域非空但无法生成覆盖站位')
        _, pos, mask = best
        channels = [ch for ch in unknown if self.coverage.gain(ch, mask)]
        self.exploring = True
        self.scan_pending = True
        return Action('measure', pos, channels[0])

    @staticmethod
    def _route(start, nodes):
        """最近邻起点 + 开放路径 2-opt，只规划剩余目标，不要求返回原点。"""
        remaining = list(nodes)
        route = []
        pos = start
        while remaining:
            node = min(remaining, key=lambda node: distance(pos, node[1]))
            remaining.remove(node)
            route.append(node)
            pos = node[1]
        for _ in range(12):
            changed = False
            for i in range(len(route) - 1):
                prev = start if i == 0 else route[i - 1][1]
                for j in range(i + 1, len(route)):
                    old = distance(prev, route[i][1])
                    new = distance(prev, route[j][1])
                    if j + 1 < len(route):
                        old += distance(route[j][1], route[j + 1][1])
                        new += distance(route[i][1], route[j + 1][1])
                    if new + 1e-6 < old:
                        route[i:j + 1] = reversed(route[i:j + 1])
                        changed = True
            if not changed:
                break
        return route

    def _next_destination(self, state):
        known = [(ch, track.estimate) for ch, track in self.tracks.items()
                 if ch not in self.cleared]
        unknown = self._unknown()
        if not unknown:
            route = self._route(state.pos, known)
            return route[0] if route else None
        remaining = np.any([self.coverage.remaining[ch] for ch in unknown], axis=0)
        return self._cover_route(state, known, remaining, len(unknown))

    def _cover_route(self, state, known, remaining, n_unknown):
        routes = [self._build_cover_route(state, known, remaining, n_unknown, power)
                  for power in (0.6, 1.0, 1.6)]
        def cost(route):
            previous = state.pos
            value = 0.0
            for ch, pos in route:
                value += distance(previous, pos) / 5
                if ch < 0:
                    value += 6 * n_unknown
                previous = pos
            return value
        route = min(routes, key=cost)
        return route[0] if route else None

    def _build_cover_route(self, state, known, remaining, n_unknown, power):
        """在已知目标路径中插入补充扫描点，以新增覆盖量/额外耗时选点。"""
        route = self._route(state.pos, known)
        uncovered = remaining.copy()
        for _, pos in known:
            uncovered &= ~self.coverage.covered_at(pos)
        candidates = self.station_positions
        masks = self.station_masks
        for _ in range(12):
            if not np.any(uncovered):
                break
            best = None
            for i, (pos, mask) in enumerate(zip(candidates, masks)):
                gain = int(np.count_nonzero(uncovered & mask))
                if not gain:
                    continue
                costs = []
                previous = state.pos
                for _, nxt in route:
                    costs.append(distance(previous, pos) + distance(pos, nxt)
                                 - distance(previous, nxt))
                    previous = nxt
                costs.append(distance(previous, pos))
                insertion = min(range(len(costs)), key=costs.__getitem__)
                score = gain**power / (costs[insertion] / 5 + 6 * n_unknown + 1)
                if best is None or score > best[0]:
                    best = (score, i, insertion)
            if best is None:
                break
            _, i, insertion = best
            route.insert(insertion, (-(i + 1), candidates[i]))
            uncovered &= ~masks[i]
        if np.any(uncovered):
            # 极小边缘格由通用探索补足；不能据此宣告完成。
            for index in np.flatnonzero(uncovered)[::max(1, np.count_nonzero(uncovered) // 10)]:
                pos = tuple(self.coverage.points[index])
                route.append((-10000 - int(index), pos))
                uncovered &= ~self.coverage.covered_at(pos)
                if not np.any(uncovered):
                    break
        route = self._route(state.pos, route)
        for _ in range(3):
            changed = False
            for node in list(route):
                if node[0] > 0:
                    continue
                i = route.index(node)
                other_masks = [self.coverage.covered_at(n[1]) for n in route if n != node]
                covered = np.any(other_masks, axis=0) if other_masks else np.zeros_like(remaining)
                critical = remaining & ~covered
                if not np.any(critical):
                    route.remove(node)
                    changed = True
                    continue
                valid = np.flatnonzero(np.all(masks[:, critical], axis=1))
                previous = state.pos if i == 0 else route[i - 1][1]
                nxt = route[i + 1][1] if i + 1 < len(route) else None
                def added(pos):
                    return distance(previous, pos) + (distance(pos, nxt) if nxt else 0)
                best = min(valid, key=lambda j: added(candidates[j]), default=None)
                if best is not None and added(candidates[best]) + 1e-5 < added(node[1]):
                    route[i] = (-(int(best) + 1), candidates[best])
                    changed = True
                refined = self._refine_station(previous, nxt, route[i][1],
                                               self.coverage.points[critical])
                if added(refined) + 1e-5 < added(route[i][1]):
                    route[i] = (route[i][0], refined)
                    changed = True
            if not changed:
                break
            route = self._route(state.pos, route)
        return route

    def _refine_station(self, previous, nxt, current, critical):
        """在保持所有独占覆盖格的前提下，将站位向行进路线移动。

        允许站位是以独占格心为圆心的圆盘交集；每个候选都重新检查全部
        约束，仅接受缩短路径的可行点。这里只优化站位，不改变排除证据。
        """
        points = np.asarray(convex_hull([tuple(p) for p in critical]))
        radius = self.coverage.cover_radius - 1e-5
        radius2 = radius * radius
        start = np.asarray(previous)
        finish = np.asarray(nxt) if nxt is not None else start
        best = np.asarray(current)
        def cost(pos):
            return distance(previous, pos) + (distance(pos, nxt) if nxt else 0)
        def project(pos):
            pos = np.asarray(pos).copy()
            for _ in range(40):
                delta = pos - points
                norm2 = np.sum(delta * delta, axis=1)
                i = int(np.argmax(norm2))
                if norm2[i] <= radius2 + 1e-7:
                    return pos
                pos = points[i] + delta[i] * radius / math.sqrt(norm2[i])
            return None
        for fraction in (0.0, 0.25, 0.5, 0.75, 1.0):
            pos = project(start + fraction * (finish - start))
            if pos is not None and cost(pos) < cost(best):
                best = pos
        for step in (100., 25., 5.):
            for _ in range(4):
                delta = best - start
                grad = delta / max(1e-6, float(np.linalg.norm(delta)))
                if nxt:
                    delta = best - finish
                    grad += delta / max(1e-6, float(np.linalg.norm(delta)))
                pos = project(best - step * grad)
                if pos is None or cost(pos) >= cost(best) - 1e-5:
                    break
                best = pos
        # 当前离散站位始终作为保底，不用尚未验证可行的连续解替换它。
        if np.all(np.sum((points - best)**2, axis=1) <= self.coverage.cover_radius**2):
            return tuple(float(v) for v in best)
        return current

    def step(self, state: State):
        if self.done:
            return Action('done')
        if len(self.cleared | state.cleared) == 16:
            self.done = True
            self.completion_reason = 'cleared_maximum_16'
            return Action('done')
        if self.scan_pending:
            self._start_scan(state, force=self.exploring or not self.tracks)
            self.exploring = False
        while self.scan_queue:
            ch = self.scan_queue.pop(0)
            if ch not in self.cleared and ch not in self.absent:
                return Action('measure', state.pos, ch)

        if self.active_ch in self.tracks and self.active_ch not in self.cleared:
            return self._local_action(self.active_ch, state)
        destination = self._next_destination(state)
        if destination is not None:
            ch, pos = destination
            if ch > 0:
                self.active_ch = ch
                return self._local_action(ch, state)
            mask = self.coverage.covered_at(pos)
            channels = [ch for ch in self._unknown() if self.coverage.gain(ch, mask)]
            self.exploring = True
            self.scan_pending = True
            return Action('measure', pos, channels[0])
        action = self._explore_action(state)
        if action is not None:
            return action
        self.done = True
        self.completion_reason = 'all_channels_cleared_or_covered'
        return Action('done')

    def on_measure(self, state, ch, result, svd):
        if result == 'no_signal':
            self.negative_positions[ch].append(state.pos)
            self.coverage.exclude(ch, state.pos)
            if ch not in self.tracks and self.coverage.absent(ch):
                self.absent.add(ch)
            # 已发现的源不因失去信号而被永久丢弃。保留之前的有效测向。
            if ch == self.active_ch and ch in self.tracks:
                track = self.tracks[ch]
                poly = exclude_disk(track.poly, state.pos)
                if poly:
                    poly = p1.hpi_intersect(poly, reception_halfplanes(
                        [p for p, _ in track.records], [state.pos]))
                if not poly:
                    raise RuntimeError(f'频道 {ch} 的无信号响应与先前测向矛盾')
                track.poly = poly
                track.center, track.radius = enclosing_circle(poly)
                track.estimate = centroid(poly)
            return
        if result == 'near':
            track = self.tracks.setdefault(ch, Track())
            track.near = state.pos
            track.center, track.radius = state.pos, 5.0
            track.estimate = state.pos
            track.probe_after_fail = False
            return
        if result != 'direction' or svd is None or not math.isfinite(svd):
            raise ValueError(f'非法测向响应: {result}, {svd}')
        track = self.tracks.setdefault(ch, Track())
        if any(distance(state.pos, p) < 1e-5 for p, _ in track.records):
            return
        track.records.append((state.pos, svd))
        if track.poly:
            poly = track.poly
        else:
            poly = p1.clip_polygon_by_disk(
                [(-1800., -1800.), (1800., -1800.),
                 (1800., 1800.), (-1800., 1800.)], (0., 0.), ARENA_R, 64)
        poly = p1.hpi_intersect(poly, p1.make_sector_halfplanes(state.pos, svd, EPS_DEG))
        if poly:
            poly = p1.clip_polygon_by_disk(poly, state.pos, RECEIVE_MAX, 32)
        if poly and self.negative_positions[ch]:
            poly = p1.hpi_intersect(poly, reception_halfplanes(
                [p for p, _ in track.records], self.negative_positions[ch]))
        if poly:
            for pos in self.negative_positions[ch]:
                poly = exclude_disk(poly, pos)
                if not poly:
                    break
        if not poly:
            self.geometry_conflicts += 1
            raise RuntimeError(f'频道 {ch} 的有效测向约束无交集，需检查接口或误差模型')
        track.poly = poly
        track.center, track.radius = enclosing_circle(poly)
        track.estimate = centroid(poly)
        track.probe_after_fail = False

    def on_clear(self, state, ch, success):
        if success:
            self.cleared.add(ch)
            state.cleared.add(ch)
            self.active_ch = None
            self.scan_pending = True
        else:
            track = self.tracks[ch]
            track.failures += 1
            track.near = None
            # 不重试同一个清除点，也不把未清除源标记为完成。
            if track.failures > 3:
                raise RuntimeError(f'频道 {ch} 清除连续失败，检查真实设备规则')
            self.active_ch = ch
            track.probe_after_fail = True


def make_strategy(name='AdaptiveV8'):
    if name not in ('AdaptiveV8', 'v8'):
        raise ValueError(f'未知 v8 策略: {name}')
    return AdaptiveV8()
\end{lstlisting}
%</paper-block>

%<paper-block id="appendix.p012" type="paragraph">
\noindent\textbf{文件：}\texttt{02\_\allowbreak{}最终代码/\allowbreak{}第三问/\allowbreak{}strategy\_\allowbreak{}fast.py}\\
\textbf{功能：}问题三正式链路使用的搜索参数基类
%</paper-block>
%<paper-block id="appendix.code010" type="code">
\begin{lstlisting}[style=appendixcode]
"""演练候选参数；保留连续区域、逐频道覆盖和失败恢复逻辑。"""
from strategy_v8 import AdaptiveV8

class FastScanP3(AdaptiveV8):
    name='FastScanP3'
    def __init__(self):
        super().__init__(cell_size=60.0,scan_fraction=0.20,probe_radius=80.0)

class CautiousP3(AdaptiveV8):
    name='CautiousP3'
    def __init__(self):
        super().__init__(probe_radius=35.0)

class CompactP3(AdaptiveV8):
    name='CompactP3'
    def __init__(self):
        super().__init__(cell_size=60.0,scan_fraction=0.20,probe_radius=40.0)
\end{lstlisting}
%</paper-block>

%<paper-block id="appendix.p013" type="paragraph">
\noindent\textbf{文件：}\texttt{02\_\allowbreak{}最终代码/\allowbreak{}第三问/\allowbreak{}strategy\_\allowbreak{}night.py}\\
\textbf{功能：}问题三连续区域覆盖证书与搜索基类
%</paper-block>
%<paper-block id="appendix.code011" type="code">
\begin{lstlisting}[style=appendixcode]
"""P3 nightly candidates; observations only, no simulator internals or known source count."""
from functools import lru_cache
import math
import numpy as np
from shapely.geometry import Polygon
from shapely.ops import unary_union
from strategy_fast import FastScanP3
from strategy_v8 import ARENA_R, RECEIVE_MIN, distance

POLYGON_SIDES=128
MARGIN_M=0.01
UNIT=tuple((math.cos(2*math.pi*i/POLYGON_SIDES),math.sin(2*math.pi*i/POLYGON_SIDES))
           for i in range(POLYGON_SIDES))
# The target circle is inside this polygon; exclusion polygons are inside 1000m disks.
OUTER_TARGET=Polygon([(x*(ARENA_R+MARGIN_M)/math.cos(math.pi/POLYGON_SIDES),
                       y*(ARENA_R+MARGIN_M)/math.cos(math.pi/POLYGON_SIDES)) for x,y in UNIT])

def polygon_coverage_certificate(negative_positions):
    if not negative_positions:return False
    radius=RECEIVE_MIN-MARGIN_M
    disks=[Polygon([(cx+radius*x,cy+radius*y) for x,y in UNIT]) for cx,cy in negative_positions]
    return bool(unary_union(disks).covers(OUTER_TARGET))

class CachedScanP3(FastScanP3):
    name='CachedScanP3'
    def __init__(self):
        super().__init__()
        original=self.coverage.covered_at
        self._cached_mask=lru_cache(maxsize=2048)(original)
        self.coverage.covered_at=lambda pos:self._cached_mask((float(pos[0]),float(pos[1])))

class CoverP3(CachedScanP3):
    name='CoverP3'
    def __init__(self):
        super().__init__()
        self.coverage_certificates={}

    def on_measure(self,state,ch,result,svd):
        super().on_measure(state,ch,result,svd)
        if result!='no_signal' or ch in self.tracks or ch in self.absent:return
        # Ordinary whole-cell coverage remains the primary test. Only certify the small tail.
        if np.count_nonzero(self.coverage.remaining[ch])>len(self.coverage.points)*0.12:return
        positions=self.negative_positions[ch]
        if polygon_coverage_certificate(positions):
            self.absent.add(ch)
            self.coverage.remaining[ch][:]=False
            self.coverage_certificates[ch]=list(positions)

    def diagnostics(self):
        info=self._cached_mask.cache_info()
        return dict(mask_cache_hits=info.hits,mask_cache_misses=info.misses,
            coverage_certificate_method='outer-target128 covered by union of inner-exclusion128, 0.01m margins',
            coverage_certificates=self.coverage_certificates)

class HopP3(CoverP3):
    name='HopP3'
    def __init__(self):
        super().__init__()
        self.first_hop=500.0
        self.lateral=100.0

class MidpointP3(CoverP3):
    name='MidpointP3'
    def __init__(self):
        super().__init__()
        self.first_hop=750.0
        self.lateral=100.0

class RouteP3(MidpointP3):
    name='RouteP3'
    @staticmethod
    def _route(start,nodes):
        route=FastScanP3._route(start,nodes)
        # A single-node relocation can improve a 2-opt local minimum. Keep the base route as fallback.
        for _ in range(4):
            best_delta=-1e-6;best=None
            for i,node in enumerate(route):
                prev=start if i==0 else route[i-1][1]
                nxt=route[i+1][1] if i+1<len(route) else None
                removed=distance(prev,node[1])
                if nxt is not None:removed+=distance(node[1],nxt)-distance(prev,nxt)
                rest=route[:i]+route[i+1:]
                for j in range(len(rest)+1):
                    if j==i:continue
                    before=start if j==0 else rest[j-1][1]
                    after=rest[j][1] if j<len(rest) else None
                    added=distance(before,node[1])
                    if after is not None:added+=distance(node[1],after)-distance(before,after)
                    if added-removed<best_delta:best_delta=added-removed;best=(i,j)
            if best is None:break
            node=route.pop(best[0]);route.insert(best[1],node)
        return route
\end{lstlisting}
%</paper-block>

%<paper-block id="appendix.p014" type="paragraph">
\noindent\textbf{文件：}\texttt{02\_\allowbreak{}最终代码/\allowbreak{}第三问/\allowbreak{}strategy\_\allowbreak{}transit.py}\\
\textbf{功能：}问题三剩余区域更新与排空判定
%</paper-block>
%<paper-block id="appendix.code012" type="code">
\begin{lstlisting}[style=appendixcode]
"""P3 practice: ResidualP3 is selected; transit/tail variants remain offline experiments."""
from functools import lru_cache
import numpy as np
from shapely import box, intersects
from shapely.geometry import Polygon
from shapely.ops import unary_union

from strategy import Action
from strategy_night import HopP3, OUTER_TARGET, UNIT, MARGIN_M
from strategy_v8 import RECEIVE_MIN, distance


class ResidualP3(HopP3):
    name = 'ResidualP3'
    residual_fraction = 0.35

    def __init__(self):
        super().__init__()
        # Coverage was constructed with exactly 60m cells by FastScanP3.
        points = self.coverage.points
        self._cells = box(points[:, 0] - 30., points[:, 1] - 30.,
                          points[:, 0] + 30., points[:, 1] + 30.)
        self._residual_mask = lru_cache(maxsize=128)(self._compute_residual_mask)
        self.residual_updates = 0
        self.residual_channels = set()
        self.residual_certificates = {}

    def _compute_residual_mask(self, positions):
        radius = RECEIVE_MIN - MARGIN_M
        exclusions = [Polygon([(cx + radius*x, cy + radius*y) for x, y in UNIT])
                      for cx, cy in positions]
        residual = OUTER_TARGET.difference(unary_union(exclusions))
        # Closed cell intersection also retains touching cells and a small outward guard.
        return np.asarray(intersects(residual.buffer(0.0001), self._cells), dtype=bool)

    def on_measure(self, state, ch, result, svd):
        super().on_measure(state, ch, result, svd)
        if result == 'no_signal' and ch in self.absent and ch in self.residual_channels:
            self.residual_certificates[ch] = list(self.negative_positions[ch])
        if result != 'no_signal' or ch in self.tracks or ch in self.absent:
            return
        remaining = self.coverage.remaining[ch]
        if np.count_nonzero(remaining) > len(remaining) * self.residual_fraction:
            return
        positions = tuple(tuple(p) for p in self.negative_positions[ch])
        residual_mask = self._residual_mask(positions)
        before = np.count_nonzero(remaining)
        remaining &= residual_mask
        changed = np.count_nonzero(remaining) < before
        self.residual_updates += int(changed)
        if changed:
            self.residual_channels.add(ch)
        if not np.any(remaining):
            self.absent.add(ch)
            self.coverage_certificates[ch] = list(positions)
            self.residual_certificates[ch] = list(positions)

    def diagnostics(self):
        result = super().diagnostics()
        result.update(residual_updates=self.residual_updates,
                      residual_method='whole 60m cells disjoint from outward-guarded polygon difference',
                      residual_cache_hits=self._residual_mask.cache_info().hits,
                      residual_certificates=self.residual_certificates)
        return result


class TransitP3(HopP3):
    name = 'TransitP3'
    transit_fraction = 0.12
    minimum_leg = 700.

    def __init__(self):
        super().__init__()
        self._transit_pending = None
        self._transit_channels = []
        self.transit_stops = []

    def step(self, state):
        if self._transit_pending is not None:
            while self._transit_channels:
                ch = self._transit_channels.pop(0)
                if ch not in self.cleared and ch not in self.absent and ch not in self.tracks:
                    return Action('measure', state.pos, ch)
            action = self._transit_pending
            self._transit_pending = None
            return action
        action = super().step(state)
        unknown = self._unknown()
        if (action.kind not in ('measure', 'clear') or not unknown
                or distance(state.pos, action.pos) < self.minimum_leg):
            return action
        endpoint = self.coverage.covered_at(action.pos)
        start = np.asarray(state.pos)
        finish = np.asarray(action.pos)
        best = None
        for fraction in (0.25, 0.5, 0.75):
            pos = tuple(float(v) for v in start + fraction * (finish-start))
            mask = self.coverage.covered_at(pos) & ~endpoint
            channels, gain = [], 0
            for ch in unknown:
                extra = self.coverage.gain(ch, mask)
                threshold = max(1, int(min(len(mask)*self.transit_fraction,
                                          np.count_nonzero(self.coverage.remaining[ch])*0.8)))
                if extra >= threshold:
                    channels.append(ch)
                    gain += extra
            if not channels:
                continue
            score = gain/(6*len(channels)+1)
            if best is None or score > best[0]:
                best = score, pos, channels
        if best is None:
            return action
        _, pos, channels = best
        self._transit_pending = action
        self._transit_channels = sorted(channels, key=lambda ch: (ch != state.ch, ch))
        self.transit_stops.append(dict(start=state.pos, finish=action.pos, pos=pos,
                                       channels=list(self._transit_channels)))
        return Action('measure', pos, self._transit_channels.pop(0))

    def diagnostics(self):
        result = super().diagnostics()
        result['transit_stops'] = self.transit_stops
        return result


class TransitResidualP3(TransitP3, ResidualP3):
    name = 'TransitResidualP3'


class TransitSparseP3(TransitP3):
    name = 'TransitSparseP3'
    transit_fraction = 0.20
    minimum_leg = 900.


class TailResidualP3(ResidualP3):
    name = 'TailResidualP3'
    residual_fraction = 0.12


class TransitEagerP3(TransitP3):
    name = 'TransitEagerP3'
    transit_fraction = 0.025
    minimum_leg = 500.


class TailTransitP3(TransitEagerP3, TailResidualP3):
    name = 'TailTransitP3'
\end{lstlisting}
%</paper-block>

%<paper-block id="appendix.p015" type="paragraph">
\noindent\textbf{文件：}\texttt{02\_\allowbreak{}最终代码/\allowbreak{}第三问/\allowbreak{}strategy\_\allowbreak{}joint.py}\\
\textbf{功能：}问题三补测点成本评估与选择
%</paper-block>
%<paper-block id="appendix.code013" type="code">
\begin{lstlisting}[style=appendixcode]
"""P3 observation-only joint localization/routing experiments."""
import math

from strategy import Action, State, p1
from strategy_transit import ResidualP3
from strategy_v8 import EPS_DEG, distance, centroid, enclosing_circle


class ActionRouteP3(ResidualP3):
    name = 'ActionRouteP3'

    def _visit_cost(self, previous, node):
        ch, finish = node
        if ch <= 0:
            return distance(previous, finish)
        # The first localization measurement may not lie on the route to the estimate.
        action = self._local_action(ch, State(pos=previous, ch=ch))
        return distance(previous, action.pos) + distance(action.pos, finish)

    def _route(self, start, nodes):
        if len(nodes) < 2:
            return list(nodes)
        n = len(nodes)
        positions = [node[1] for node in nodes] + [start]
        costs = [[self._visit_cost(pos, node) if i != j else 0.
                  for j, node in enumerate(nodes)] for i, pos in enumerate(positions)]
        remaining = set(range(n))
        order, previous = [], n
        while remaining:
            chosen = min(remaining, key=lambda j: (costs[previous][j], j))
            order.append(chosen)
            remaining.remove(chosen)
            previous = chosen
        # Directed 2-opt: account for every reversed internal arc, not only endpoint edges.
        for _ in range(12):
            best = (-1e-6, None)
            prefix = [0.]
            for a, b in zip(order, order[1:]):
                prefix.append(prefix[-1] + costs[b][a] - costs[a][b])
            for i in range(n-1):
                prev = n if i == 0 else order[i-1]
                for j in range(i+1, n):
                    delta = costs[prev][order[j]] - costs[prev][order[i]] + prefix[j]-prefix[i]
                    if j+1 < n:
                        nxt = order[j+1]
                        delta += costs[order[i]][nxt] - costs[order[j]][nxt]
                    if delta < best[0]:
                        best = delta, (i, j)
            if best[1] is None:
                break
            i, j = best[1]
            order[i:j+1] = reversed(order[i:j+1])
        return [nodes[i] for i in order]

    def _cover_route(self, state, known, remaining, n_unknown):
        routes = [self._build_cover_route(state, known, remaining, n_unknown, power)
                  for power in (0.6, 1.0, 1.6)]
        def cost(route):
            previous, total = state.pos, 0.
            for node in route:
                total += self._visit_cost(previous, node)/5 + (6*n_unknown if node[0] < 0 else 0.)
                previous = node[1]
            return total
        route = min(routes, key=cost)
        return route[0] if route else None


class SurveyP3(ResidualP3):
    name = 'SurveyP3'

    def __init__(self):
        super().__init__()
        self._survey_measure = None
        self.local_survey_stops = 0

    def step(self, state):
        action = super().step(state)
        self._survey_measure = None
        if (action.kind == 'measure' and action.ch == self.active_ch
                and not self.scan_queue and distance(state.pos, action.pos) > 100.):
            self._survey_measure = action.ch
        return action

    def on_measure(self, state, ch, result, svd):
        super().on_measure(state, ch, result, svd)
        if self._survey_measure == ch:
            self.scan_pending = True
            self.local_survey_stops += 1
        self._survey_measure = None

    def diagnostics(self):
        result = super().diagnostics()
        result['local_survey_stops'] = self.local_survey_stops
        return result


class JointP3(SurveyP3, ActionRouteP3):
    name = 'JointP3'


class ProbePlanP3(ResidualP3):
    """One-step expected-cost probe selection; sampled beliefs are not source truth."""
    name = 'ProbePlanP3'

    def __init__(self):
        super().__init__()
        self.probe_plans = []

    @staticmethod
    def _belief_samples(poly):
        center = centroid(poly)
        samples = []
        for a, b in zip(poly, poly[1:]+poly[:1]):
            weight = abs((a[0]-center[0])*(b[1]-center[1]) - (a[1]-center[1])*(b[0]-center[0]))
            if weight > 1e-9:
                samples.append((((center[0]+a[0]+b[0])/3,
                                 (center[1]+a[1]+b[1])/3), weight))
        total = sum(w for _, w in samples)
        return [(point, w/total) for point, w in samples]

    @staticmethod
    def _probe_cost(current, point, poly, samples):
        value = distance(current, point)/5 + 5.
        for truth, weight in samples:
            d = distance(point, truth)
            if d < 5.:
                value += weight*5.
                continue
            # Conservative planning penalty for a possible no-signal response at the old probe.
            if d > 1000.:
                value += weight*(d/5 + 30.)
                continue
            angle = math.degrees(math.atan2(truth[1]-point[1], truth[0]-point[0])) % 360
            for error, ew in ((-1., .25), (0., .5), (1., .25)):
                observed = round((angle+error) % 360, 2) % 360
                posterior = p1.hpi_intersect(poly, p1.make_sector_halfplanes(point, observed, EPS_DEG))
                if not posterior:
                    return math.inf
                center, radius = enclosing_circle(posterior)
                if radius <= 20. - 1e-5:
                    continuation = max(0., distance(point, center)-(20-radius-1e-5))/5+5.
                else:
                    estimate = centroid(posterior)
                    miss = distance(estimate, truth)
                    continuation = distance(point, estimate)/5
                    if radius <= 80. and miss <= 20.:
                        continuation += 5.
                    else:
                        continuation += 10. + max(0., miss-20.)/5
                        if radius <= 80.:
                            continuation += 3.
                value += weight*ew*continuation
        return value

    def step(self, state):
        action = super().step(state)
        if action.kind != 'measure' or action.ch != self.active_ch:
            return action
        track = self.tracks.get(action.ch)
        if (track is None or len(track.records) != 1 or track.probe_after_fail
                or distance(state.pos, action.pos) < 1. or not track.poly):
            return action
        old, angle = track.records[0]
        theta = math.radians(angle)
        ux, uy = math.cos(theta), math.sin(theta)
        length = distance(old, track.estimate)
        samples = self._belief_samples(track.poly)
        best_point = action.pos
        baseline = self._probe_cost(state.pos, action.pos, track.poly, samples)
        best_cost = baseline
        for fraction in (.45, .65, .85, 1.):
            for lateral in (50., 100., 200.):
                options = [(old[0]+fraction*length*ux-sign*lateral*uy,
                            old[1]+fraction*length*uy+sign*lateral*ux) for sign in (-1, 1)]
                point = min(options, key=lambda p: distance(state.pos, p))
                if max(distance(point, vertex) for vertex in track.poly) > 999.99:
                    continue
                if any(distance(point, p) < 10. for p, _ in track.records):
                    continue
                score = self._probe_cost(state.pos, point, track.poly, samples)
                if score < best_cost-1e-6:
                    best_point, best_cost = point, score
        if best_point != action.pos:
            self.probe_plans.append(dict(ch=action.ch, old=action.pos, chosen=best_point,
                                         predicted_old_s=baseline, predicted_new_s=best_cost))
            return Action('measure', best_point, action.ch)
        return action

    def diagnostics(self):
        result = super().diagnostics()
        result['probe_plans'] = self.probe_plans
        return result


class PlanJointP3(ProbePlanP3, JointP3):
    name = 'PlanJointP3'
\end{lstlisting}
%</paper-block>

%<paper-block id="appendix.p016" type="paragraph">
\noindent\textbf{文件：}\texttt{02\_\allowbreak{}最终代码/\allowbreak{}第三问/\allowbreak{}strategy\_\allowbreak{}task.py}\\
\textbf{功能：}问题三最终 OpticalTaskP3：数量上界与有限光学覆盖
%</paper-block>
%<paper-block id="appendix.code014" type="code">
\begin{lstlisting}[style=appendixcode]
"""P3 task-cost candidates: certify all sources, then minimize actual action time.

Uses only observations. Optical strips cover a continuous feasible polygon with
20m disks, not sampled targets. Belief quadrature ranks plans; it never certifies them.
"""
import math

from strategy import Action, State
from strategy_joint import ProbePlanP3
from strategy_v8 import (CLEAR_R, EPS_DEG, centroid, distance, enclosing_circle,
                         exclude_disk, p1)


def optical_strip(poly, current, samples, maximum_shots=6, axis=None):
    """A bounded continuous cover of a polygon's oriented bounding rectangle.

    At normal half-width w, each radius-r disk covers a rectangle of axial
    half-length sqrt(r*r-w*w). Adjacent rectangles tile the bounding rectangle.
    A 0.01m inward guard separates execution from the mathematical 20m boundary.
    """
    if len(poly)<3:
        return None
    if axis is None:
        _,a,b=p1.polygon_diameter(poly)
        # A longest diagonal can make a narrow rectangle appear too wide.
        # Retain it as a fallback, and compare all polygon-edge orientations.
        axes=[(a,b)]+list(zip(poly,poly[1:]+poly[:1]))
        plans=[optical_strip(poly,current,samples,maximum_shots,direction) for direction in axes]
        feasible=[p for p in plans if p is not None]
        return min(feasible,key=lambda p:p['expected_s']) if feasible else None
    a,b=axis
    length=distance(a,b)
    if length<1e-9:
        return None
    ux,uy=(b[0]-a[0])/length,(b[1]-a[1])/length
    vx,vy=-uy,ux
    along=[x*ux+y*uy for x,y in poly]
    across=[x*vx+y*vy for x,y in poly]
    low,high=min(along),max(along);middle=(min(across)+max(across))/2
    half_width=(max(across)-min(across))/2
    r=CLEAR_R-.01
    if half_width>=r:
        return None
    half_length=math.sqrt(r*r-half_width*half_width)
    n=max(1,math.ceil((high-low)/(2*half_length)))
    if n>maximum_shots:
        return None
    centers=[(low+high)/2] if n==1 else [low+half_length+i*(high-low-2*half_length)/(n-1) for i in range(n)]
    points=[(s*ux+middle*vx,s*uy+middle*vy) for s in centers]
    def expected(order):
        result=0.
        for target,weight in samples:
            prev=current;cost=0.;found=False
            for point in order:
                cost+=distance(prev,point)/5
                if distance(target,point)<=CLEAR_R:
                    cost+=5;found=True;break
                cost+=3;prev=point
            if not found:
                return math.inf
            result+=weight*cost
        return result
    reverse=list(reversed(points))
    forward_cost,reverse_cost=expected(points),expected(reverse)
    if reverse_cost<forward_cost:
        points=reverse;forward_cost=reverse_cost
    return dict(points=points, expected_s=forward_cost, half_width=half_width,
                half_length=half_length, axial_bounds=[low,high], radius_guard=r)


class CardinalityP3(ProbePlanP3):
    name='CardinalityP3'

    def __init__(self):
        super().__init__()
        self.cardinality_certificates=[]

    def _retire_unknown_at_maximum(self):
        # The count is inferred from distinct positively observed channels, not the
        # simulator's hidden N. Finding 16 is enough; clearing all 16 is not needed.
        known=set(self.tracks)|self.cleared
        if len(known)>16:
            raise RuntimeError('More observed source channels than the problem permits')
        if len(known)==16:
            unknown=[ch for ch in range(1,21) if ch not in known and ch not in self.absent]
            if unknown:
                self.absent.update(unknown)
                self.cardinality_certificates.append(dict(observed_channels=sorted(known),absent_channels=unknown))

    def on_measure(self,state,ch,result,svd):
        super().on_measure(state,ch,result,svd)
        self._retire_unknown_at_maximum()

    def diagnostics(self):
        result=super().diagnostics()
        result['cardinality_certificates']=self.cardinality_certificates
        return result


class OpticalTaskP3(CardinalityP3):
    name='OpticalTaskP3'
    optical_gain_margin_s=1.
    maximum_optical_shots=6

    def __init__(self):
        super().__init__()
        self.optical_pending=None
        self.optical_plans=[]

    @staticmethod
    def _old_action_cost(action,poly,samples,current):
        if action.kind=='measure':
            return ProbePlanP3._probe_cost(current,action.pos,poly,samples)
        if action.kind!='clear':
            return math.inf
        missed=[(point,w) for point,w in samples if distance(point,action.pos)>CLEAR_R]
        probability=sum(w for _,w in missed)
        cost=distance(current,action.pos)/5+(1-probability)*5
        if probability:
            conditional=[(point,w/probability) for point,w in missed]
            cost+=probability*(3+ProbePlanP3._probe_cost(action.pos,action.pos,poly,conditional))
        return cost

    def step(self,state):
        if self.optical_pending is not None:
            plan=self.optical_pending
            if not plan['points']:
                raise RuntimeError('Continuous optical cover exhausted without clearing: inconsistent observations')
            return Action('clear',plan['points'][0],plan['ch'])
        action=super().step(state)
        if action.kind not in ('clear','measure') or action.ch!=self.active_ch:
            return action
        track=self.tracks.get(action.ch)
        if (track is None or track.near is not None or len(track.records)<2
                or track.probe_after_fail or track.radius<=CLEAR_R or not track.poly):
            return action
        samples=self._belief_samples(track.poly)
        plan=optical_strip(track.poly,state.pos,samples,self.maximum_optical_shots)
        if plan is None:
            return action
        old_cost=self._old_action_cost(action,track.poly,samples,state.pos)
        if plan['expected_s']+self.optical_gain_margin_s>=old_cost:
            return action
        self.optical_pending=dict(ch=action.ch,points=list(plan['points']))
        self.optical_plans.append(dict(ch=action.ch,**plan,previous_expected_s=old_cost,
                                       feasible_polygon=list(track.poly)))
        return Action('clear',plan['points'][0],action.ch)

    def on_clear(self,state,ch,success):
        if self.optical_pending is None:
            return super().on_clear(state,ch,success)
        if ch!=self.optical_pending['ch']:
            raise RuntimeError('Optical plan channel mismatch')
        if success:
            self.optical_pending=None
            return super().on_clear(state,ch,True)
        self.optical_pending['points'].pop(0)
        track=self.tracks[ch]
        track.failures+=1
        track.near=None
        # A failed optical action is useful information. Preserve a conservative
        # hull; the precomputed continuous cover stays valid for any subset.
        poly=exclude_disk(track.poly,state.pos,CLEAR_R)
        if not poly:
            raise RuntimeError('Optical failure contradicts the feasible source region')
        track.poly=poly;track.center,track.radius=enclosing_circle(poly)
        track.estimate=centroid(poly);track.probe_after_fail=False

    def diagnostics(self):
        result=super().diagnostics()
        result['optical_plans']=self.optical_plans
        return result


class OpticalProbeP3(OpticalTaskP3):
    """Also price the terminal optical plan when choosing the second bearing point."""
    name='OpticalProbeP3'

    @staticmethod
    def _probe_cost(current,point,poly,samples):
        value=distance(current,point)/5+5.
        for hypothesis,weight in samples:
            d=distance(point,hypothesis)
            if d<=5:
                value+=weight*5;continue
            if d>1000:
                value+=weight*(d/5+30);continue
            angle=math.degrees(math.atan2(hypothesis[1]-point[1],hypothesis[0]-point[0]))%360
            for error,ew in ((-1.,.25),(0.,.5),(1.,.25)):
                observed=round((angle+error)%360,2)%360
                posterior=p1.hpi_intersect(poly,p1.make_sector_halfplanes(point,observed,EPS_DEG))
                if not posterior:return math.inf
                center,radius=enclosing_circle(posterior)
                if radius<=20-1e-5:
                    continuation=max(0.,distance(point,center)-(20-radius-1e-5))/5+5
                else:
                    estimate=centroid(posterior);miss=distance(estimate,hypothesis)
                    continuation=distance(point,estimate)/5
                    if radius<=80 and miss<=20:
                        continuation+=5
                    else:
                        continuation+=10+max(0.,miss-20)/5+(3 if radius<=80 else 0)
                    plan=optical_strip(posterior,point,ProbePlanP3._belief_samples(posterior))
                    if plan is not None:
                        continuation=min(continuation,plan['expected_s'])
                value+=weight*ew*continuation
        return value


class SharedBearingP3(CardinalityP3):
    """Buy a second channel's bearing at an already paid-for localization stop.

    Expected avoided future localization cost must exceed 5s detection + up to
    1s switching, plus a 2s planning margin. No extra movement is inserted.
    """
    name='SharedBearingP3'

    def __init__(self):
        super().__init__()
        self.shared_queue=[]
        self.shared_trigger=None
        self.shared_measurements=[]

    @staticmethod
    def _future_finish_cost(previous,poly,hypothesis):
        center,radius=enclosing_circle(poly)
        if radius<=20-1e-5:
            return max(0.,distance(previous,center)-(20-radius-1e-5))/5+5
        estimate=centroid(poly);error=distance(estimate,hypothesis)
        cost=distance(previous,estimate)/5
        if radius<=80 and error<=20:return cost+5
        return cost+10+max(0.,error-20)/5+(3 if radius<=80 else 0)

    def _sharing_candidates(self,state):
        current_track=self.tracks.get(self.active_ch)
        previous=current_track.estimate if current_track else state.pos
        nodes=[(ch,t.estimate) for ch,t in self.tracks.items() if ch not in self.cleared and ch!=self.active_ch]
        candidates=[]
        for ch,estimate in self._route(previous,nodes):
            track=self.tracks[ch]
            origin=previous;previous=estimate
            if (track.near is not None or track.probe_after_fail or track.radius<=20
                    or not track.records or not track.poly):continue
            if min(distance(state.pos,p) for p,_ in track.records)<100:continue
            if max(distance(state.pos,v) for v in track.poly)>999.99:continue
            samples=self._belief_samples(track.poly)
            old=self._local_action(ch,State(pos=origin,ch=ch))
            before=OpticalTaskP3._old_action_cost(old,track.poly,samples,origin)
            after=0.
            for hypothesis,weight in samples:
                if distance(state.pos,hypothesis)<=5:
                    after+=weight*(distance(origin,state.pos)/5+5);continue
                angle=math.degrees(math.atan2(hypothesis[1]-state.pos[1],hypothesis[0]-state.pos[0]))%360
                for error,ew in ((-1.,.25),(0.,.5),(1.,.25)):
                    observed=round((angle+error)%360,2)%360
                    posterior=p1.hpi_intersect(track.poly,p1.make_sector_halfplanes(state.pos,observed,EPS_DEG))
                    if not posterior:after=math.inf;break
                    after+=weight*ew*self._future_finish_cost(origin,posterior,hypothesis)
            saved=before-after-6.
            if saved>2.:
                candidates.append((saved,ch))
        return sorted(candidates,reverse=True)

    def step(self,state):
        self.shared_trigger=None
        while self.shared_queue:
            ch=self.shared_queue.pop(0)
            if ch in self.cleared or ch in self.absent:continue
            return Action('measure',state.pos,ch)
        action=super().step(state)
        if (action.kind=='measure' and action.ch==self.active_ch
                and distance(state.pos,action.pos)>100):
            self.shared_trigger=action.ch
        return action

    def on_measure(self,state,ch,result,svd):
        super().on_measure(state,ch,result,svd)
        if self.shared_trigger==ch:
            candidates=self._sharing_candidates(state)
            self.shared_queue=[other for _,other in candidates]
            if candidates:
                self.shared_measurements.append(dict(pos=state.pos,channels=list(self.shared_queue),
                                                      predicted_net_savings=[s for s,_ in candidates]))
        self.shared_trigger=None

    def diagnostics(self):
        result=super().diagnostics()
        result['shared_bearing_stops']=self.shared_measurements
        return result


class SharedOpticalP3(SharedBearingP3,OpticalTaskP3):
    name='SharedOpticalP3'


class BoundedOpticalP3(OpticalTaskP3):
    """Only short optical finishes; require a larger cost margin than one action tick."""
    name='BoundedOpticalP3'
    maximum_optical_shots=3
    optical_gain_margin_s=3.
\end{lstlisting}
%</paper-block>

%<paper-block id="appendix.h005" type="heading">
\subsection*{B.3 第四问最终策略程序}
%</paper-block>
%<paper-block id="appendix.p017" type="paragraph">
\noindent\textbf{文件：}\texttt{02\_\allowbreak{}最终代码/\allowbreak{}第四问/\allowbreak{}strategy.py}\\
\textbf{功能：}问题四动作、状态及基础定位接口
%</paper-block>
%<paper-block id="appendix.code015" type="code">
\begin{lstlisting}[style=appendixcode]
# -*- coding: utf-8 -*-
# strategy.py  (v2 重写版)
# 第三问：4 套定位清除策略（统一接口）
#
# 设计原则：
#   - 每个频道 ch 有状态：NEW / M1 / READY / CLEARED / SKIP
#   - 策略 step() 在外层 cursor 上推进，对当前 ch 调用 _next_action_for_ch(ch)
#   - 状态机不会回退：M1 → 一定要补第 2 个 bearing，然后切到 READY（再试定位）
#   - 失败的 clear 不会无限重试：超过 N 次失败 → SKIP
#
# 4 套策略的区别 = 第二个测向点的位置选取 + 清除前是否再移动补测：
#   1) StaticScan：固定在 (1000, 0)，第 2 点 = 沿 θ 走 800 m
#   2) Patrol：4 个站点轮换
#   3) BlinkHop：第 1 点随机在 (0,0) 起步，每测到 direction 沿 θ 跳 800 m
#   4) Hybrid：早期闪烁跳跃广撒网，后期 StaticScan 精定位
#
# 复用 Q1 v4 的 solve_problem_1 / localization_polygon。

import math
import os
import importlib.util as _ilu
from dataclasses import dataclass, field
from typing import List, Tuple, Optional, Dict, Set

# 复用 Q1 v4 算法 —— 优先用内联版本 problem1_v4_inline.py（自包含），
# 没有内联文件时再用外部多级 fallback 找 problem1_v4.py。
# 这样 v5 文件夹可以单独复制到任何机器，不需要 Q1 v4 源码。
def _locate_problem1_v4():
    """优先内联；找不到再 fallback 找外部 problem1_v4.py。
    返回 (module, source_label)。"""
    here = os.path.dirname(os.path.abspath(__file__))
    # 1) 内联版本（v5 自带）
    inline = os.path.join(here, 'problem1_v4_inline.py')
    if os.path.isfile(inline):
        return inline, 'inline'

    # 2) 外部多级 fallback
    fname = 'problem1_v4.py'
    candidates = []
    # 2a) 与本文件相对
    candidates.append(os.path.normpath(os.path.join(
        here, '..', '..', '第一问', 'v4', fname)))
    # 2b) 从 cwd 向上找
    cwd = os.getcwd()
    cur = cwd
    for _ in range(8):
        cand = os.path.join(cur, '第一问', 'v4', fname)
        if os.path.isfile(cand):
            candidates.append(os.path.abspath(cand))
            break
        parent = os.path.dirname(cur)
        if parent == cur:
            break
        cur = parent
    # 2c) 环境变量
    env = os.environ.get('CUMCM2026_ROOT')
    if env:
        candidates.append(os.path.normpath(os.path.join(env, '第一问', 'v4', fname)))
    # 2d) 常见路径
    common = [
        r'C:\Users\LENOVO\Desktop\CUMCM2026Problems',
        r'D:\CUMCM2026Problems',
        r'/mnt/c/Users/LENOVO/Desktop/CUMCM2026Problems',
        os.path.expanduser('~/Desktop/CUMCM2026Problems'),
        os.path.expanduser('~/CUMCM2026Problems'),
    ]
    for c in common:
        candidates.append(os.path.join(c, '第一问', 'v4', fname))

    for c in candidates:
        if os.path.isfile(c):
            return os.path.abspath(c), 'external'

    raise FileNotFoundError(
        '找不到 Q1 v4 算法源。请确认 v5 文件夹里有 problem1_v4_inline.py，'
        '或设置 CUMCM2026_ROOT 指向 CUMCM2026Problems 根目录。\n'
        '尝试过的路径：\n  ' + '\n  '.join(candidates))


_FIRST_FILE, _SOURCE = _locate_problem1_v4()
_spec = _ilu.spec_from_file_location('problem1_v4', _FIRST_FILE)
p1 = _ilu.module_from_spec(_spec)
_spec.loader.exec_module(p1)
print(f'[strategy] Q1 v4 算法源 = {_SOURCE}: {os.path.basename(_FIRST_FILE)}')

Vec = Tuple[float, float]
CHANNELS = list(range(1, 21))
ARENA_R = 1800.0
EPS_DEG = 1.0
CLEAR_R = 20.0
SPEED = 5.0
HOP_DIST = 800.0           # 闪烁跳跃：沿 direction 走多远
STATIC_POS: Vec = (1000.0, 0.0)
PATROL_POS: List[Vec] = [
    (1000.0, 0.0), (0.0, 1000.0), (-1000.0, 0.0), (0.0, -1000.0)
]
MAX_CLEAR_FAILS = 3        # 同一频道清除失败 ≥ N 次 → SKIP

# 每个频道的内部状态
CH_NEW = 'new'
CH_M1 = 'm1'
CH_READY = 'ready'
CH_CLEARED = 'cleared'
CH_SKIP = 'skip'


# ============================================================
# Action
# ============================================================
@dataclass
class Action:
    kind: str            # 'measure' | 'clear' | 'done'
    pos: Vec = (0.0, 0.0)
    ch: int = 1

    def __repr__(self):
        if self.kind == 'done':
            return "Action(done)"
        return f"Action({self.kind}, pos={self.pos}, ch={self.ch})"


@dataclass
class State:
    pos: Vec
    ch: int
    cleared: Set[int] = field(default_factory=set)
    skipped: Set[int] = field(default_factory=set)
    bearing_records: Dict[int, List[Tuple[Vec, float]]] = field(default_factory=dict)
    virtual_time: float = 0.0


# ============================================================
# 工具：定位 + 清除
# ============================================================
def localize(records: List[Tuple[Vec, float]],
             eps: float = EPS_DEG) -> Optional[Tuple[Vec, float]]:
    """对一组 (S, θ) 用 Q1 算法定位，返回 MEC 圆心和半径。"""
    if len(records) < 2:
        return None
    dets = [r[0] for r in records]
    thetas = [r[1] for r in records]
    res = p1.solve_problem_1(dets, thetas, eps_deg=eps,
                             R_eff=1500.0, R_target=ARENA_R, N_disk=32)
    if res['n_vertices'] < 3 or res['D'] < 1e-3:
        return None
    return res['mec_center'], res['mec_radius']


def second_pos_from_bearing(S: Vec, theta: float, dist: float = HOP_DIST) -> Vec:
    """从 S 沿 θ 方向走 dist 米。"""
    rad = math.radians(theta)
    return (S[0] + dist * math.cos(rad),
            S[1] + dist * math.sin(rad))


# ============================================================
# 基类
# ============================================================
class BaseStrategy:
    name: str = "base"

    def __init__(self):
        self.ch_state: Dict[int, str] = {ch: CH_NEW for ch in CHANNELS}
        self.bearings: Dict[int, List[Tuple[Vec, float]]] = {ch: [] for ch in CHANNELS}
        self.clear_fails: Dict[int, int] = {ch: 0 for ch in CHANNELS}
        self.cursor = 0
        self.done = False
        self._attempts = 0  # 总步数上限

    # 子类重写：返回某个频道的第 1 / 第 2 个 measure 位置
    def first_pos(self, ch: int, state: State) -> Vec:
        return STATIC_POS

    def second_pos(self, ch: int, state: State,
                   first_record: Tuple[Vec, float]) -> Vec:
        S, theta = first_record
        return second_pos_from_bearing(S, theta, HOP_DIST)

    def next_action_for_ch(self, ch: int, state: State) -> Optional[Action]:
        """对某个频道返回下一个动作；None 表示切到下个频道。"""
        st = self.ch_state[ch]
        if st == CH_CLEARED or st == CH_SKIP:
            return None
        recs = self.bearings[ch]
        if st == CH_NEW:
            return Action('measure', self.first_pos(ch, state), ch)
        if st == CH_M1:
            assert len(recs) == 1
            return Action('measure',
                          self.second_pos(ch, state, recs[0]), ch)
        if st == CH_READY:
            # 已经有 ≥ 2 个 bearing；尝试定位清除
            loc = localize(recs)
            if loc is None:
                # 几何退化（两扇形反向等），跳过
                self.ch_state[ch] = CH_SKIP
                state.skipped.add(ch)
                return None
            c, r = loc
            if r <= CLEAR_R:
                return Action('clear', c, ch)
            # 定位还不够准，再补一个 bearing
            # 简单策略：在上次某个 bearing 附近 + 沿另一方向 600m
            extra = self._extra_pos(ch, recs)
            return Action('measure', extra, ch)
        return None

    def _extra_pos(self, ch: int, recs: List[Tuple[Vec, float]]) -> Vec:
        """第 3+ 个 bearing 的位置：取已有 bearing 中心 + 沿某方向 600m。"""
        cx = sum(r[0][0] for r in recs) / len(recs)
        cy = sum(r[0][1] for r in recs) / len(recs)
        # 沿"已收 bearing 角度差最大"的方向
        thetas = [r[1] for r in recs]
        a = math.radians(sum(thetas) / len(thetas) + 90)  # 垂直方向
        return (cx + 600 * math.cos(a), cy + 600 * math.sin(a))

    def step(self, state: State) -> Action:
        if self.done:
            return Action('done')
        self._attempts += 1
        if self._attempts > 8000:
            self.done = True
            return Action('done')

        # 扫一圈找下一个需要处理的频道
        for _ in range(len(CHANNELS)):
            ch = CHANNELS[self.cursor % len(CHANNELS)]
            self.cursor += 1
            if ch in state.cleared or ch in state.skipped:
                continue
            a = self.next_action_for_ch(ch, state)
            if a is not None:
                return a
        # 全部清完或无新动作
        self.done = True
        return Action('done')

    # --------- 事件回调 ---------
    def on_measure(self, state: State, ch: int, result: str,
                   svd: Optional[float]):
        if result == 'direction' and svd is not None:
            recs = self.bearings[ch]
            # 避免重复（同一 pos 的 bearing 只记一次）
            pos = state.pos
            if not any((abs(S[0] - pos[0]) < 1e-3 and
                        abs(S[1] - pos[1]) < 1e-3) for S, _ in recs):
                recs.append((pos, svd))
            # 更新状态
            if self.ch_state[ch] == CH_NEW:
                self.ch_state[ch] = CH_M1
            elif self.ch_state[ch] in (CH_M1, CH_READY):
                # M1 → READY（如果 ≥ 2 个）；READY 再多收也无所谓
                if len(recs) >= 2:
                    self.ch_state[ch] = CH_READY
        elif result == 'no_signal':
            self.ch_state[ch] = CH_SKIP
            state.skipped.add(ch)
        elif result == 'near':
            # 已经在源附近 5m：直接进入"尝试清除"状态（用当前位置）
            self.ch_state[ch] = CH_READY
            # 不需要再测，把 state.pos 作为一个"假 bearing"也加入
            self.bearings[ch].append((state.pos, 0.0))
            self.bearings[ch].append((state.pos, 90.0))

    def on_clear(self, state: State, ch: int, success: bool):
        if success:
            self.ch_state[ch] = CH_CLEARED
            state.cleared.add(ch)
        else:
            self.clear_fails[ch] += 1
            if self.clear_fails[ch] >= MAX_CLEAR_FAILS:
                self.ch_state[ch] = CH_SKIP
                state.skipped.add(ch)


# ============================================================
# 具体策略
# ============================================================
class StaticScan(BaseStrategy):
    """固定在 (1000, 0)，对每个频道：
       1) measure(ch) @ (1000, 0)
       2) 若 direction：沿 θ 走 800m，再 measure(ch)
       3) 定位 → clear
    """
    name = "StaticScan"

    def first_pos(self, ch: int, state: State) -> Vec:
        return STATIC_POS


class Patrol(BaseStrategy):
    """在 4 个站位之间轮换：每站位扫所有频道。"""
    name = "Patrol"

    def __init__(self):
        super().__init__()
        self.site_idx = 0

    def first_pos(self, ch: int, state: State) -> Vec:
        return PATROL_POS[self.site_idx]

    def step(self, state: State) -> Action:
        # 在每站位扫一轮（len(CHANNELS) 个动作）
        a = super().step(state)
        # 当 cursor 回到 0，说明本轮结束 → 换站位
        if self.cursor % len(CHANNELS) == 0:
            self.site_idx = (self.site_idx + 1) % len(PATROL_POS)
        return a


class BlinkHop(BaseStrategy):
    """在初始 (0,0) 起步，每收 1 个 direction 就沿 θ 跳 800m 当作下个测点。"""
    name = "BlinkHop"

    def first_pos(self, ch: int, state: State) -> Vec:
        return state.pos  # 沿用当前位置

    def second_pos(self, ch: int, state: State,
                   first_record: Tuple[Vec, float]) -> Vec:
        S, theta = first_record
        return second_pos_from_bearing(S, theta, HOP_DIST)


class Hybrid(BaseStrategy):
    """早期：闪烁跳跃（pos 跟着 direction 走）；后期：固定到 STATIC_POS 精定位。"""
    name = "Hybrid"

    def __init__(self):
        super().__init__()
        self._phase = 'hop'   # 'hop' | 'static'
        self._switch_at = 5   # 已清除 ≥ 5 切换到 static

    def first_pos(self, ch: int, state: State) -> Vec:
        if self._phase == 'hop':
            return state.pos
        return STATIC_POS

    def second_pos(self, ch: int, state: State,
                   first_record: Tuple[Vec, float]) -> Vec:
        if self._phase == 'hop':
            S, theta = first_record
            return second_pos_from_bearing(S, theta, HOP_DIST)
        # static：固定基线，沿 θ 反方向 800m 作第 2 点
        S, theta = first_record
        rad = math.radians(theta)
        # 取与 first_pos 反向 + 800m
        return (STATIC_POS[0] - 800 * math.cos(rad),
                STATIC_POS[1] - 800 * math.sin(rad))

    def step(self, state: State) -> Action:
        # 阶段切换
        if self._phase == 'hop' and len(state.cleared) >= self._switch_at:
            self._phase = 'static'
        return super().step(state)


# ============================================================
# 工厂
# ============================================================
def make_strategy(name: str) -> BaseStrategy:
    return {
        'BlinkHop': BlinkHop,
        'StaticScan': StaticScan,
        'Patrol': Patrol,
        'Hybrid': Hybrid,
    }[name]()


if __name__ == "__main__":
    s = make_strategy('Hybrid')
    st = State(pos=(0.0, 0.0), ch=1)
    for i in range(8):
        a = s.step(st)
        print(i, a)
        if a.kind == 'done':
            break
\end{lstlisting}
%</paper-block>

%<paper-block id="appendix.p018" type="paragraph">
\noindent\textbf{文件：}\texttt{02\_\allowbreak{}最终代码/\allowbreak{}第四问/\allowbreak{}problem1\_\allowbreak{}v4\_\allowbreak{}inline.py}\\
\textbf{功能：}问题四内联交会定位几何模块
%</paper-block>
%<paper-block id="appendix.code016" type="code">
\begin{lstlisting}[style=appendixcode]
# -*- coding: utf-8 -*-
# problem1_v4_inline.py  （v5 自包含的 Q1 v4 算法副本）
#
# 这是 Q1 v4（problem1_v4.py）的完整内联版，src 在 解题库/第一问/v4/。
# v5 不依赖外部 Q1 v4 源码；本文件已经包含：
#   - HalfPlane、make_sector_halfplanes（楔形张角）
#   - hpi_intersect、halfplane_intersect、clip_polygon_by_disk（Sutherland-Hodgman）
#   - polygon_diameter、mec_welzl、diameter_circle_covers（直径 + Welzl MEC）
#   - localization_polygon、solve_problem_1（定位 + 直径圆覆盖判据）

import math
import random
from typing import List, Tuple, Optional, Sequence

Vec = Tuple[float, float]
Poly = List[Vec]


class HalfPlane:
    """半平面 n·P <= d（外向法向 n，闭合侧为内侧）。"""
    __slots__ = ('n', 'd')

    def __init__(self, n: Vec, d: float):
        self.n = n
        self.d = d

    def contains(self, P: Vec) -> bool:
        return self.n[0] * P[0] + self.n[1] * P[1] <= self.d + 1e-9


def make_sector_halfplanes(S: Vec, theta_deg: float, eps_deg: float
                           ) -> Tuple[HalfPlane, HalfPlane]:
    """扇形构造：点 P 与 S 的连线方向 ∈ [θ−ε, θ+ε]。"""
    lo = math.radians(theta_deg - eps_deg)
    hi = math.radians(theta_deg + eps_deg)
    u_lo = (math.cos(lo), math.sin(lo))
    u_hi = (math.cos(hi), math.sin(hi))
    n_lo = (u_lo[1], -u_lo[0])
    n_hi = (-u_hi[1], u_hi[0])
    d_lo = n_lo[0] * S[0] + n_lo[1] * S[1]
    d_hi = n_hi[0] * S[0] + n_hi[1] * S[1]
    return HalfPlane(n_lo, d_lo), HalfPlane(n_hi, d_hi)


def _seg_line_intersect(p1: Vec, p2: Vec, n: Vec, d: float) -> Vec:
    dx = p2[0] - p1[0]
    dy = p2[1] - p1[1]
    denom = n[0] * dx + n[1] * dy
    if abs(denom) < 1e-12:
        return p1
    t = (d - n[0] * p1[0] - n[1] * p1[1]) / denom
    return (p1[0] + t * dx, p1[1] + t * dy)


def hpi_intersect(poly: Poly, halfplanes: Sequence[HalfPlane]
                  ) -> Optional[Poly]:
    """Sutherland-Hodgman：给定初始凸多边形 poly，依次用各半平面裁剪。"""
    for hp in halfplanes:
        if poly is None or len(poly) == 0:
            return None
        new_poly: List[Vec] = []
        n, d = hp.n, hp.d
        m = len(poly)
        for i in range(m):
            cur = poly[i]
            prv = poly[i - 1]
            cur_in = (n[0] * cur[0] + n[1] * cur[1]) <= d + 1e-9
            prv_in = (n[0] * prv[0] + n[1] * prv[1]) <= d + 1e-9
            if cur_in:
                if not prv_in:
                    new_poly.append(_seg_line_intersect(prv, cur, n, d))
                new_poly.append(cur)
            elif prv_in:
                new_poly.append(_seg_line_intersect(prv, cur, n, d))
        poly = new_poly
    if poly is None or len(poly) < 3:
        return None
    return poly


def _bbox_polygon(bbox: float) -> Poly:
    b = float(bbox)
    return [(-b, -b), (b, -b), (b, b), (-b, b)]


def halfplane_intersect(halfplanes: Sequence[HalfPlane], bbox: float = 5000.0
                        ) -> Optional[Poly]:
    return hpi_intersect(_bbox_polygon(bbox), halfplanes)


def clip_polygon_by_disk(poly: Poly, c: Vec, r: float, N: int = 64
                         ) -> Optional[Poly]:
    """用 N 条线段近似圆盘 |P−c| ≤ r，截多边形 poly。"""
    if poly is None or len(poly) < 3:
        return None
    hps = []
    for k in range(N):
        a = 2 * math.pi * k / N
        n = (math.cos(a), math.sin(a))
        d = r + n[0] * c[0] + n[1] * c[1]
        hps.append(HalfPlane(n, d))
    return hpi_intersect(poly, hps)


def polygon_diameter(poly: Poly) -> Tuple[float, Vec, Vec]:
    """返回 (D, A, B)。O(n²) 暴力。"""
    n = len(poly)
    if n < 2:
        p = poly[0] if poly else (0.0, 0.0)
        return 0.0, p, p
    best_d2 = -1.0
    best_i, best_j = 0, 1
    for i in range(n):
        xi, yi = poly[i]
        for j in range(i + 1, n):
            dx = xi - poly[j][0]
            dy = yi - poly[j][1]
            d2 = dx * dx + dy * dy
            if d2 > best_d2:
                best_d2, best_i, best_j = d2, i, j
    return math.sqrt(best_d2), poly[best_i], poly[best_j]


def _circle_from_2(a: Vec, b: Vec) -> Tuple[Vec, float]:
    c = ((a[0] + b[0]) * 0.5, (a[1] + b[1]) * 0.5)
    r2 = (a[0] - c[0]) ** 2 + (a[1] - c[1]) ** 2
    return c, r2


def _circle_from_3(a: Vec, b: Vec, c: Vec) -> Optional[Tuple[Vec, float]]:
    ax, ay = a
    bx, by = b
    cx, cy = c
    d = 2 * (ax * (by - cy) + bx * (cy - ay) + cx * (ay - by))
    if abs(d) < 1e-12:
        return None
    ux = ((ax * ax + ay * ay) * (by - cy) +
          (bx * bx + by * by) * (cy - ay) +
          (cx * cx + cy * cy) * (ay - by)) / d
    uy = ((ax * ax + ay * ay) * (cx - bx) +
          (bx * bx + by * by) * (ax - cx) +
          (cx * cx + cy * cy) * (bx - ax)) / d
    ctr = (ux, uy)
    r2 = (ax - ux) ** 2 + (ay - uy) ** 2
    return ctr, r2


def mec_welzl(points: Sequence[Vec]) -> Tuple[Vec, float]:
    """Welzl 最小包围圆（迭代版，O(n) 期望时间）。"""
    import sys
    sys.setrecursionlimit(10000)
    P = list(points)
    random.shuffle(P)

    def circle_trivial(R: Sequence[Vec]) -> Tuple[Vec, float]:
        if len(R) == 0:
            return (0.0, 0.0), 0.0
        if len(R) == 1:
            return R[0], 0.0
        if len(R) == 2:
            return _circle_from_2(R[0], R[1])
        c, r2 = _circle_from_3(R[0], R[1], R[2])
        if c is None:
            return _circle_from_2(R[0], R[1])[0], max(
                (R[0][0] - R[1][0]) ** 2 + (R[0][1] - R[1][1]) ** 2,
                (R[0][0] - R[2][0]) ** 2 + (R[0][1] - R[2][1]) ** 2,
                (R[1][0] - R[2][0]) ** 2 + (R[1][1] - R[2][1]) ** 2,
            ) * 0.25
        return c, r2

    def mec(P: List[Vec], R: List[Vec], n: int) -> Tuple[Vec, float]:
        if n == 0 or len(R) == 3:
            return circle_trivial(R)
        p = P[n - 1]
        c, r2 = mec(P, R, n - 1)
        if (p[0] - c[0]) ** 2 + (p[1] - c[1]) ** 2 <= r2 + 1e-9:
            return c, r2
        return mec(P, R + [p], n - 1)

    c, r2 = mec(P, [], len(P))
    return c, math.sqrt(r2)


def diameter_circle_covers(poly: Poly, A: Vec, B: Vec
                           ) -> Tuple[bool, float, Tuple[Vec, float]]:
    """判断以 AB 为直径的圆是否能覆盖多边形。"""
    mid = ((A[0] + B[0]) * 0.5, (A[1] + B[1]) * 0.5)
    D = math.hypot(A[0] - B[0], A[1] - B[1])
    R = D / 2
    off = -1e9
    for p in poly:
        d = math.hypot(p[0] - mid[0], p[1] - mid[1])
        off = max(off, d - R)
    covered = off <= 1e-9
    c_mec, r_mec = mec_welzl(poly)
    return covered, off, (c_mec, r_mec)


def localization_polygon(dets: Sequence[Vec], thetas: Sequence[float],
                         eps_deg: float = 1.0,
                         R_eff: float = 1500.0,
                         R_target: float = 1800.0,
                         N_disk: int = 64) -> Optional[Poly]:
    """构造定位多边形 P = ∩扇形_i ∩ D(S_i, R_eff) ∩ D(O, R_target)。"""
    hps = []
    for S, th in zip(dets, thetas):
        hp_lo, hp_hi = make_sector_halfplanes(S, th, eps_deg)
        hps.extend([hp_lo, hp_hi])

    poly = halfplane_intersect(hps, bbox=5000.0)
    if poly is None:
        return None

    poly = clip_polygon_by_disk(poly, (0.0, 0.0), R_target, N_disk)
    if poly is None:
        return None

    for S in dets:
        poly = clip_polygon_by_disk(poly, S, R_eff, N_disk)
        if poly is None:
            return None

    return poly


def solve_problem_1(dets: Sequence[Vec], thetas: Sequence[float],
                    eps_deg: float = 1.0,
                    R_eff: float = 1500.0,
                    R_target: float = 1800.0,
                    N_disk: int = 64) -> dict:
    """完整求解第一问。返回含 mec_center, mec_radius, D 等的字典。"""
    poly = localization_polygon(dets, thetas, eps_deg, R_eff, R_target, N_disk)
    if poly is None:
        return {
            'D': 0.0, 'R_star': 0.0, 'ratio_2R_D': 0.0,
            'covered': False, 'max_offset': 0.0,
            'A': (0, 0), 'B': (0, 0), 'mec_center': (0, 0), 'mec_radius': 0.0,
            'n_vertices': 0, 'poly': [],
        }
    D, A, B = polygon_diameter(poly)
    covered, off, (c_mec, r_mec) = diameter_circle_covers(poly, A, B)
    return {
        'D': D,
        'R_star': r_mec,
        'ratio_2R_D': 2 * r_mec / D if D > 0 else 0.0,
        'covered': covered,
        'max_offset': off,
        'A': A,
        'B': B,
        'mec_center': c_mec,
        'mec_radius': r_mec,
        'n_vertices': len(poly),
        'poly': poly,
    }
\end{lstlisting}
%</paper-block>

%<paper-block id="appendix.p019" type="paragraph">
\noindent\textbf{文件：}\texttt{02\_\allowbreak{}最终代码/\allowbreak{}第四问/\allowbreak{}strategy\_\allowbreak{}p4.py}\\
\textbf{功能：}问题四定向源观测更新与基础搜索逻辑
%</paper-block>
%<paper-block id="appendix.code017" type="code">
\begin{lstlisting}[style=appendixcode]
# -*- coding: utf-8 -*-
"""问题4 v6：对定向源安全的发现、定位与清除策略。

策略不读取模拟器真值。发现阶段使用 700m 基线网格和 350m 错位外围网格，
在每一个站位检测全部频道。只有整张发现网格完成后，仍然从未出现
``direction``/``near`` 的频道才会被标记为 absent。

这里的 ``no_signal`` 永远不是定位负约束：

* 发现阶段不把单次 no_signal 从覆盖区域中排除；
* 已经有正向测量的频道不修改既有定位多边形；
* 局部测量丢失时，只回到最近的正向测量点恢复。

一旦有正向测量，交会定位、MEC 和 20m 清除判据复用
``problem1_v4_inline.py``。额外的局部探测点只取已验证正向测量点的凸组合，
因此不会把一个猜测的 no_signal 当成几何事实。
"""

import math
from dataclasses import dataclass, field
from typing import Dict, List, Optional, Set, Tuple

from strategy import Action, State, p1


Vec = Tuple[float, float]
CHANNELS = tuple(range(1, 21))
ARENA_R = 1800.0
R_EFF = 1500.0
EPS_DEG = 1.0050001
CLEAR_R = 20.0
GRID_STEP = 700.0
GRID_OFFSET = 350.0
MIN_DISCOVERY_RECORDS = 5


def _snake_grid(axis: Tuple[int, ...]) -> List[Vec]:
    """生成固定顺序的 700m 方格路线，不依赖源位置。"""
    points: List[Vec] = []
    for row, y in enumerate(axis):
        xs = axis if row % 2 == 0 else tuple(reversed(axis))
        points.extend((float(x), float(y)) for x in xs)
    return points


def _discovery_stations() -> List[Vec]:
    """700m 基线网格 + 错位外围网格。

    基线覆盖 [-2100, 2100]^2；外围错位网格覆盖
    [-2450, 2450]^2。两者都保留在目标圆外的必要站位，保证源在目标圆
    边缘且定向半平面朝外时仍有近距离正向检测机会。
    """
    base_axis = tuple(range(-2100, 2101, 700))
    offset_axis = tuple(range(-2450, 2451, 700))
    ordered = _snake_grid(base_axis) + _snake_grid(offset_axis)

    # 站点集合保持不变；完成全清验证后，只优化访问顺序以减少移动时间。
    origin = (0.0, 0.0)
    seen = {origin}
    points = []
    for point in ordered:
        if point not in seen:
            seen.add(point)
            points.append(point)

    # 确定性最近邻开放路径；它不使用任何源信息，也不改变发现覆盖集合。
    result = [origin]
    current = origin
    while points:
        next_point = min(points, key=lambda point: (
            math.hypot(current[0] - point[0], current[1] - point[1]),
            point[1], point[0]))
        result.append(next_point)
        points.remove(next_point)
        current = next_point
    return result


DISCOVERY_STATIONS = tuple(_discovery_stations())


def distance(a: Vec, b: Vec) -> float:
    return math.hypot(a[0] - b[0], a[1] - b[1])


def _same_point(a: Vec, b: Vec, tol: float = 1e-3) -> bool:
    return distance(a, b) <= tol


def _localize(records: List[Tuple[Vec, float]]) -> Optional[dict]:
    """只用 direction 记录交会定位并计算 MEC。"""
    if len(records) < 2:
        return None
    dets = [pos for pos, _ in records]
    thetas = [angle for _, angle in records]
    result = p1.solve_problem_1(
        dets,
        thetas,
        eps_deg=EPS_DEG,
        R_eff=R_EFF,
        R_target=ARENA_R,
        N_disk=64,
    )
    if (result['n_vertices'] < 3 or result['D'] <= 1e-6
            or not all(math.isfinite(v) for v in result['mec_center'])
            or not math.isfinite(result['mec_radius'])):
        return None
    return result


@dataclass
class Track:
    records: List[Tuple[Vec, float]] = field(default_factory=list)
    near: Optional[Vec] = None
    poly: List[Vec] = field(default_factory=list)
    center: Vec = (0.0, 0.0)
    radius: float = math.inf
    estimate: Vec = (0.0, 0.0)
    recovery_required: bool = False
    clear_failures: int = 0
    probe_round: int = 0
    probe_history: List[Vec] = field(default_factory=list)


class AdaptiveP4:
    """问题4独立策略：完整安全发现后，再逐频道定位和清除。"""

    name = 'AdaptiveP4'

    def __init__(self, discovery_stations: Optional[List[Vec]] = None,
                 min_positive_records: int = MIN_DISCOVERY_RECORDS):
        self.discovery_stations = list(discovery_stations or DISCOVERY_STATIONS)
        self.min_positive_records = max(2, int(min_positive_records))
        self.discovery_station_index = 0
        self.discovery_channel_index = 0
        self._station_channels: List[int] = []
        self.phase = 'discovery'
        self.discovery_complete = False

        self.tracks: Dict[int, Track] = {}
        self.absent: Set[int] = set()
        self.cleared: Set[int] = set()
        self.work_channels: List[int] = []
        self.work_index = 0

        self.done = False
        self.completion_reason: Optional[str] = None
        self._attempts = 0
        self.geometry_conflicts = 0

    # ------------------------------------------------------------
    # 发现阶段
    # ------------------------------------------------------------
    def _finish_discovery(self) -> None:
        if self.discovery_complete:
            return
        self.discovery_complete = True
        # 只有完整网格完成后才允许做这一次集合判定；no_signal 本身不参与。
        self.absent = set(CHANNELS) - set(self.tracks)
        self.work_channels = sorted(self.tracks)
        self.work_index = 0
        self.phase = 'localization'

    def _next_discovery_action(self) -> Optional[Action]:
        if (self._station_channels
                and self.discovery_channel_index >= len(self._station_channels)):
            self.discovery_station_index += 1
            self.discovery_channel_index = 0
            self._station_channels = []

        while self.discovery_station_index < len(self.discovery_stations):
            if self.discovery_channel_index >= len(self._station_channels):
                station = self.discovery_stations[self.discovery_station_index]
                self._station_channels = [
                    channel for channel in CHANNELS
                    if channel not in self.tracks
                    or (self.tracks[channel].near is None
                        and len(self.tracks[channel].records)
                        < self.min_positive_records)
                ]
                self.discovery_channel_index = 0
                if not self._station_channels:
                    self.discovery_station_index += 1
                    continue

            station = self.discovery_stations[self.discovery_station_index]
            channel = self._station_channels[self.discovery_channel_index]
            self.discovery_channel_index += 1
            return Action('measure', station, channel)

        self._finish_discovery()
        return None

    # ------------------------------------------------------------
    # 正向记录与安全探测
    # ------------------------------------------------------------
    @staticmethod
    def _nearest_record(track: Track, pos: Vec) -> Vec:
        if not track.records:
            return pos
        return min(track.records, key=lambda row: distance(row[0], pos))[0]

    def _safe_probe_candidates(self, track: Track, state: State) -> List[Vec]:
        """从正向测量点凸组合产生候选点。

        对同一源，正向可接收集合是有效接收圆盘与定向半平面的交集，
        是凸集；其内的凸组合仍然是正向候选。候选永远不由 no_signal 生成。
        """
        positions: List[Vec] = []
        for pos, _ in track.records:
            if not any(_same_point(pos, old) for old in positions):
                positions.append(pos)
        if len(positions) < 2:
            return positions[:]

        pairs = []
        for i in range(len(positions)):
            for j in range(i + 1, len(positions)):
                pair_distance = distance(positions[i], positions[j])
                if pair_distance > 1e-3:
                    pairs.append((pair_distance, positions[i], positions[j]))
        pairs.sort(reverse=True, key=lambda item: item[0])

        candidates: List[Vec] = []
        # 优先长基线，再在基线内部取多个新的安全点。
        for _, a, b in pairs[:80]:
            for fraction in (0.12, 0.25, 0.40, 0.60, 0.75, 0.88):
                point = (
                    a[0] * (1.0 - fraction) + b[0] * fraction,
                    a[1] * (1.0 - fraction) + b[1] * fraction,
                )
                if any(_same_point(point, old, 0.5)
                       for old in positions + track.probe_history):
                    continue
                candidates.append(point)

        # 加入三点平均值，避免所有候选都落在同一条长边上。
        for group in (positions[:3], positions[-3:], positions[::2][:3]):
            if len(group) == 3:
                point = tuple(sum(p[i] for p in group) / 3.0 for i in (0, 1))
                if not any(_same_point(point, old, 0.5)
                           for old in positions + track.probe_history):
                    candidates.append(point)

        if not candidates:
            return []

        center = track.center if math.isfinite(track.radius) else state.pos

        def angular_separation(point: Vec) -> float:
            if distance(point, center) < 1e-6:
                return math.pi
            new_angle = math.atan2(center[1] - point[1], center[0] - point[0])
            separations = []
            for old, _ in track.records:
                if distance(old, center) < 1e-6:
                    continue
                old_angle = math.atan2(center[1] - old[1], center[0] - old[0])
                separations.append(abs(math.sin(new_angle - old_angle)))
            return min(separations) if separations else 0.0

        def score(point: Vec) -> Tuple[float, float, float]:
            # 先取新方向最不同的点；再偏好靠近当前估计、且不重复的位置。
            return (
                angular_separation(point),
                -distance(point, center),
                -distance(point, state.pos),
            )

        return sorted(candidates, key=score, reverse=True)

    def _next_probe(self, track: Track, state: State) -> Action:
        candidates = self._safe_probe_candidates(track, state)
        if candidates:
            point = candidates[0]
            track.probe_history.append(point)
            track.probe_round += 1
            return Action('measure', point, self._active_channel)

        # 仅在没有第二个不同正向点时重复已验证正向位置；这仍然不会产生
        # 空间负约束，新的独立噪声测量有机会收紧同一测点的角度区间。
        point = self._nearest_record(track, state.pos)
        track.probe_round += 1
        return Action('measure', point, self._active_channel)

    # ------------------------------------------------------------
    # 局部定位与清除阶段
    # ------------------------------------------------------------
    def _next_local_action(self, channel: int, state: State) -> Action:
        track = self.tracks[channel]
        self._active_channel = channel

        if track.near is not None:
            return Action('clear', track.near, channel)

        if track.recovery_required:
            # 消费恢复请求；若该点再次 no_signal，on_measure 会重新置位。
            track.recovery_required = False
            return Action('measure', self._nearest_record(track, state.pos), channel)

        localized = _localize(track.records)
        if localized is not None:
            track.poly = localized['poly']
            track.center = localized['mec_center']
            track.radius = localized['mec_radius']
            track.estimate = localized['mec_center']

            if track.radius <= CLEAR_R:
                return Action('clear', track.center, channel)

        return self._next_probe(track, state)

    # ------------------------------------------------------------
    # 驱动器接口
    # ------------------------------------------------------------
    def step(self, state: State) -> Action:
        if self.done:
            return Action('done')
        self._attempts += 1
        if self._attempts > 30000:
            self.done = True
            self.completion_reason = 'step_limit'
            return Action('done')

        if self.phase == 'discovery':
            action = self._next_discovery_action()
            if action is not None:
                return action

        while self.work_index < len(self.work_channels):
            channel = self.work_channels[self.work_index]
            if channel in self.cleared:
                self.work_index += 1
                continue
            return self._next_local_action(channel, state)

        self.done = True
        self.completion_reason = 'all_channels_cleared_or_covered'
        return Action('done')

    def on_measure(self, state: State, channel: int, result: str,
                   svd: Optional[float]) -> None:
        if result in ('direction', 'near'):
            track = self.tracks.setdefault(channel, Track())
            if result == 'near':
                track.near = state.pos
                track.recovery_required = False
                return
            if svd is None or not math.isfinite(svd):
                return
            # discovery 阶段每个站位只测一次；局部恢复允许同一点重复测量，
            # 作为新的独立噪声样本，但从不把 no_signal 作为约束。
            track.records.append((state.pos, float(svd) % 360.0))
            track.recovery_required = False
            return

        if result == 'no_signal' and self.phase == 'localization':
            track = self.tracks.get(channel)
            if track is not None and track.records:
                # 保留 poly、center、radius 和全部正向 records 原样。
                track.recovery_required = True

    def on_clear(self, state: State, channel: int, success: bool) -> None:
        track = self.tracks[channel]
        if success:
            self.cleared.add(channel)
            state.cleared.add(channel)
            track.recovery_required = False
            self.work_index += 1
            return

        # 清除失败也不能将频道标 absent 或 skip；回到正向点继续收集几何信息。
        track.clear_failures += 1
        track.recovery_required = True


def make_strategy() -> AdaptiveP4:
    return AdaptiveP4()


if __name__ == '__main__':
    strategy = AdaptiveP4()
    state = State(pos=(0.0, 0.0), ch=1)
    print(strategy.name, len(strategy.discovery_stations), strategy.step(state))
\end{lstlisting}
%</paper-block>

%<paper-block id="appendix.p020" type="paragraph">
\noindent\textbf{文件：}\texttt{02\_\allowbreak{}最终代码/\allowbreak{}第四问/\allowbreak{}strategy\_\allowbreak{}fast.py}\\
\textbf{功能：}问题四搜索与定位调度基类
%</paper-block>
%<paper-block id="appendix.code018" type="code">
\begin{lstlisting}[style=appendixcode]
"""演练候选：先完成有证明的49站发现网格，仅为已知难定位源补扫。"""
import math
from strategy import Action
from strategy_p4 import AdaptiveP4, CHANNELS, DISCOVERY_STATIONS, _localize, distance

def nearest_path(points,start=(0.0,0.0)):
    remaining=set(points);out=[]
    while remaining:
        nxt=min(remaining,key=lambda p:(distance(start,p),p[1],p[0]))
        out.append(nxt);remaining.remove(nxt);start=nxt
    return out

CORE_STATIONS=tuple(nearest_path([(float(x),float(y)) for x in range(-2100,2101,700)
                                for y in range(-2100,2101,700)]))

class FastP4(AdaptiveP4):
    name='FastP4_grid49'
    def __init__(self):
        extras=set(DISCOVERY_STATIONS)-set(CORE_STATIONS)
        super().__init__(list(CORE_STATIONS)+nearest_path(extras,CORE_STATIONS[-1]))
        self.core_complete=False
        self._geometry_counts={}
        self._ordered=False

    def _update_geometry(self,ch):
        track=self.tracks[ch]
        if len(track.records)<2 or self._geometry_counts.get(ch)==len(track.records):return
        self._geometry_counts[ch]=len(track.records)
        geom=_localize(track.records)
        if geom is not None:
            track.poly=geom['poly'];track.center=geom['mec_center']
            track.radius=geom['mec_radius'];track.estimate=geom['mec_center']

    def _needs_sample(self,ch):
        if ch in self.cleared or ch in self.absent:return False
        track=self.tracks.get(ch)
        if track is None:return not self.core_complete
        return (track.near is None and track.radius>19.99999
                and len(track.records)<self.min_positive_records)

    def _next_discovery_action(self):
        while self.discovery_station_index<len(self.discovery_stations):
            if self._station_channels and self.discovery_channel_index>=len(self._station_channels):
                self.discovery_station_index+=1;self.discovery_channel_index=0;self._station_channels=[]
                continue
            if self.discovery_station_index>=len(CORE_STATIONS) and not self.core_complete:
                # Every target lies in a 700m square whose four corners are <=700sqrt(2)<1000m.
                # Any closed half-plane through the target contains at least one corner.
                # Thus only AFTER these 49 stations can an unseen channel be declared absent.
                self.core_complete=True;self.absent=set(CHANNELS)-set(self.tracks)
            if not self._station_channels:
                self._station_channels=[ch for ch in CHANNELS if self._needs_sample(ch)]
                self.discovery_channel_index=0
            if not self._station_channels:
                self.discovery_station_index+=1;continue
            ch=self._station_channels[self.discovery_channel_index]
            self.discovery_channel_index+=1
            if not self._needs_sample(ch):continue
            return Action('measure',self.discovery_stations[self.discovery_station_index],ch)
        self._finish_discovery()
        return None

    def on_measure(self,state,channel,result,svd):
        # Repeated readings at one point must not be treated as independent spatial information.
        track=self.tracks.get(channel)
        if result=='direction' and track and any(distance(state.pos,p)<1e-6 for p,_ in track.records):
            track.recovery_required=False;return
        super().on_measure(state,channel,result,svd)
        if result=='direction':self._update_geometry(channel)

    def _target(self,ch):
        track=self.tracks[ch]
        if track.near is not None:return track.near
        if math.isfinite(track.radius):return track.center
        return track.records[-1][0]

    def _next_local_action(self,channel,state):
        # This hook also runs on the discovery -> localization transition.
        if not self._ordered:
            # Reorder only between completed channels, never interrupt a localization recovery.
            remaining=[ch for ch in self.work_channels if ch not in self.cleared]
            order=[];pos=state.pos
            while remaining:
                ch=min(remaining,key=lambda c:(distance(pos,self._target(c)),c))
                order.append(ch);remaining.remove(ch);pos=self._target(ch)
            self.work_channels=order;self.work_index=0;self._ordered=True
            channel=order[0]
        action=super()._next_local_action(channel,state)
        track=self.tracks[channel]
        if action.kind=='clear' and track.near is None and track.radius<19.99999:
            d=distance(state.pos,track.center)
            shift=min(d,19.99999-track.radius)
            if d>1e-9:
                point=tuple(track.center[i]+(state.pos[i]-track.center[i])*shift/d for i in (0,1))
                return Action('clear',point,channel)
        return action

    def on_clear(self,state,channel,success):
        super().on_clear(state,channel,success)
        if success:self._ordered=False
\end{lstlisting}
%</paper-block>

%<paper-block id="appendix.p021" type="paragraph">
\noindent\textbf{文件：}\texttt{02\_\allowbreak{}最终代码/\allowbreak{}第四问/\allowbreak{}strategy\_\allowbreak{}cost.py}\\
\textbf{功能：}问题四有限光学覆盖函数及成本计算
%</paper-block>
%<paper-block id="appendix.code019" type="code">
\begin{lstlisting}[style=appendixcode]
"""One bounded P4 candidate; all switches off delegates exactly to frozen FastP4.

Geometry uses only positive bearings, the 1500m upper radius, arena, and
failed optical clears. No direction-dependent no_signal spatial exclusion.
"""
import itertools
import math
from collections import Counter
from shapely.geometry import Point, Polygon
from shapely.ops import unary_union
from strategy import Action
from strategy_fast import FastP4
from strategy_p4 import CHANNELS, distance

OPTICAL_R = 19.99999


def optical_cover(region, start, limit=6):
    """Cover a continuous superset by <=6 inscribed optical disks.

    Tile an oriented bounding rectangle into rectangles of diagonal <40m.
    Exact polygon difference against inscribed disks certifies full coverage;
    no target samples, area tolerance, or missing slivers are accepted.
    """
    if region is None or region.is_empty or not region.is_valid:
        return None
    rect = list(region.minimum_rotated_rectangle.exterior.coords)
    axes = [(1., 0.)]
    if len(rect) >= 2:
        dx, dy = rect[1][0]-rect[0][0], rect[1][1]-rect[0][1]
        d = math.hypot(dx, dy)
        if d > 1e-9:
            axes.append((dx/d, dy/d))
    coords = list(region.convex_hull.exterior.coords)
    best = None
    for ux, uy in axes:
        vx, vy = -uy, ux
        xs = [x*ux+y*uy for x,y in coords]
        ys = [x*vx+y*vy for x,y in coords]
        x0,x1,y0,y1 = min(xs),max(xs),min(ys),max(ys)
        for nx in range(1,limit+1):
            for ny in range(1,limit//nx+1):
                if math.hypot((x1-x0)/nx,(y1-y0)/ny)/2 > 19.99:
                    continue
                centers = []
                for i in range(nx):
                    for j in range(ny):
                        a=x0+(i+.5)*(x1-x0)/nx; b=y0+(j+.5)*(y1-y0)/ny
                        centers.append((a*ux+b*vx,a*uy+b*vy))
                disks = unary_union([Point(p).buffer(OPTICAL_R,quad_segs=32) for p in centers])
                if not region.difference(disks).is_empty:
                    continue
                # k-1 failures, final success; clear never switches measurement channel.
                path = min(itertools.permutations(centers),key=lambda ps: sum(
                    distance(a,b) for a,b in zip((start,)+ps,ps)))
                movement = sum(distance(a,b)/5 for a,b in zip((start,)+path,path))
                cost = movement+3*(len(path)-1)+5
                if best is None or cost < best['total_s']:
                    best=dict(points=list(path),total_s=cost,move_s=movement,
                              measure_s=0.,switch_s=0.,clear_success_s=5.,
                              clear_failure_s=3.*(len(path)-1))
    return best


class CostAwareP4(FastP4):
    def __init__(self, enable_a=True, enable_b=True, enable_c=False):
        super().__init__()
        self.enable_a,self.enable_b,self.enable_c=enable_a,enable_b,enable_c
        self.name='CostAwareP4_'+(''.join(k for k,v in zip('ABC',(enable_a,enable_b,enable_c)) if v) or 'off')
        self.ever_observed_channels=set()
        self.failed_clear_positions={}
        self.optical_plans={}
        self.events=Counter()
        self.cost_comparisons=[]
        self._last_action=None
        self._same_action_count=0
        self._local_counts=Counter()
        self.action_stage='discovery'

    def diagnostics(self):
        return dict(switches=dict(A=self.enable_a,B=self.enable_b,C=self.enable_c),
                    triggers=dict(self.events),ever_observed_channels=sorted(self.ever_observed_channels),
                    failed_clear_positions=self.failed_clear_positions,cost_comparisons=self.cost_comparisons)

    def _needs_sample(self,ch):
        if self.enable_a and len(self.ever_observed_channels)==16 and ch not in self.ever_observed_channels:
            return False
        return super()._needs_sample(ch)

    def on_measure(self,state,channel,result,svd):
        if result in ('direction','near'):
            self.ever_observed_channels.add(channel)
            if self.enable_a and len(self.ever_observed_channels)==16 and not self.events['A_cardinality_proof']:
                self.events['A_cardinality_proof']=1
                self.absent.update(set(CHANNELS)-self.ever_observed_channels)
        super().on_measure(state,channel,result,svd)

    def region(self,ch):
        """Outer polygon: tangent disk halfplanes + bearing error incl rounding.

        P4 _localize uses eps=1.0050001 and circumscribed (not inscribed)
        1500m/1800m disks. Removing INSCRIBED failed-clear disks remains outer.
        This is recomputed with all stored exclusions after every localization.
        """
        t=self.tracks[ch]
        if len(t.poly)<3:
            return None
        g=Polygon(t.poly)
        if not g.is_valid:
            # Unusable numerical polygon is not an optical coverage certificate.
            # Fall back to the original action instead of repairing/shrinking it.
            self.events['invalid_geometry_fallback']+=1
            return None
        for p in self.failed_clear_positions.get(ch,[]):
            g=g.difference(Point(p).buffer(OPTICAL_R,quad_segs=32))
        if g.is_empty:
            raise RuntimeError('geometry_contradiction: empty feasible region after failed clears')
        return g

    def _instant(self,state):
        for ch,t in sorted(self.tracks.items()):
            if ch in self.cleared:
                continue
            near=t.near is not None and distance(state.pos,t.near)<1e-7
            g=None if near else self.region(ch)
            covered=g is not None and g.difference(Point(state.pos).buffer(OPTICAL_R,quad_segs=32)).is_empty
            if near or covered:
                self.events['B_near' if near else 'B_region']=self.events['B_near' if near else 'B_region']+1
                return Action('clear',state.pos,ch)
        return None

    def step(self,state):
        if not any((self.enable_a,self.enable_b,self.enable_c)):
            return super().step(state)
        action=None
        if self.enable_b and self.phase=='discovery':
            action=self._instant(state)
        if action is not None:
            self.action_stage='instant_clear'
        else:
            action=super().step(state)
            self.action_stage='optical' if action.ch in self.optical_plans else self.phase
        if action.kind!='done':
            signature=(action.kind,tuple(action.pos),action.ch)
            self._same_action_count=self._same_action_count+1 if signature==self._last_action else 1
            self._last_action=signature
            if self._same_action_count>2:
                self.events['no_progress_stop']+=1
                raise RuntimeError('no_progress: repeated identical action')
        return action

    def _next_local_action(self,channel,state):
        # FastP4 may reorder the active channel; respect its selected action.ch.
        if self.enable_c and self._ordered and channel in self.optical_plans:
            plan=self.optical_plans[channel]
            if not plan:
                self.events['geometry_contradiction']+=1
                raise RuntimeError('geometry_contradiction: optical cover exhausted without success')
            return Action('clear',plan[0],channel)
        action=super()._next_local_action(channel,state)
        channel=action.ch
        if not self.enable_c:
            return action
        self._local_counts[channel]+=1
        if self._local_counts[channel]>256:
            raise RuntimeError('no_progress: bounded localization budget exhausted')
        g=self.region(channel)
        if action.kind!='measure' or g is None:
            return action
        plan=optical_cover(g,state.pos)
        if plan is None:
            self.events['C_no_cover']+=1
            return action
        candidates=self._safe_probe_candidates(self.tracks[channel],state)+[action.pos]
        # Optimistic LOWER bound for safe-probe-and-clear: travel, detection,
        # channel switch, minimum subsequent travel, one successful clear.
        # Further measurements/failures cost >=0. Optical uses worst-case misses.
        # Only replace if even that lower bound is more expensive.
        def lower(p):
            return distance(state.pos,p)/5+5+float(state.ch!=channel)+max(0.,Point(p).distance(g)-20.)/5+5
        safe=min(lower(p) for p in candidates)
        self.cost_comparisons.append(dict(channel=channel,optical=plan,safe_lower_bound_s=safe,
                                          selected=plan['total_s']+1e-6<safe))
        if plan['total_s']+1e-6>=safe:
            self.events['C_no_advantage']+=1
            return action
        self.optical_plans[channel]=plan['points'][:]
        self.events['C_plan_started']+=1
        return Action('clear',self.optical_plans[channel][0],channel)

    def on_clear(self,state,channel,success):
        if not any((self.enable_a,self.enable_b,self.enable_c)):
            return super().on_clear(state,channel,success)
        before=self.work_index
        if not success:
            self.failed_clear_positions.setdefault(channel,[]).append(tuple(state.pos))
            self.events['failed_clear_exclusion']+=1
        super().on_clear(state,channel,success)
        # Parent advances local index even during discovery; undo only that increment.
        if self.phase=='discovery':
            self.work_index=before
        if channel in self.optical_plans:
            if success:
                del self.optical_plans[channel]
                self.events['C_plan_success']+=1
            else:
                self.optical_plans[channel].pop(0)
                self.tracks[channel].recovery_required=False
                self.events['C_plan_failure_advance']+=1
                if not self.optical_plans[channel]:
                    self.events['geometry_contradiction']+=1
                    raise RuntimeError('geometry_contradiction: optical cover exhausted without success')
        elif not success and self.action_stage=='instant_clear':
            raise RuntimeError('geometry_contradiction: certified instantaneous clear failed')
\end{lstlisting}
%</paper-block>

%<paper-block id="appendix.p022" type="paragraph">
\noindent\textbf{文件：}\texttt{02\_\allowbreak{}最终代码/\allowbreak{}第四问/\allowbreak{}strategy\_\allowbreak{}finish.py}\\
\textbf{功能：}问题四终局光学覆盖与清除控制
%</paper-block>
%<paper-block id="appendix.code020" type="code">
\begin{lstlisting}[style=appendixcode]
"""P4 finish-cost transfer: stop extra survey only with a finite optical certificate.

Inherits FastP4 directly. No CostAware policy chain and no P3 imports.
The complete 49-station discovery proof and all simulator physics stay unchanged.
"""
import math
from collections import Counter
from shapely.geometry import Polygon,Point
from strategy import Action
from strategy_fast import FastP4, CORE_STATIONS, nearest_path
from strategy_p4 import CHANNELS, DISCOVERY_STATIONS, distance
from strategy_cost import optical_cover


def optimized_core():
    points=[p for p in CORE_STATIONS if not (abs(p[0])==2100 and abs(p[1])==2100)]
    route=nearest_path(points)
    # Fixed-start OPEN route: deterministic bounded 2-opt; no case observations.
    for _ in range(80):
        best=None
        for i in range(1,len(route)-1):
            for j in range(i+1,len(route)):
                before=distance(route[i-1],route[i])
                after=distance(route[i-1],route[j])
                if j+1<len(route):
                    before+=distance(route[j],route[j+1]);after+=distance(route[i],route[j+1])
                gain=before-after
                if gain>1e-7 and (best is None or gain>best[0]):best=(gain,i,j)
        if best is None:break
        _,i,j=best;route[i:j+1]=reversed(route[i:j+1])
    return tuple(route)

TRIMMED_CORE=optimized_core()


class FinishP4(FastP4):
    name='FinishP4_certified_tail'

    def __init__(self, enabled=True, trim_corners=False):
        super().__init__()
        self.enabled=enabled
        self.trim_corners=trim_corners
        self.name='FinishP4_'+('tail' if enabled else '')+('_grid45' if trim_corners else '')
        if trim_corners:
            extras=set(DISCOVERY_STATIONS)-set(CORE_STATIONS)
            self.discovery_stations=list(TRIMMED_CORE)+nearest_path(extras,TRIMMED_CORE[-1])
        self.certificates={}
        self.plans={}
        self.events=Counter()
        self._certificate_checks={}
        self.failed_clears={}
        self.action_stage='discovery'

    def _next_discovery_action(self):
        if not self.trim_corners:return super()._next_discovery_action()
        while self.discovery_station_index<len(self.discovery_stations):
            if self._station_channels and self.discovery_channel_index>=len(self._station_channels):
                self.discovery_station_index+=1;self.discovery_channel_index=0;self._station_channels=[]
                continue
            if self.discovery_station_index>=len(TRIMMED_CORE) and not self.core_complete:
                # Every 700m cell intersecting the 1800m disk retains all corners.
                # The four discarded corners touch only cells >=sqrt(2)*1400 away.
                self.core_complete=True;self.absent=set(CHANNELS)-set(self.tracks)
            if not self._station_channels:
                self._station_channels=[ch for ch in CHANNELS if self._needs_sample(ch)]
                self.discovery_channel_index=0
            if not self._station_channels:
                self.discovery_station_index+=1;continue
            ch=self._station_channels[self.discovery_channel_index];self.discovery_channel_index+=1
            if not self._needs_sample(ch):continue
            return Action('measure',self.discovery_stations[self.discovery_station_index],ch)
        self._finish_discovery()
        return None

    def _certificate(self,ch):
        track=self.tracks[ch]
        key=(len(track.records),tuple(track.poly))
        if self._certificate_checks.get(ch)==key:
            return self.certificates.get(ch)
        self._certificate_checks[ch]=key
        self.certificates.pop(ch,None)
        if len(track.records)<2 or len(track.poly)<3:return None
        region=Polygon(track.poly)
        # Never repair invalid polygons into smaller regions; keep original policy.
        if not region.is_valid or region.is_empty:
            self.events['invalid_polygon_fallback']+=1
            return None
        # Starting at the existing target estimate isolates the terminal extra cost;
        # actual routing is replanned from the real position when this channel runs.
        plan=optical_cover(region,track.center,limit=6)
        if plan is None:return None
        # Bounded terminal surcharge (including up to 5 misses) must be cheaper
        # than one 700m grid move + measurement. This is a selection heuristic,
        # not a guarantee of globally improved route cost; paired gate is required.
        if plan['total_s']>700./5+5:
            self.events['finish_cost_rejected']+=1
            return None
        proof=dict(polygon=list(track.poly),points=plan['points'],
                   worst_finish_from_center_s=plan['total_s'],records=len(track.records))
        self.certificates[ch]=proof
        return proof

    def _needs_sample(self,ch):
        needed=super()._needs_sample(ch)
        if not self.enabled or not needed or not self.core_complete:return needed
        # If any channel still needs the outer tour, keep cheap opportunistic
        # readings of every eligible channel along that same tour.
        remaining=[c for c in self.tracks if FastP4._needs_sample(self,c)]
        proofs={c:self._certificate(c) for c in remaining}
        if any(proof is None for proof in proofs.values()):return needed
        for c in remaining:
            if c not in self.plans:
                self.events['extra_scan_channel_retired']+=1
                self.plans[c]=None
        return False

    def step(self,state):
        action=super().step(state)
        self.action_stage='optical_finish' if action.ch in self.plans and action.kind=='clear' else self.phase
        return action

    def _next_local_action(self,ch,state):
        if self.enabled and self._ordered and self.plans.get(ch):
            return Action('clear',self.plans[ch][0],ch)
        # FastP4 can reorder channels in this hook, so first let it select action.ch.
        action=super()._next_local_action(ch,state)
        ch=action.ch
        if not self.enabled or ch not in self.plans:return action
        if self.plans[ch] is None:
            region=Polygon(self.certificates[ch]['polygon'])
            plan=optical_cover(region,state.pos,limit=6)
            if plan is None:raise RuntimeError('certificate_plan_unavailable')
            self.plans[ch]=plan['points'][:]
            self.events['optical_plan_started']+=1
        if not self.plans[ch]:raise RuntimeError('optical_cover_exhausted: inconsistent observations')
        return Action('clear',self.plans[ch][0],ch)

    def on_clear(self,state,ch,success):
        if ch not in self.plans:return super().on_clear(state,ch,success)
        track=self.tracks[ch]
        if success:
            super().on_clear(state,ch,True)
            del self.plans[ch]
            self.events['optical_success']+=1
            return
        # Immutable full-region cover stays valid for its shrinking subset.
        # Store negative evidence; no relocalization occurs mid-plan.
        self.failed_clears.setdefault(ch,[]).append(tuple(state.pos))
        self.plans[ch].pop(0)
        track.clear_failures+=1;track.recovery_required=False
        self.events['optical_miss_advance']+=1
        if not self.plans[ch]:raise RuntimeError('optical_cover_exhausted: inconsistent observations')

    def diagnostics(self):
        return dict(enabled=self.enabled,trim_corners=self.trim_corners,triggers=dict(self.events),certificates=self.certificates,
                    failed_clear_positions=self.failed_clears)
\end{lstlisting}
%</paper-block>

%<paper-block id="appendix.p023" type="paragraph">
\noindent\textbf{文件：}\texttt{02\_\allowbreak{}最终代码/\allowbreak{}第四问/\allowbreak{}discovery\_\allowbreak{}mesh.py}\\
\textbf{功能：}问题四连续覆盖核验所需几何网格辅助
%</paper-block>
%<paper-block id="appendix.code021" type="code">
\begin{lstlisting}[style=appendixcode]
"""A continuous, directional-discovery certificate using 25 fixed stations.

This module is independent of simulator state and of every P3 module.  The
certificate uses the official lower reception radius (1000 m), not a source's
hidden radius.  It says nothing about no_signal localization exclusions.

Construction/proof
------------------
Tile the whole plane with equilateral triangles of side 999.999 m.  Keep all
triangles whose CLOSED area intersects the CLOSED 1800 m target disk, and keep
all their vertices.  Triangle/disk intersection is decided by exact geometric
formulas (with conservative floating-point padding), never a point sample.

For every possible source p, at least one retained triangle contains p.  Each
vertex is at distance <= the triangle's diameter < 1000 m from p.  For every
unit direction n, p=sum(w_i*v_i), w_i>=0, sum(w_i)=1 gives
sum(w_i * n.dot(v_i-p))=0, so at least one vertex is in the closed transmitting
half-plane n.dot(v_i-p)>=0.  Measuring that source's channel at all retained
vertices therefore discovers it regardless of its direction/radius.

The 25-point route below is an open Hamilton path made of lattice edges, from
the closest station to the origin.  Distinct lattice vertices are >= one edge
apart, so its length is the EXACT optimum for this particular fixed station
set and origin start.  This is not a lower bound for all possible P4 policies.
"""

import math


ARENA_RADIUS = 1800.0
RECEPTION_LOWER_BOUND = 1000.0
EDGE = 999.999
HEIGHT = math.sqrt(3.0) * EDGE / 2.0
OFFSET = (1.0 / 12.0, 1.0 / 12.0)
INTERSECTION_PADDING = 1e-7


def lattice_point(index):
    """Map the integer triangular-lattice index to a fixed physical point."""
    i, j = index
    u, v = OFFSET
    return (EDGE * (i - u + (j - v) / 2.0), HEIGHT * (j - v))


def _segment_distance_from_origin(a, b):
    dx, dy = b[0] - a[0], b[1] - a[1]
    t = -(a[0] * dx + a[1] * dy) / (dx * dx + dy * dy)
    t = max(0.0, min(1.0, t))
    return math.hypot(a[0] + t * dx, a[1] + t * dy)


def triangle_distance_from_origin(vertices):
    """Minimum Euclidean norm over the complete closed triangle."""
    a, b, c = vertices
    edges = ((a, b), (b, c), (c, a))
    crosses = [x[0] * y[1] - x[1] * y[0] for x, y in edges]
    if min(crosses) >= 0.0 or max(crosses) <= 0.0:
        return 0.0
    return min(_segment_distance_from_origin(x, y) for x, y in edges)


def _retained_triangle_indices():
    # If a triangle meets the radius-R disk, every vertex is within R+EDGE.
    # For q=a*e1+b*e2, |a| and |b| <= |q|/HEIGHT.  This gives a finite
    # enumeration bound for ALL such vertices; the extra unit includes either
    # choice of a cell's lower index.  No omitted cell can meet the target disk.
    bound = math.ceil((ARENA_RADIUS + EDGE) / HEIGHT + max(map(abs, OFFSET))) + 1
    retained = []
    for i in range(-bound, bound + 1):
        for j in range(-bound, bound + 1):
            for ids in (
                ((i, j), (i + 1, j), (i, j + 1)),
                ((i + 1, j), (i + 1, j + 1), (i, j + 1)),
            ):
                vertices = tuple(lattice_point(k) for k in ids)
                if triangle_distance_from_origin(vertices) <= ARENA_RADIUS + INTERSECTION_PADDING:
                    retained.append(ids)
    return tuple(retained)


MESH_TRIANGLE_INDICES = _retained_triangle_indices()
MESH_TRIANGLES = tuple(tuple(lattice_point(i) for i in tri) for tri in MESH_TRIANGLE_INDICES)

# This route visits each retained vertex once.  Every consecutive pair differs
# by one of (+/-1,0), (0,+/-1), (+/-1,-/+1), hence is exactly one lattice edge.
MESH_ROUTE_INDICES = (
    (0, 0), (-1, 0), (-2, 0), (-2, 1), (-2, 2), (-2, 3),
    (-1, 3), (-1, 2), (-1, 1), (0, 1), (0, 2), (1, 2),
    (2, 1), (1, 1), (1, 0), (2, 0), (3, -1), (3, -2),
    (2, -1), (2, -2), (1, -1), (1, -2), (0, -2), (-1, -1), (0, -1),
)
MESH_STATIONS = tuple(lattice_point(i) for i in MESH_ROUTE_INDICES)


def route_length(stations=MESH_STATIONS, start=(0.0, 0.0)):
    total = 0.0
    for station in stations:
        total += math.dist(start, station)
        start = station
    return total


MESH_CERTIFICATE = {
    "kind": "complete_equilateral_triangle_intersection",
    "arena_radius_m": ARENA_RADIUS,
    "guaranteed_reception_radius_m": RECEPTION_LOWER_BOUND,
    "edge_m": EDGE,
    "lattice_offset": OFFSET,
    "station_count": len(MESH_STATIONS),
    "retained_triangle_count": len(MESH_TRIANGLES),
    "open_route_from_origin_m": route_length(),
    "fixed_point_set_route_lower_bound_m": (
        (len(MESH_STATIONS) - 1) * EDGE + min(math.hypot(*p) for p in MESH_STATIONS)
    ),
    "coverage_basis": "continuous_triangle_diameter_and_convex_half_plane",
    "simulation_truth_used": False,
    "sampling_used_for_coverage": False,
    "completion_condition": "each still-unseen channel actually measured at every retained vertex",
}


def discovery_stations():
    """A fresh list, so a caller may own its scan progress without mutation here."""
    return list(MESH_STATIONS)


if __name__ == "__main__":
    import json
    print(json.dumps(MESH_CERTIFICATE, indent=2, ensure_ascii=False))
\end{lstlisting}
%</paper-block>

%<paper-block id="appendix.p024" type="paragraph">
\noindent\textbf{文件：}\texttt{02\_\allowbreak{}最终代码/\allowbreak{}第四问/\allowbreak{}discovery\_\allowbreak{}rings.py}\\
\textbf{功能：}问题四搜索基类加载的环形站点构造模块
%</paper-block>
%<paper-block id="appendix.code022" type="code">
\begin{lstlisting}[style=appendixcode]
"""Continuous directional discovery with 25 stations and a 17.94 km route.

The radius-1800 m target disk lies strictly inside a regular 16-gon whose
apothem is 1800.0001 m. A center, an 8-point inner ring and the outer polygon's
16 vertices triangulate that entire polygon into 32 triangles. Every triangle
edge is strictly shorter than the official 1000 m reception lower bound.

For any target p in a triangle, all three vertices are within 1000 m of p by
convexity of distance. Writing p as a convex combination of the vertices also
shows that at least one vertex lies in every closed directional half-plane
through p. Measuring each still-unseen channel at ALL its required stations
therefore guarantees discovery. No no_signal localization exclusion follows
from this proof. No simulator truth or sampled-point coverage is used.

This module uses only the Python standard library and has no P3 dependency.
"""
import math


ARENA_RADIUS=1800.0
RECEPTION_LOWER_BOUND=1000.0
INNER_RADIUS=999.999
OUTER_APOTHEM=1800.0001
OUTER_RADIUS=OUTER_APOTHEM/math.cos(math.pi/16)
CENTER=(0.0,0.0)


def _ring(radius,count):
    return tuple((radius*math.cos(math.tau*k/count),radius*math.sin(math.tau*k/count))
                 for k in range(count))


INNER_POINTS=_ring(INNER_RADIUS,8)
OUTER_POINTS=_ring(OUTER_RADIUS,16)

# Start at the origin; traverse the inner ring counterclockwise. Inner point 7
# aligns with outer point 14, so transfer radially and traverse the outer ring.
RING_STATIONS=(CENTER,)+INNER_POINTS+tuple(OUTER_POINTS[(14+k)%16] for k in range(16))


def _triangles():
    result=[]
    for a in range(8):
        ia,ib=INNER_POINTS[a],INNER_POINTS[(a+1)%8]
        oa,om,ob=(OUTER_POINTS[k%16] for k in (2*a,2*a+1,2*a+2))
        result.extend(((CENTER,ia,ib),(ia,oa,om),(ia,om,ib),(ib,om,ob)))
    return tuple(result)


RING_TRIANGLES=_triangles()


def route_length(stations=RING_STATIONS,start=CENTER):
    total=0.0
    for point in stations:
        total+=math.dist(start,point)
        start=point
    return total


RING_CERTIFICATE={
    'kind':'complete_concentric_ring_triangulation',
    'arena_radius_m':ARENA_RADIUS,
    'guaranteed_reception_radius_m':RECEPTION_LOWER_BOUND,
    'inner_radius_m':INNER_RADIUS,
    'outer_radius_m':OUTER_RADIUS,
    'outer_polygon_apothem_m':OUTER_APOTHEM,
    'station_count':len(RING_STATIONS),
    'retained_triangle_count':len(RING_TRIANGLES),
    'edge_m':max(math.dist(tri[i],tri[(i+1)%3]) for tri in RING_TRIANGLES for i in range(3)),
    'open_route_from_origin_m':route_length(),
    'coverage_basis':'entire_outer_polygon_triangulated; continuous_triangle_diameter_and_convex_half_plane',
    'simulation_truth_used':False,
    'sampling_used_for_coverage':False,
    'completion_condition':'each still-unseen channel actually measured at every ring station',
}


def discovery_stations():
    """Return a fresh list; scan progress remains owned by the caller."""
    return list(RING_STATIONS)


if __name__=='__main__':
    import json
    print(json.dumps(RING_CERTIFICATE,ensure_ascii=False,indent=2))
\end{lstlisting}
%</paper-block>

%<paper-block id="appendix.p025" type="paragraph">
\noindent\textbf{文件：}\texttt{02\_\allowbreak{}最终代码/\allowbreak{}第四问/\allowbreak{}discovery\_\allowbreak{}coverage.py}\\
\textbf{功能：}问题四未知朝向下的连续发现覆盖核验
%</paper-block>
%<paper-block id="appendix.code023" type="code">
\begin{lstlisting}[style=appendixcode]
"""Conservative continuous directional-discovery certificates for P4.

For observations at S, a target p is guaranteed detectable for every transmitting
half-plane iff p is in conv({s in S: |s-p| <= 1000}).  This verifier covers a
polygon strictly OUTSIDE the target disk with finite convex cells.  On each cell
Q it uses only stations within 1000 m of EVERY vertex of Q, then verifies
Q <= conv(those stations).  Convexity proves the whole cell, not a point sample.

Cell clipping and ambiguous orientation predicates use exact rational arithmetic
on the supplied floating-point coordinates.  Distance filtering keeps a 10 um
inward reserve.  Cells partition the outer polygon; failure to prove any cell,
including finite-budget exhaustion, returns False.  No simulator truth and no
single-no_signal localization exclusion are involved.

A certifier caches continuous cell witnesses, so additional observations cannot
invalidate a previously established certificate.  Instances own their caches;
use one per stream of observations (or identical unknown-channel point sets).
"""

from dataclasses import dataclass, field
from fractions import Fraction
from functools import lru_cache
import math

import numpy as np

from discovery_mesh import ARENA_RADIUS, lattice_point, triangle_distance_from_origin


RADIUS = 1000.0
DISTANCE_RESERVE = 1e-5
OUTER_RESERVE = 1e-5
OUTER_SIDES = 96


@lru_cache(maxsize=8192)
def _exact_point(point):
    return (Fraction.from_float(point[0]), Fraction.from_float(point[1]))


def _cross_exact(a, b, p):
    return (b[0] - a[0]) * (p[1] - a[1]) - (b[1] - a[1]) * (p[0] - a[0])


def _orientation(a, b, p, exact_p=None):
    x = (b[0] - a[0]) * (p[1] - a[1])
    y = (b[1] - a[1]) * (p[0] - a[0])
    cross = x - y
    uncertainty = 1e-12 * (abs(x) + abs(y) + 1.0)
    if cross > uncertainty:
        return 1
    if cross < -uncertainty:
        return -1
    value = _cross_exact(_exact_point(a), _exact_point(b), exact_p or _exact_point(p))
    return (value > 0) - (value < 0)


def _hull(points):
    points = sorted(set(points))
    if len(points) < 3:
        return points
    lower, upper = [], []
    for p in points:
        while len(lower) >= 2 and _orientation(lower[-2], lower[-1], p) <= 0:
            lower.pop()
        lower.append(p)
    for p in reversed(points):
        while len(upper) >= 2 and _orientation(upper[-2], upper[-1], p) <= 0:
            upper.pop()
        upper.append(p)
    return lower[:-1] + upper[:-1]


def _inside_hull(hull, vertices, exact_vertices):
    if len(hull) < 3:
        return False
    for a, b in zip(hull, hull[1:] + hull[:1]):
        for p, exact_p in zip(vertices, exact_vertices):
            if _orientation(a, b, tuple(p), exact_p) < 0:
                return False
    return True


def _clip(poly, a, b):
    """Exact closed left-half-plane clipping of a convex rational polygon."""
    if not poly:
        return ()
    result = []
    previous = poly[-1]
    previous_d = _cross_exact(a, b, previous)
    for current in poly:
        current_d = _cross_exact(a, b, current)
        if (current_d >= 0) != (previous_d >= 0):
            t = previous_d / (previous_d - current_d)
            result.append(tuple(previous[k] + t * (current[k] - previous[k]) for k in (0, 1)))
        if current_d >= 0:
            result.append(current)
        previous, previous_d = current, current_d
    dedup = []
    for p in result:
        if not dedup or p != dedup[-1]:
            dedup.append(p)
    if len(dedup) > 1 and dedup[0] == dedup[-1]:
        dedup.pop()
    return tuple(dedup)


def _positive_area(poly):
    if len(poly) < 3:
        return False
    return sum(a[0] * b[1] - a[1] * b[0] for a, b in zip(poly, poly[1:] + poly[:1])) > 0


@lru_cache(maxsize=1)
def _initial_geometry():
    radius = (ARENA_RADIUS + OUTER_RESERVE) / math.cos(math.pi / OUTER_SIDES)
    outer = tuple((radius * math.cos((2 * i + 1) * math.pi / OUTER_SIDES),
                   radius * math.sin((2 * i + 1) * math.pi / OUTER_SIDES))
                  for i in range(OUTER_SIDES))
    # The regular polygon has all edge supports > ARENA_RADIUS by 0.01 mm,
    # much more than the floating-point construction error.
    for a, b in zip(outer, outer[1:] + outer[:1]):
        support = (a[0] * b[1] - a[1] * b[0]) / math.dist(a, b)
        if support <= ARENA_RADIUS + OUTER_RESERVE / 2:
            raise ArithmeticError("outer polygon failed its strict enclosure bound")
    exact_outer = tuple(_exact_point(p) for p in outer)
    cells = []
    # At radius <1802 m, intersecting lattice triangles have every vertex
    # within 2802 m.  Their inverse lattice coordinates are <4 in magnitude.
    # This larger finite box contains all cells that could intersect the polygon.
    for i in range(-6, 7):
        for j in range(-6, 7):
            for ids in (((i, j), (i + 1, j), (i, j + 1)),
                        ((i + 1, j), (i + 1, j + 1), (i, j + 1))):
                triangle = tuple(lattice_point(index) for index in ids)
                if triangle_distance_from_origin(triangle) > radius + 1e-7:
                    continue
                exact_triangle = tuple(_exact_point(p) for p in triangle)
                poly = exact_outer
                for a, b in zip(exact_triangle, exact_triangle[1:] + exact_triangle[:1]):
                    poly = _clip(poly, a, b)
                if _positive_area(poly):
                    cells.append(poly)
    return outer, exact_outer, tuple(cells)


@dataclass
class _Cell:
    exact: tuple
    depth: int = 0
    witnesses: list = field(default_factory=list)
    children: tuple = ()

    def __post_init__(self):
        self.vertices = np.array([[float(x), float(y)] for x, y in self.exact])

    def split(self):
        if self.children:
            return self.children
        bounds = [(min(p[k] for p in self.exact), max(p[k] for p in self.exact)) for k in (0, 1)]
        axis = max((0, 1), key=lambda k: bounds[k][1] - bounds[k][0])
        midpoint = sum(bounds[axis], Fraction()) / 2
        if axis == 0:
            a, b = (midpoint, Fraction(0)), (midpoint, Fraction(1))
        else:
            a, b = (Fraction(1), midpoint), (Fraction(0), midpoint)
        halves = (_clip(self.exact, a, b), _clip(self.exact, b, a))
        self.children = tuple(_Cell(p, self.depth + 1) for p in halves if _positive_area(p))
        return self.children


class CoverageCertifier:
    """Finite sufficient certificate; False means unproved, never safe to stop."""

    def __init__(self, max_cells=2400, max_depth=48):
        self.max_cells = int(max_cells)
        self.max_depth = int(max_depth)
        self.outer, self.exact_outer, cells = _initial_geometry()
        self.cells = tuple(_Cell(poly) for poly in cells)
        self.completed_witnesses = []
        self.last_diagnostics = {}

    def certified_complete(self, points, *, max_cells=None):
        points = tuple(sorted(set(tuple(map(float, p)) for p in points)))
        if any(len(p) != 2 or not all(math.isfinite(x) for x in p) for p in points):
            raise ValueError("observation points must be finite (x, y) pairs")
        point_set = frozenset(points)
        stats = dict(complete=False, point_count=len(points), cells_checked=0,
                     witness_cache_hits=0, cells_split=0, max_depth=0,
                     reason="not_started", coverage_uses_sampling=False)
        self.last_diagnostics = stats
        if any(witness <= point_set for witness in self.completed_witnesses):
            stats.update(complete=True, reason="complete_witness_cache", witness_cache_hits=1)
            return True
        if len(points) < 3 or not _inside_hull(_hull(points), self.outer, self.exact_outer):
            stats["reason"] = "outer_polygon_not_in_global_hull"
            return False
        budget = self.max_cells if max_cells is None else int(max_cells)
        array = np.asarray(points)
        stack = list(reversed(self.cells))
        while stack:
            if stats["cells_checked"] >= budget:
                stats["reason"] = "finite_cell_budget_exhausted"
                return False
            cell = stack.pop()
            stats["cells_checked"] += 1
            stats["max_depth"] = max(stats["max_depth"], cell.depth)
            if any(witness <= point_set for witness in cell.witnesses):
                stats["witness_cache_hits"] += 1
                continue
            delta = array[:, None, :] - cell.vertices[None, :, :]
            maximum_squared = np.max(np.sum(delta * delta, axis=2), axis=1)
            common = [p for p, eligible in zip(points, maximum_squared <= (RADIUS - DISTANCE_RESERVE) ** 2) if eligible]
            hull = _hull(common)
            if _inside_hull(hull, cell.vertices, cell.exact):
                # Only the hull stations are needed as a replayable whole-cell
                # witness.  Their range was bounded over every cell vertex.
                witness = frozenset(hull)
                cell.witnesses = [witness] + cell.witnesses[:1]
                continue
            if cell.depth >= self.max_depth:
                stats["reason"] = "finite_depth_budget_exhausted"
                return False
            children = cell.split()
            if not children:
                stats["reason"] = "non_subdividable_unproved_cell"
                return False
            stats["cells_split"] += 1
            stack.extend(reversed(children))
        self.completed_witnesses = [point_set] + self.completed_witnesses[:7]
        stats.update(complete=True, reason="all_continuous_cells_proved")
        return True

    def certified_omit(self, points, remaining_stations, station_to_omit, *, max_cells=None):
        """Certify the combined actual+future set after omitting one station.

        A True result proves a planned route sufficient.  It does NOT prove that
        future measurements have already happened, or mark any channel absent.
        """
        omitted = tuple(map(float, station_to_omit))
        remaining = [tuple(map(float, p)) for p in remaining_stations]
        if omitted not in remaining:
            raise ValueError("station_to_omit is not in remaining_stations")
        return self.certified_complete(list(points) + [p for p in remaining if p != omitted], max_cells=max_cells)


@lru_cache(maxsize=1)
def _default_certifier():
    return CoverageCertifier()


def certified_complete(points, *, max_cells=None):
    return _default_certifier().certified_complete(points, max_cells=max_cells)


def certified_omit(points, remaining_stations, station_to_omit, *, max_cells=None):
    return _default_certifier().certified_omit(points, remaining_stations, station_to_omit, max_cells=max_cells)
\end{lstlisting}
%</paper-block>

%<paper-block id="appendix.p026" type="paragraph">
\noindent\textbf{文件：}\texttt{02\_\allowbreak{}最终代码/\allowbreak{}第四问/\allowbreak{}discovery\_\allowbreak{}compact.py}\\
\textbf{功能：}问题四21站紧凑发现构造
%</paper-block>
%<paper-block id="appendix.code024" type="code">
\begin{lstlisting}[style=appendixcode]
"""Continuously certified P4 discovery with 21 fixed stations.

Construction: center + eight inner points at 999.999 m + twelve outer points
whose regular polygon apothem is 1800.0001 m.  Unlike the earlier ring25 proof,
this construction does not claim every triangulation edge is below 1000 m.
Its certificate is the stronger common-near-station convex-hull condition,
proved over every cell of an outer enclosure of the continuous target disk.

Production import uses only the standard library and does not rerun the proof.
build_coverage_certifier() reproduces its exact domain clipping without any
search script or result JSON.  Reproduction dependencies are the local
discovery_coverage.py, discovery_mesh.py and NumPy; no P3 module is needed.
"""

import hashlib
import json
import math


ARENA_RADIUS = 1800.0
RECEPTION_LOWER_BOUND = 1000.0
INNER_RADIUS = 999.999
INNER_COUNT = 8
OUTER_COUNT = 12
OUTER_APOTHEM = 1800.0001
OUTER_RADIUS = OUTER_APOTHEM / math.cos(math.pi / OUTER_COUNT)
CENTER = (0.0, 0.0)


def _ring(radius, count):
    return tuple((radius * math.cos(k * math.tau / count),
                  radius * math.sin(k * math.tau / count)) for k in range(count))


INNER_POINTS = _ring(INNER_RADIUS, INNER_COUNT)
OUTER_POINTS = _ring(OUTER_RADIUS, OUTER_COUNT)
COMPACT_POINTS = (CENTER,) + INNER_POINTS + OUTER_POINTS
COMPACT_ROUTE_INDICES = (0, 2, 1, 8, 7, 6, 5, 4, 3, 12, 13, 14, 15, 16, 17, 18, 19, 20, 9, 10, 11)
COMPACT_STATIONS = tuple(COMPACT_POINTS[index] for index in COMPACT_ROUTE_INDICES)


def route_length(stations=COMPACT_STATIONS, start=CENTER):
    total = 0.0
    for point in stations:
        total += math.dist(start, point)
        start = point
    return total


PROOF_PARAMETERS = {
    "kind": "continuous_common_range_convex_hull_cells",
    "arena_radius_m": ARENA_RADIUS,
    "reception_lower_bound_m": RECEPTION_LOWER_BOUND,
    "outer_polygon_sides": 96,
    "outer_polygon_apothem_reserve_m": 1e-5,
    "distance_inward_reserve_m": 1e-5,
    "domain": "exact intersection of outer 96-gon and strictly enclosing station hull",
    "max_cells": 10000,
    "max_depth": 56,
    "clipping": "exact rational closed half planes",
    "acceptance": "every continuous cell lies in hull of stations within 1000m of every cell vertex",
}


def _hash(value):
    data = json.dumps(value, ensure_ascii=True, sort_keys=True, separators=(",", ":")).encode("ascii")
    return hashlib.sha256(data).hexdigest()


# Coordinate hex strings identify the exact binary floats used in the proof.
GEOMETRY_SHA256 = _hash({"points_hex": [[x.hex(), y.hex()] for x, y in COMPACT_POINTS],
                       "route_indices": COMPACT_ROUTE_INDICES})
CERTIFICATE_SHA256 = _hash({"geometry_sha256": GEOMETRY_SHA256, "proof_parameters": PROOF_PARAMETERS})
COMPACT_CERTIFICATE = {
    **PROOF_PARAMETERS,
    "inner_count": INNER_COUNT,
    "outer_count": OUTER_COUNT,
    "inner_radius_m": INNER_RADIUS,
    "outer_polygon_apothem_m": OUTER_APOTHEM,
    "outer_radius_m": OUTER_RADIUS,
    "station_count": len(COMPACT_STATIONS),
    "open_route_from_origin_m": route_length(),
    "geometry_sha256": GEOMETRY_SHA256,
    "certificate_sha256": CERTIFICATE_SHA256,
    "coverage_uses_sampling": False,
    "simulation_truth_used": False,
    "completion_condition": "each still-unseen channel actually measured at every required station",
}


def discovery_stations():
    return list(COMPACT_STATIONS)


def build_coverage_certifier(*, max_cells=10000, max_depth=56):
    """Create a reusable verifier, independently of search outputs.

    Both enclosures contain the entire radius-1800 disk; their exact convex
    intersection still contains it.  The candidate's hull reserve is checked
    before it is allowed to restrict the proof domain.  A caller can later test
    ANY actual point set against this fixed enclosing domain.
    """
    from discovery_coverage import (CoverageCertifier, DISTANCE_RESERVE,
                                    OUTER_RESERVE, OUTER_SIDES, _Cell, _clip,
                                    _exact_point, _hull, _initial_geometry, _positive_area)

    if (OUTER_SIDES != PROOF_PARAMETERS["outer_polygon_sides"]
            or OUTER_RESERVE != PROOF_PARAMETERS["outer_polygon_apothem_reserve_m"]
            or DISTANCE_RESERVE != PROOF_PARAMETERS["distance_inward_reserve_m"]):
        raise RuntimeError("coverage implementation changed; refresh the compact proof parameters")
    hull = _hull(COMPACT_POINTS)
    for a, b in zip(hull, hull[1:] + hull[:1]):
        support = (a[0] * b[1] - a[1] * b[0]) / math.dist(a, b)
        if support <= ARENA_RADIUS + 0.00005:
            raise ArithmeticError("compact station hull no longer strictly encloses the target disk")
    exact_hull = tuple(_exact_point(p) for p in hull)
    _, original_outer, original_cells = _initial_geometry()

    def intersect(poly):
        for a, b in zip(exact_hull, exact_hull[1:] + exact_hull[:1]):
            poly = _clip(poly, a, b)
        return poly

    checker = CoverageCertifier(max_cells=max_cells, max_depth=max_depth)
    checker.exact_outer = intersect(original_outer)
    checker.outer = tuple(tuple(map(float, p)) for p in checker.exact_outer)
    checker.cells = tuple(_Cell(poly) for cell in original_cells
                          if _positive_area(poly := intersect(cell)))
    return checker


if __name__ == "__main__":
    import argparse
    parser = argparse.ArgumentParser()
    parser.add_argument("--verify", action="store_true")
    args = parser.parse_args()
    output = dict(COMPACT_CERTIFICATE)
    if args.verify:
        certifier = build_coverage_certifier()
        output["reproved_complete"] = certifier.certified_complete(COMPACT_STATIONS)
        output["reproof_diagnostics"] = certifier.last_diagnostics
        if not output["reproved_complete"]:
            raise SystemExit(json.dumps(output, indent=2))
    print(json.dumps(output, indent=2))
\end{lstlisting}
%</paper-block>

%<paper-block id="appendix.p027" type="paragraph">
\noindent\textbf{文件：}\texttt{02\_\allowbreak{}最终代码/\allowbreak{}第四问/\allowbreak{}discovery\_\allowbreak{}prune.py}\\
\textbf{功能：}问题四带覆盖证书的后续扫描站删减
%</paper-block>
%<paper-block id="appendix.code025" type="code">
\begin{lstlisting}[style=appendixcode]
"""Finite, continuous-certified replacement of future P4 discovery stations.

This component proves a FUTURE plan only.  A clear at extra_point is not a
measurement of any unknown channel.  Before deleting a pending station, the
caller must arrange real measurements at extra_point for EVERY unknown channel
whose plan relies on it.  Discovery completion still requires separate evidence
from each channel's actual completed measurements.

There is no simulator/source-count input, no source truth and no individual
no_signal localization exclusion.  Geometry uncertainty or finite-budget failure
keeps the original future plan.  At most eight deletion candidates are tested.
"""

import math

from discovery_compact import build_coverage_certifier


MAX_CANDIDATES = 8
MAX_CELLS_PER_CERTIFICATE = 2400
MAX_CERTIFICATE_DEPTH = 48
SAME_POINT_TOLERANCE = 1e-6


def _point(value):
    result = tuple(map(float, value))
    if len(result) != 2 or not all(math.isfinite(x) for x in result):
        raise ValueError("points must be finite (x, y) pairs")
    return result


def _unique_points(points):
    # Tiny coordinate serialization differences do not constitute new evidence.
    unique = []
    for point in points:
        if not any(math.dist(point, existing) <= SAME_POINT_TOLERANCE for existing in unique):
            unique.append(point)
    return unique


def route_length(points, start):
    length = 0.0
    for point in points:
        length += math.dist(start, point)
        start = point
    return length


def _saving(entries, index, start):
    previous = entries[index - 1][1] if index else start
    point = entries[index][1]
    saving = math.dist(previous, point)
    if index + 1 < len(entries):
        following = entries[index + 1][1]
        saving += math.dist(point, following) - math.dist(previous, following)
    return max(0.0, saving)


def choose_replacement(actual_points, pending, extra_point, start, checker=None):
    """Return {removed, remaining, diagnostics}; no actual evidence is mutated.

    actual_points must contain ONLY locations at which ALL currently unknown
    channels have already really been measured.  extra_point is a hypothetical
    common measurement location, even if the robot has already cleared there.
    pending order is preserved.  Accepted deletions are tested cumulatively.
    """
    actual = _unique_points([_point(point) for point in actual_points])
    original = [_point(point) for point in pending]
    extra = _point(extra_point)
    start = _point(start)
    equivalent = next((point for point in actual if math.dist(point, extra) <= SAME_POINT_TOLERANCE), None)
    proof_extra = equivalent if equivalent is not None else extra
    hypothetical_common = actual + ([] if equivalent is not None else [proof_extra])
    entries = list(enumerate(original))
    removed = []
    stats = {
        "kind": "continuous_future_plan_replacement",
        "future_plan_only": True,
        "actual_completion": "not_assessed",
        "actual_common_unique_points": len(actual),
        "extra_point": extra,
        "extra_is_new_location": equivalent is None,
        "extra_requires_actual_unknown_channel_measurements": equivalent is None,
        "pending_before": len(original),
        "pending_after": len(original),
        "max_candidate_attempts": MAX_CANDIDATES,
        "max_cells_per_certificate": MAX_CELLS_PER_CERTIFICATE,
        "attempts": [],
        "future_plan_proved": False,
        "route_before_m": route_length(original, start),
        "route_after_m": route_length(original, start),
        "route_saving_m": 0.0,
        "coverage_uses_sampling": False,
    }
    if checker is None:
        try:
            checker = build_coverage_certifier(max_cells=MAX_CELLS_PER_CERTIFICATE,
                                              max_depth=MAX_CERTIFICATE_DEPTH)
        except (ArithmeticError, ValueError, RuntimeError) as exc:
            stats["initial_future_certificate"] = {"reason": "certificate_construction_error", "error": str(exc)}
            stats["reason"] = "initial_future_plan_not_proved"
            return dict(removed=[], remaining=list(original), diagnostics=stats)

    def prove(points):
        try:
            passed = checker.certified_complete(points, max_cells=MAX_CELLS_PER_CERTIFICATE)
            diagnostic = dict(checker.last_diagnostics)
            return bool(passed), diagnostic
        except (ArithmeticError, ValueError, RuntimeError) as exc:
            # A certificate construction/calculation failure supplies no license
            # to remove a station.  The caller gets the reason rather than a crash.
            return False, {"reason": "certificate_error", "error": str(exc)}

    initial_passed, initial_diagnostic = prove(hypothetical_common + original)
    stats["initial_future_certificate"] = initial_diagnostic
    if not initial_passed:
        stats["reason"] = "initial_future_plan_not_proved"
        return dict(removed=[], remaining=list(original), diagnostics=stats)
    stats["future_plan_proved"] = True
    attempted_ids = set()
    while entries and len(attempted_ids) < MAX_CANDIDATES:
        candidates = [i for i, (identity, _) in enumerate(entries) if identity not in attempted_ids]
        if not candidates:
            break
        # Nearby interior points can be replaced by a clear stop. Outer support
        # vertices often have a large route saving but cannot be replaced by an
        # interior point; do not spend the entire finite budget on those first.
        index = min(candidates, key=lambda i: (math.dist(extra, entries[i][1]),
                                               -_saving(entries, i, start), entries[i][0]))
        identity, point = entries[index]
        attempted_ids.add(identity)
        estimated_saving = _saving(entries, index, start)
        tentative = [p for j, (_, p) in enumerate(entries) if j != index]
        passed, diagnostic = prove(hypothetical_common + tentative)
        stats["attempts"].append({
            "original_pending_index": identity,
            "point": point,
            "prospective_route_saving_m": estimated_saving,
            "distance_to_extra_m": math.dist(extra, point),
            "accepted": passed,
            "certificate": diagnostic,
        })
        if passed:
            removed.append(point)
            entries.pop(index)
    remaining = [point for _, point in entries]
    stats["pending_after"] = len(remaining)
    stats["route_after_m"] = route_length(remaining, start)
    stats["route_saving_m"] = stats["route_before_m"] - stats["route_after_m"]
    stats["reason"] = "certified_future_replacement" if removed else "no_deletion_proved"
    return dict(removed=removed, remaining=remaining, diagnostics=stats)
\end{lstlisting}
%</paper-block>

%<paper-block id="appendix.p028" type="paragraph">
\noindent\textbf{文件：}\texttt{02\_\allowbreak{}最终代码/\allowbreak{}第四问/\allowbreak{}local\_\allowbreak{}hunt.py}\\
\textbf{功能：}问题四有限局部补测规划
%</paper-block>
%<paper-block id="appendix.code026" type="code">
\begin{lstlisting}[style=appendixcode]
"""Bounded P4 local-bearing planner; imports only the bundled P4 modules.

The caller owns measurement results, Tracks, clear plans and the global survey
cursor. This component only proposes at most eight measurements per channel.
One-bearing moves are explicitly speculative, never coverage certificates.
With two positive positions, every proposed point is a convex combination of
positive positions, hence remains in the convex reception disk/half-plane.
No ``no_signal`` result is accepted here and no spatial exclusions are inferred.
"""
import math
from collections import Counter

from strategy import Action, p1
from strategy_p4 import ARENA_R, EPS_DEG, R_EFF, distance


def _unique(points, tolerance=0.5):
    result = []
    for point in points:
        point = tuple(point)
        if not any(distance(point, old) <= tolerance for old in result):
            result.append(point)
    return result


def _ray_interval(origin, direction):
    """Central-bearing interval in the public arena/range; ranking only."""
    dot = origin[0] * direction[0] + origin[1] * direction[1]
    delta = dot * dot + ARENA_R * ARENA_R - sum(v * v for v in origin)
    if delta < 0:
        return 0.0, R_EFF
    root = math.sqrt(delta)
    lo, hi = max(0.0, -dot - root), min(R_EFF, -dot + root)
    return (lo, hi) if hi > lo else (0.0, R_EFF)


def _centroid(poly):
    # Vertex average stays in a convex polygon. It is only a cost hypothesis.
    return tuple(sum(p[k] for p in poly) / len(poly) for k in (0, 1))


def _finish_estimate(poly, position):
    """Ranking estimate includes travel and success/miss clear costs.

    This is deliberately NOT an optical coverage guarantee. The caller must
    still obtain a full continuous coverage certificate before using a plan.
    """
    center = _centroid(poly)
    radius = max(distance(center, q) for q in poly)
    diameter, _, _ = p1.polygon_diameter(poly)
    count = max(1, math.ceil(diameter / 32.0))
    move = max(0.0, distance(position, center) - max(0.0, 20.0 - radius)) / 5.0
    # A lower-dimensional strip estimate is enough for ranking candidate
    # measurements; cost/correctness claims never use this hypothetical count.
    return move + 5.0 + 3.0 * (count - 1) + max(0, count - 1) * 24.0 / 5.0


class LocalHuntPlanner:
    """Simple interface: ``next_probe(track, state, channel) -> Action | None``.

    The issued action itself consumes budget. No callback is required: the next
    call reads the caller's current positive records. Calling again after a
    no-signal reply therefore advances to a distinct point without modifying
    the feasible region. ``None`` means hand the channel back to the existing
    global survey/localization fallback, never clear/skip/mark it done.
    """

    def __init__(self, max_probes=8):
        if not 1 <= int(max_probes) <= 8:
            raise ValueError('max_probes must be within 1..8')
        self.max_probes = int(max_probes)
        self.history = {}
        self.events = Counter()
        self.last_reason = {}
        self.cost_comparisons = []

    def _first_candidates(self, track):
        origin, bearing = track.records[0]
        angle = math.radians(bearing)
        u = (math.cos(angle), math.sin(angle))
        v = (-u[1], u[0])
        lo, hi = _ray_interval(origin, u)
        hop = lo + 0.4 * (hi - lo)
        hop = max(100.0, min(1100.0, hop))
        lateral = min(180.0, max(60.0, hop * 0.3))
        # Opposite lateral side first after a miss; then retreat along the
        # original ray. None of these no-signal attempts shortens that ray.
        offsets = ((hop, lateral), (hop, -lateral),
                   (hop * 0.5, -lateral * 0.75),
                   (hop * 0.5, lateral * 0.75),
                   (hop / 6.0, lateral * 0.5),
                   (hop / 6.0, -lateral * 0.5))
        return [(origin[0] + forward * u[0] + side * v[0],
                 origin[1] + forward * u[1] + side * v[1])
                for forward, side in offsets]

    def _safe_candidates(self, track, used):
        positions = _unique(pos for pos, _ in track.records)
        pairs = sorted(((distance(a, b), a, b)
                        for i, a in enumerate(positions)
                        for b in positions[i + 1:]), reverse=True)
        candidates = []
        for _, a, b in pairs[:12]:
            for fraction in (0.2, 0.4, 0.6, 0.8):
                candidates.append(tuple(a[k] * (1.0 - fraction) + b[k] * fraction
                                        for k in (0, 1)))
        if len(positions) >= 3:
            for group in (positions[:3], positions[-3:], positions[::2][:3]):
                if len(group) == 3:
                    candidates.append(_centroid(group))
        return [p for p in _unique(candidates)
                if all(distance(p, old) > 0.5 for old in positions + used)]

    def _score(self, point, track, state, channel):
        immediate = distance(state.pos, point) / 5.0 + 5.0 + float(state.ch != channel)
        poly = track.poly
        if len(poly) < 3:
            # No available polygon: ranking by measured travel remains valid.
            return immediate + distance(point, track.records[-1][0]) / 5.0 + 5.0
        center = _centroid(poly)
        # Spatial sample hypotheses select a cheap action only. All safety and
        # any later optical coverage use the unchanged complete input polygon.
        _, a, b = p1.polygon_diameter(poly)
        hypotheses = (center, tuple(0.75 * center[k] + 0.25 * a[k] for k in (0, 1)),
                      tuple(0.75 * center[k] + 0.25 * b[k] for k in (0, 1)))
        finishes = []
        for target in hypotheses:
            if distance(point, target) < 1e-6:
                finishes.append(5.0)
                continue
            theta = math.degrees(math.atan2(target[1] - point[1], target[0] - point[0]))
            hypothetical = p1.hpi_intersect(poly[:], p1.make_sector_halfplanes(point, theta, EPS_DEG))
            finishes.append(_finish_estimate(hypothetical or poly, point))
        return immediate + sum(finishes) / len(finishes)

    def next_probe(self, track, state, channel):
        used = self.history.setdefault(channel, [])
        if not track.records or track.near is not None:
            self.last_reason[channel] = 'no_bearing_or_near'
            return None
        if len(used) >= self.max_probes:
            self.last_reason[channel] = 'probe_budget_exhausted'
            self.events['budget_fallback'] += 1
            return None
        positions = _unique(pos for pos, _ in track.records)
        excluded = used + list(getattr(track, 'probe_history', [])) + positions
        if len(positions) == 1:
            candidates = [p for p in self._first_candidates(track)
                          if all(distance(p, old) > 0.5 for old in excluded)]
            # Keep bounded planned side/retreat order after a miss. Initial
            # reflection is selected by travel from the actual current point.
            if not used and len(candidates) >= 2:
                candidates[:2] = sorted(candidates[:2], key=lambda p: (distance(state.pos, p), p))
            kind = 'speculative_first_pair'
        else:
            candidates = self._safe_candidates(track, excluded)
            candidates.sort(key=lambda p: (self._score(p, track, state, channel), p))
            kind = 'safe_convex_probe'
        if not candidates:
            self.last_reason[channel] = 'distinct_candidates_exhausted'
            self.events['candidate_fallback'] += 1
            return None
        point = candidates[0]
        used.append(point)
        self.last_reason[channel] = kind
        self.events[kind] += 1
        self.cost_comparisons.append(dict(channel=channel, kind=kind, point=point,
                                         estimated_total_s=self._score(point, track, state, channel),
                                         issued_index=len(used)))
        return Action('measure', point, channel)

    def diagnostics(self):
        return dict(max_probes=self.max_probes, triggers=dict(self.events),
                    history={ch: list(points) for ch, points in self.history.items()},
                    last_reason=dict(self.last_reason),
                    cost_comparisons=list(self.cost_comparisons))
\end{lstlisting}
%</paper-block>

%<paper-block id="appendix.p029" type="paragraph">
\noindent\textbf{文件：}\texttt{02\_\allowbreak{}最终代码/\allowbreak{}第四问/\allowbreak{}strategy\_\allowbreak{}hunt.py}\\
\textbf{功能：}问题四发现、局部定位与清除协同
%</paper-block>
%<paper-block id="appendix.code027" type="code">
\begin{lstlisting}[style=appendixcode]
"""P4 candidate: certified triangular discovery and bounded early local work.

No hidden source count/radius, P3 imports, or no-signal localization exclusions.
The old FastP4 and FinishP4 are left runnable unchanged.
"""
import math
from collections import Counter
from shapely.geometry import Polygon

from strategy import Action
from strategy_fast import FastP4, nearest_path
from strategy_finish import FinishP4
from strategy_p4 import CHANNELS, DISCOVERY_STATIONS, distance
from strategy_cost import optical_cover
from discovery_mesh import MESH_STATIONS, MESH_CERTIFICATE
from discovery_rings import RING_STATIONS, RING_CERTIFICATE
from discovery_compact import COMPACT_STATIONS, COMPACT_CERTIFICATE, build_coverage_certifier
from local_hunt import LocalHuntPlanner
from discovery_coverage import CoverageCertifier


class HuntP4(FinishP4):
    def __init__(self, hunt=True, early_hunt=False, grid='compact', adapt=True, insertion=True):
        super().__init__(enabled=True, trim_corners=False)
        self.hunt = bool(hunt)
        self.early_hunt = bool(early_hunt)
        self.adapt = bool(adapt and hunt)
        self.insertion = bool(insertion)
        self.coverage_certifier = build_coverage_certifier() if grid == 'compact' else CoverageCertifier()
        self.reanchored = False
        self._anchor_checks = set()
        if grid not in ('mesh', 'rings', 'compact'):
            raise ValueError('grid must be mesh, rings or compact')
        self.grid = grid
        self.core_stations, self.core_certificate = {
            'mesh': (MESH_STATIONS, MESH_CERTIFICATE),
            'rings': (RING_STATIONS, RING_CERTIFICATE),
            'compact': (COMPACT_STATIONS, COMPACT_CERTIFICATE),
        }[grid]
        self.name = ('HuntP4_' if hunt else 'MeshP4_') + grid + str(len(self.core_stations))
        self.discovery_stations = list(self.core_stations) + nearest_path(
            set(DISCOVERY_STATIONS), self.core_stations[-1])
        self.ever_observed_channels = set()
        self.actual_scan_points = {ch: set() for ch in CHANNELS}
        self.local_planner = LocalHuntPlanner(max_probes=8)
        self.active_hunt = None
        self.early_plan = None
        self.held_action = None
        self.hunt_versions = {}
        self.hunt_events = Counter()
        self._issued_local = set()

    def _needs_sample(self, ch):
        if len(self.ever_observed_channels) >= 16 and ch not in self.ever_observed_channels:
            return False
        if self.hunt and ch in self.tracks and ch not in self.cleared:
            certificate = self._certificate(ch)
            if certificate and len(certificate['points']) <= 3:
                return False
        return super()._needs_sample(ch)

    def _next_discovery_action(self):
        while self.discovery_station_index < len(self.discovery_stations):
            if self._station_channels and self.discovery_channel_index >= len(self._station_channels):
                self.discovery_station_index += 1
                self.discovery_channel_index = 0
                self._station_channels = []
                continue
            if self.discovery_station_index >= len(self.core_stations) and not self.core_complete:
                # Check actual per-channel evidence, never merely a cursor value.
                for ch in set(CHANNELS) - self.ever_observed_channels:
                    if len(self.ever_observed_channels) < 16 and not set(self.core_stations) <= self.actual_scan_points[ch]:
                        if not self.reanchored or not self.coverage_certifier.certified_complete(self.actual_scan_points[ch]):
                            raise RuntimeError('incomplete_directional_discovery_certificate')
                self.core_complete = True
                self.absent = set(CHANNELS) - self.ever_observed_channels
                self.hunt_events['mesh_coverage_completed'] += 1
                if self.hunt:
                    self._finish_discovery()
                    return None
            if not self._station_channels:
                self._station_channels = [ch for ch in CHANNELS if self._needs_sample(ch)]
                self.discovery_channel_index = 0
            if not self._station_channels:
                self.discovery_station_index += 1
                continue
            ch = self._station_channels[self.discovery_channel_index]
            self.discovery_channel_index += 1
            if self._needs_sample(ch):
                return Action('measure', self.discovery_stations[self.discovery_station_index], ch)
        self._finish_discovery()
        return None

    def _cover(self, ch, pos):
        track = self.tracks[ch]
        if track.near is not None:
            return dict(points=[tuple(track.near)], total_s=distance(pos, track.near) / 5.0 + 5.0)
        if len(track.poly) < 3:
            return None
        region = Polygon(track.poly)
        if not region.is_valid or region.is_empty:
            return None
        # All vertices in the convex clear disk proves the COMPLETE polygon.
        if max(distance(pos, v) for v in track.poly) < 19.99999:
            return dict(points=[tuple(pos)], total_s=5.0)
        if track.radius < 19.99999:
            d = distance(pos, track.center)
            shift = min(d, 19.99999 - track.radius)
            point = tuple(track.center[i] + (pos[i] - track.center[i]) * shift / d for i in (0, 1)) if d else tuple(pos)
            return dict(points=[point], total_s=distance(pos, point) / 5.0 + 5.0)
        return optical_cover(region, pos, limit=6)

    def _start_plan(self, ch, plan):
        if not plan or not plan['points']:
            raise RuntimeError('empty_early_optical_plan')
        self.early_plan = dict(channel=ch, points=list(plan['points']))
        self.hunt_events['early_plan_started'] += 1
        return Action('clear', self.early_plan['points'][0], ch)

    def _hunt_action(self, state):
        ch = self.active_hunt
        if ch in self.cleared:
            self.active_hunt = None
            return None
        plan = self._cover(ch, state.pos)
        if plan is not None:
            return self._start_plan(ch, plan)
        action = self.local_planner.next_probe(self.tracks[ch], state, ch)
        if action is None:
            self.hunt_events['bounded_hunt_fallback'] += 1
            self.hunt_versions[ch] = len(self.tracks[ch].records)
            self.active_hunt = None
        return action

    def _reanchor_scan(self, state):
        action = self.held_action
        index = self.discovery_station_index
        if (not self.adapt or action is None or action.kind != 'measure'
                or self.phase != 'discovery' or index >= len(self.core_stations)
                or self.discovery_channel_index > 1 or distance(state.pos, action.pos) < 1e-6):
            return
        key = (index, tuple(state.pos))
        if key in self._anchor_checks:
            return
        self._anchor_checks.add(key)
        # Source-clear detours change the starting position. Reorder only whole,
        # still-unvisited scan stations; every channel's required set is preserved.
        tail = self.discovery_stations[index:len(self.core_stations)]
        def path_length(points):
            return sum(distance(a, b) for a, b in zip([state.pos] + points, points))
        best = tail[:]
        trial = nearest_path(tail, state.pos)
        if path_length(trial) < path_length(best):
            best = trial
        for _ in range(24):
            improvement = None
            for i in range(len(best) - 1):
                a = state.pos if i == 0 else best[i - 1]
                for j in range(i + 1, len(best)):
                    gain = distance(a, best[i]) - distance(a, best[j])
                    if j + 1 < len(best):
                        gain += distance(best[j], best[j + 1]) - distance(best[i], best[j + 1])
                    if gain > 1e-7 and (improvement is None or gain > improvement[0]):
                        improvement = (gain, i, j)
            if improvement is None:
                break
            _, i, j = improvement
            best[i:j + 1] = reversed(best[i:j + 1])
        if path_length(best) + 1e-7 < path_length(tail):
            self.discovery_stations[index:len(self.core_stations)] = best
            action = Action('measure', best[0], action.ch)
            self.held_action = action
            self.hunt_events['remaining_scan_route_shortened'] += 1
        unknown = set(CHANNELS) - self.ever_observed_channels
        if not unknown or len(self.ever_observed_channels) >= 16:
            return
        actual = set.intersection(*(self.actual_scan_points[ch] for ch in unknown))
        future = set(self.discovery_stations[index + 1:len(self.core_stations)])
        # This certificate only changes a future plan. Absence later requires
        # a second certificate built solely from actual channel measurements.
        if self.coverage_certifier.certified_complete(actual | future | {tuple(state.pos)}):
            self.discovery_stations[index] = tuple(state.pos)
            self.held_action = Action('measure', tuple(state.pos), action.ch)
            self.reanchored = True
            self.hunt_events['clear_stop_replaces_scan_station'] += 1

    def step(self, state):
        if self.done:
            return Action('done')
        if self.early_plan:
            self.action_stage = 'early_clear'
            return Action('clear', self.early_plan['points'][0], self.early_plan['channel'])
        # A near reply always permits a zero-movement clear without touching scan cursors.
        for ch, track in self.tracks.items():
            if ch not in self.cleared and track.near is not None and distance(track.near, state.pos) < 1e-6:
                self.action_stage = 'near_clear'
                return self._start_plan(ch, self._cover(ch, state.pos))
        if self.active_hunt is not None:
            action = self._hunt_action(state)
            if action is not None:
                self.action_stage = 'early_clear' if action.kind == 'clear' else 'early_localization'
                return action
        self._reanchor_scan(state)
        action = self.held_action
        self.held_action = None
        if action is None or action.ch in self.cleared:
            action = super().step(state)
        # Interpose only between scan stations, after the full queued station was read.
        if self.hunt and self.phase == 'discovery' and action.kind == 'measure' and distance(state.pos, action.pos) > 1e-6:
            ready = []
            for ch, track in self.tracks.items():
                if ch in self.cleared:
                    continue
                plan = self._cover(ch, state.pos)
                if plan is not None:
                    last = plan['points'][-1]
                    extra = plan['total_s'] + (distance(last, action.pos) - distance(state.pos, action.pos)) / 5.0
                    future_best = math.inf
                    if self.insertion and not self.core_complete and len(self.ever_observed_channels) < 16:
                        target = self._target(ch)
                        remaining = self.discovery_stations[self.discovery_station_index:len(self.core_stations)]
                        for a, b in zip(remaining, remaining[1:]):
                            future_best = min(future_best, (distance(a, target) + distance(target, b) - distance(a, b)) / 5.0 + 5.0)
                    if extra <= 160.0 and extra <= future_best + 8.0:
                        ready.append((extra, ch, plan))
            if ready:
                _, ch, plan = min(ready, key=lambda row: (row[0], row[1]))
                self.held_action = action
                self.action_stage = 'early_clear'
                return self._start_plan(ch, plan)
            # A bounded local task starts close to a fresh positive observation.
            eligible = [ch for ch, t in self.tracks.items() if ch not in self.cleared
                        and t.records and self.hunt_versions.get(ch) != len(t.records)
                        and len(self.local_planner.history.get(ch, [])) < self.local_planner.max_probes
                        and min(distance(state.pos, p) for p, _ in t.records) <= 1050.0]
            if eligible and self.early_hunt:
                ch = min(eligible, key=lambda c: (distance(state.pos, self._target(c)), c))
                self.active_hunt = ch
                follow = self._hunt_action(state)
                if follow is not None:
                    self.held_action = action
                    self.action_stage = 'early_clear' if follow.kind == 'clear' else 'early_localization'
                    return follow
        self.action_stage = self.phase
        return action

    def _next_local_action(self, ch, state):
        action = super()._next_local_action(ch, state)
        ch = action.ch
        if self.hunt and ch not in self.plans:
            plan = self._cover(ch, state.pos)
            if plan is not None:
                return self._start_plan(ch, plan)
            if action.kind == 'measure':
                better = self.local_planner.next_probe(self.tracks[ch], state, ch)
                if better is not None:
                    action = better
        if action.kind == 'measure':
            key = (ch, tuple(action.pos), len(self.tracks[ch].records))
            if key in self._issued_local:
                raise RuntimeError('localization_no_progress: repeated fixed-position action')
            self._issued_local.add(key)
        return action

    def on_measure(self, state, ch, result, svd):
        self.actual_scan_points[ch].add(tuple(state.pos))
        if result in ('direction', 'near'):
            self.ever_observed_channels.add(ch)
        super().on_measure(state, ch, result, svd)

    def on_clear(self, state, ch, success):
        if self.early_plan is None or self.early_plan['channel'] != ch:
            return super().on_clear(state, ch, success)
        if success:
            old_index = self.work_index
            FastP4.on_clear(self, state, ch, True)
            if self.phase == 'discovery':
                self.work_index = old_index
            self.early_plan = None
            self.active_hunt = None
            self.plans.pop(ch, None)
            self.hunt_events['early_clear_success'] += 1
        else:
            self.failed_clears.setdefault(ch, []).append(tuple(state.pos))
            self.tracks[ch].clear_failures += 1
            self.tracks[ch].recovery_required = False
            self.early_plan['points'].pop(0)
            self.hunt_events['early_clear_miss_advance'] += 1
            if not self.early_plan['points']:
                raise RuntimeError('early_optical_cover_exhausted: inconsistent observations')

    def diagnostics(self):
        result = super().diagnostics()
        result.update(hunt=self.hunt, early_hunt=self.early_hunt, grid=self.grid, adapt=self.adapt, insertion=self.insertion,
                      mesh_certificate=self.core_certificate,
                      ever_observed_channels=sorted(self.ever_observed_channels),
                      local_planner=self.local_planner.diagnostics())
        result['triggers'].update(self.hunt_events)
        return result
\end{lstlisting}
%</paper-block>

%<paper-block id="appendix.p030" type="paragraph">
\noindent\textbf{文件：}\texttt{02\_\allowbreak{}最终代码/\allowbreak{}第四问/\allowbreak{}strategy\_\allowbreak{}route.py}\\
\textbf{功能：}问题四扫描站与清除任务联合路径规划
%</paper-block>
%<paper-block id="appendix.code028" type="code">
\begin{lstlisting}[style=appendixcode]
"""P4 online joint route for required survey stations and certified clear tasks.

Plans use only current positive observations and actual robot position. Every
unseen channel still receives a complete continuous discovery certificate.
HuntP4 remains unchanged as the deployed baseline.
"""
import math

from strategy import Action
from strategy_fast import FastP4
from strategy_hunt import HuntP4
from strategy_p4 import CHANNELS, distance
from discovery_prune import choose_replacement


def path_length(nodes, start):
    total = 0.0
    for node in nodes:
        total += distance(start, node[2])
        start = node[2]
    return total


def joint_path(stations, tasks, start):
    """Bounded deterministic open TSP heuristic; no scene information is used."""
    nodes = [('scan', i, p) for i, p in enumerate(stations)]
    nodes += [('clear', ch, p) for ch, p in tasks]
    remaining = list(nodes)
    nearest = []
    pos = start
    while remaining:
        nxt = min(remaining, key=lambda n: (distance(pos, n[2]), n[0], n[1]))
        nearest.append(nxt)
        remaining.remove(nxt)
        pos = nxt[2]
    # A scan backbone with cheapest task insertions provides a second initial
    # route. The backbone is itself merely a plan; visits are recorded separately.
    inserted = nodes[:len(stations)]
    for node in nodes[len(stations):]:
        def extra(i):
            a = start if i == 0 else inserted[i - 1][2]
            gain = distance(a, node[2])
            if i < len(inserted):
                gain += distance(node[2], inserted[i][2]) - distance(a, inserted[i][2])
            return gain
        index = min(range(len(inserted) + 1), key=lambda i: (extra(i), i))
        inserted.insert(index, node)

    def improve(route):
        route = route[:]
        for _ in range(64):
            best = None
            for i in range(len(route) - 1):
                before = start if i == 0 else route[i - 1][2]
                for j in range(i + 1, len(route)):
                    gain = distance(before, route[i][2]) - distance(before, route[j][2])
                    if j + 1 < len(route):
                        after = route[j + 1][2]
                        gain += distance(route[j][2], after) - distance(route[i][2], after)
                    if gain > 1e-7 and (best is None or gain > best[0]):
                        best = gain, i, j
            if best is None:
                break
            _, i, j = best
            route[i:j + 1] = reversed(route[i:j + 1])
        # One-node relocations address clusters that segment reversal cannot.
        for _ in range(8):
            best = None
            for i, node in enumerate(route):
                a = start if i == 0 else route[i - 1][2]
                saving = distance(a, node[2])
                if i + 1 < len(route):
                    b = route[i + 1][2]
                    saving += distance(node[2], b) - distance(a, b)
                reduced = route[:i] + route[i + 1:]
                for j in range(len(reduced) + 1):
                    before = start if j == 0 else reduced[j - 1][2]
                    cost = distance(before, node[2])
                    if j < len(reduced):
                        after = reduced[j][2]
                        cost += distance(node[2], after) - distance(before, after)
                    gain = saving - cost
                    if gain > 1e-7 and (best is None or gain > best[0]):
                        best = gain, i, j
            if best is None:
                break
            _, i, j = best
            node = route.pop(i)
            route.insert(j, node)
        return route

    alternatives = [improve(nearest), improve(inserted)]
    return min(alternatives, key=lambda route: path_length(route, start))


class RouteAwareP4(HuntP4):
    name = 'RouteAwareP4_joint21'

    def __init__(self, adaptive_scan=True, scan_investments=0):
        super().__init__(hunt=True, early_hunt=False, grid='compact', adapt=False, insertion=False)
        self.name = type(self).name
        self.pending_stations = list(self.core_stations)
        self.current_station = None
        self.completed_stations = set()
        self.route_history = []
        self._state = None
        self.adaptive_scan = bool(adaptive_scan)
        self.survey_changed = False
        self._consider_scan_stop = False
        self.replacement_history = []
        self.scan_investments = max(0, min(6, int(scan_investments)))
        self.invested_scan_points = []

    def _queued_channels(self):
        channels = [ch for ch in CHANNELS if self._needs_sample(ch)]
        # A measure at the currently selected channel avoids one paid switch.
        # Every other queued channel remains present exactly once.
        if self._state.ch in channels:
            channels.remove(self._state.ch)
            channels.insert(0, self._state.ch)
        return channels

    def _finish_with_evidence(self):
        unseen = set(CHANNELS) - self.ever_observed_channels
        if len(self.ever_observed_channels) < 16:
            actual = set.intersection(*(self.actual_scan_points[ch] for ch in unseen))
            if not set(self.core_stations) <= actual:
                if not self.survey_changed or not self.coverage_certifier.certified_complete(actual):
                    raise RuntimeError('joint_route_missing_actual_discovery_evidence')
        self.core_complete = True
        self.hunt_events['joint_discovery_completed'] += 1
        self._finish_discovery()

    def _route_next(self):
        state = self._state
        if self._consider_scan_stop and self.adaptive_scan and self.pending_stations and len(self.ever_observed_channels) < 16:
            self._consider_scan_stop = False
            unseen = set(CHANNELS) - self.ever_observed_channels
            actual = set.intersection(*(self.actual_scan_points[ch] for ch in unseen))
            replacement = choose_replacement(actual, self.pending_stations, state.pos, state.pos,
                                              checker=self.coverage_certifier)
            self.replacement_history.append(replacement)
            if replacement['removed']:
                self.pending_stations = replacement['remaining']
                self.current_station = tuple(state.pos)
                self._station_channels = self._queued_channels()
                self.discovery_channel_index = 0
                self.survey_changed = True
                self.hunt_events['clear_stop_used_for_discovery'] += 1
                self.hunt_events['proved_scan_station_removed'] += len(replacement['removed'])
                return None
            radius = math.hypot(*state.pos)
            if (len(self.invested_scan_points) < self.scan_investments
                    and 800.0 <= radius <= 1500.0
                    and all(distance(state.pos, p) >= 400.0 for p in actual)):
                # A bounded cost investment at a position already reached for
                # clear. No pending station is removed by this heuristic.
                # Later deletions still require the full continuous certificate.
                self.current_station = tuple(state.pos)
                self._station_channels = self._queued_channels()
                self.discovery_channel_index = 0
                self.invested_scan_points.append(tuple(state.pos))
                self.hunt_events['bounded_additional_scan_at_clear_stop'] += 1
                return None
        ready = {}
        for ch, track in self.tracks.items():
            if ch in self.cleared:
                continue
            if track.near is not None or self._certificate(ch) is not None:
                plan = self._cover(ch, state.pos)
                if plan is not None:
                    ready[ch] = plan
        if len(self.ever_observed_channels) >= 16 and all(ch in self.cleared or ch in ready for ch in self.ever_observed_channels):
            # This derives the public maximum solely from positive observations.
            self.pending_stations.clear()
            self.hunt_events['sixteen_sources_retire_remaining_stations'] += 1
        if not self.pending_stations:
            self._finish_with_evidence()
            return None
        route = joint_path(self.pending_stations, [(ch, self._target(ch)) for ch in ready], state.pos)
        if not route:
            raise RuntimeError('empty_joint_plan_with_pending_stations')
        self.route_history.append(dict(position=state.pos, scan_points=len(self.pending_stations),
                                       ready_channels=sorted(ready), planned_distance_m=path_length(route, state.pos),
                                       next_kind=route[0][0], next_key=route[0][1]))
        kind, key, point = route[0]
        self.hunt_events['joint_route_planned'] += 1
        if kind == 'clear':
            self.hunt_events['joint_clear_before_scan'] += 1
            return self._start_plan(key, ready[key])
        # Reordering happens between complete stations, never mid-channel queue.
        self.pending_stations = [node[2] for node in route if node[0] == 'scan']
        self.current_station = self.pending_stations.pop(0)
        self._station_channels = self._queued_channels()
        self.discovery_channel_index = 0
        return None

    def _next_discovery_action(self):
        while True:
            if self.current_station is not None:
                while self.discovery_channel_index < len(self._station_channels):
                    ch = self._station_channels[self.discovery_channel_index]
                    self.discovery_channel_index += 1
                    if self._needs_sample(ch):
                        return Action('measure', self.current_station, ch)
                self.completed_stations.add(self.current_station)
                self.current_station = None
                self.discovery_station_index += 1
                self._station_channels = []
            action = self._route_next()
            if action is not None:
                return action
            if self.discovery_complete:
                return None

    def step(self, state):
        self._state = state
        if self.done:
            return Action('done')
        if self.early_plan:
            self.action_stage = 'joint_optical_clear'
            return Action('clear', self.early_plan['points'][0], self.early_plan['channel'])
        for ch, track in self.tracks.items():
            if ch not in self.cleared and track.near is not None and distance(state.pos, track.near) < 1e-6:
                self.action_stage = 'near_clear'
                return self._start_plan(ch, self._cover(ch, state.pos))
        # The stable driver contract dispatches our scan and local hooks. Avoid
        # ancestor Hunt's independent greedy clear insertion running a second time.
        action = FastP4.step(self, state)
        self.action_stage = 'joint_optical_clear' if action.kind == 'clear' and self.early_plan else self.phase
        return action

    def diagnostics(self):
        result = super().diagnostics()
        result.update(joint_route=True, route_history=self.route_history,
                      actual_completed_station_count=len(self.completed_stations),
                      adaptive_scan=self.adaptive_scan, replacement_history=self.replacement_history,
                      scan_investments=self.scan_investments, invested_scan_points=self.invested_scan_points)
        return result

    def on_clear(self, state, ch, success):
        super().on_clear(state, ch, success)
        if success and self.phase == 'discovery':
            self._consider_scan_stop = True
\end{lstlisting}
%</paper-block>

%<paper-block id="appendix.p031" type="paragraph">
\noindent\textbf{文件：}\texttt{02\_\allowbreak{}最终代码/\allowbreak{}第四问/\allowbreak{}strategy\_\allowbreak{}route\_\allowbreak{}probe.py}\\
\textbf{功能：}问题四最终 RouteProbeP4：零绕路补测与调度
%</paper-block>
%<paper-block id="appendix.code029" type="code">
\begin{lstlisting}[style=appendixcode]
"""At most two valuable, zero-detour one-bearing probes per known P4 channel.

Parent RouteAwareP4 owns every scan queue and discovery proof. A probe is an
interior point of an already-issued travel segment; its parent action resumes
before any route replanning. Hypothetical target samples only rank information
value. They never prove reception, localization, absence or successful clear.
"""
import math
from collections import Counter

from strategy import Action
from strategy_p4 import distance
from strategy_route import RouteAwareP4


class RouteProbeP4(RouteAwareP4):
    name = 'RouteProbeP4_zero_detour'

    def __init__(self, probe_enabled=True, max_probes_per_channel=2, **kwargs):
        super().__init__(**kwargs)
        if not 0 <= int(max_probes_per_channel) <= 2:
            raise ValueError('max_probes_per_channel must lie within 0..2')
        self.probe_enabled = bool(probe_enabled)
        self.max_probes_per_channel = int(max_probes_per_channel)
        self.probe_positions = {}
        self.probe_events = Counter()
        self.probe_history = []
        self.probe_added_cost = dict(measure_time_s=0.0, switch_time_s=0.0, move_time_s=0.0)
        self._probe_held_action = None
        self._probe_context = None
        self._probe_inflight = False
        self._probe_near_channel = None

    @staticmethod
    def _single_positive(track):
        positions = []
        for position, angle in track.records:
            if not any(distance(position, old[0]) <= 1e-6 for old in positions):
                positions.append((tuple(position), float(angle)))
        return positions[0] if len(positions) == 1 else None

    def _rank_probe(self, ch, track, start, destination, old_channel, scan_channel):
        observation = self._single_positive(track)
        if observation is None:
            return None
        origin, angle = observation
        theta = math.radians(angle)
        bearing = (math.cos(theta), math.sin(theta))
        targets = [(origin[0] + r * bearing[0], origin[1] + r * bearing[1])
                   for r in (500.0, 1000.0, 1400.0)]
        # The public arena only rejects implausible ranking hypotheses.
        targets = [p for p in targets if math.hypot(*p) <= 1800.0001]
        if len(targets) < 2:
            return None
        dx, dy = destination[0] - start[0], destination[1] - start[1]
        segment2 = dx * dx + dy * dy
        if segment2 < 40.0 ** 2:
            return None
        fractions = {0.15, 0.3, 0.5, 0.7, 0.85}
        for point in targets:
            fraction = ((point[0] - start[0]) * dx + (point[1] - start[1]) * dy) / segment2
            if 0.08 <= fraction <= 0.92:
                fractions.add(fraction)
        forbidden = [p for p, _ in track.records] + self.probe_positions.get(ch, [])
        added_switch = float(old_channel != ch) + float(ch != scan_channel) - float(old_channel != scan_channel)
        added_cost = 5.0 + added_switch
        best = None
        for fraction in sorted(fractions):
            point = (start[0] + fraction * dx, start[1] + fraction * dy)
            if any(distance(point, old) <= 1.0 for old in forbidden):
                continue
            estimates = []
            useful = 0
            for target in targets:
                old_range, new_range = distance(origin, target), distance(point, target)
                if not 5.0 < new_range <= 1500.0:
                    estimates.append(0.0)
                    continue
                cross = abs(bearing[0] * (target[1] - point[1]) - bearing[1] * (target[0] - point[0])) / new_range
                if cross < 0.2:
                    estimates.append(0.0)
                    continue
                # Small-angle width is an estimate for value ranking, not the
                # actual region or a clear-radius certificate.
                predicted_radius = math.tan(math.radians(1.0050001)) * (old_range + new_range) / cross
                information = max(0.0, min(1.0, (250.0 - predicted_radius) / 230.0))
                # Saving a future distinct-position visit is worth at most this
                # bounded travel estimate. Reception remains explicitly unknown.
                estimated_saved_s = min(80.0, distance(origin, point) / 5.0) * information * 0.5
                estimates.append(estimated_saved_s)
                useful += predicted_radius <= 120.0
            expected_gain = sum(estimates) / len(targets)
            if useful < 2 or expected_gain <= max(10.0, 2.0 * added_cost):
                continue
            candidate = dict(channel=ch, point=point, segment_fraction=fraction,
                             estimated_gain_s=expected_gain, added_measure_s=5.0,
                             added_switch_s=added_switch, estimated_net_gain_s=expected_gain - added_cost,
                             useful_hypotheses=useful, hypothesis_count=len(targets),
                             reception_guaranteed=False, geometry_certificate=False)
            if best is None or (candidate['estimated_net_gain_s'], -fraction) > (best['estimated_net_gain_s'], -best['segment_fraction']):
                best = candidate
        return best

    def _candidate(self, state, action):
        if (not self.probe_enabled or self.max_probes_per_channel == 0 or self.phase != 'discovery'
                or action.kind != 'measure' or distance(state.pos, action.pos) <= 40.0):
            return None
        candidates = []
        for ch, track in sorted(self.tracks.items()):
            # The held scan channel already has an imminent paid observation;
            # avoid adding another observation to that channel on this trip.
            if (ch in self.cleared or ch == action.ch or track.near is not None
                    or len(self.probe_positions.get(ch, [])) >= self.max_probes_per_channel):
                continue
            candidate = self._rank_probe(ch, track, tuple(state.pos), tuple(action.pos), state.ch, action.ch)
            if candidate is not None:
                candidates.append(candidate)
        return max(candidates, key=lambda row: (row['estimated_net_gain_s'], -row['channel'])) if candidates else None

    def step(self, state):
        self._state = state
        if self._probe_held_action is not None:
            if self._probe_inflight:
                raise RuntimeError('route_probe_measurement_result_missing')
            ch = self._probe_context['channel']
            track = self.tracks.get(ch)
            if (ch not in self.cleared and track is not None and track.near is not None
                    and distance(track.near, state.pos) < 1e-6):
                self._probe_near_channel = ch
                self.action_stage = 'route_probe_near_clear'
                self.probe_events['near_immediate_clear'] += 1
                return self._start_plan(ch, self._cover(ch, state.pos))
            action = self._probe_held_action
            context = self._probe_context
            detour = distance(context['start'], tuple(state.pos)) + distance(tuple(state.pos), action.pos) - context['segment_distance_m']
            if abs(detour) > 1e-6:
                raise RuntimeError('route_probe_resume_would_add_movement')
            self._probe_held_action = None
            self._probe_context = None
            self.action_stage = context['held_stage']
            self.probe_events['held_scan_resumed'] += 1
            return action
        action = super().step(state)
        candidate = self._candidate(state, action)
        if candidate is None:
            return action
        point, ch = candidate['point'], candidate['channel']
        direct = distance(state.pos, action.pos)
        detour = distance(state.pos, point) + distance(point, action.pos) - direct
        assert abs(detour) <= 1e-6, 'probe point must lie on the existing movement segment'
        self._probe_held_action = action
        self._probe_context = dict(candidate, start=tuple(state.pos), destination=tuple(action.pos),
                                   segment_distance_m=direct, added_distance_m=detour,
                                   held_stage=self.action_stage,
                                   prior_consider_scan_stop=self._consider_scan_stop)
        self.probe_history.append(self._probe_context.copy())
        self.probe_positions.setdefault(ch, []).append(point)
        self._probe_inflight = True
        self.probe_events['zero_detour_probe_issued'] += 1
        self.action_stage = 'route_probe'
        return Action('measure', point, ch)

    def on_measure(self, state, ch, result, svd):
        inserted = (self._probe_inflight and self._probe_context is not None
                    and ch == self._probe_context['channel']
                    and distance(state.pos, self._probe_context['point']) < 1e-6)
        super().on_measure(state, ch, result, svd)
        if inserted:
            self._probe_inflight = False
            self.probe_events['probe_' + result] += 1
            self.probe_added_cost['measure_time_s'] += 5.0
            self.probe_added_cost['switch_time_s'] += self._probe_context['added_switch_s']
            self.probe_history[-1]['result'] = result

    def on_clear(self, state, ch, success):
        inserted_near = self._probe_near_channel == ch
        super().on_clear(state, ch, success)
        if inserted_near and success:
            # Resume the already-issued station before the parent can reinterpret
            # this intermediate clear point as a new survey stop.
            self._consider_scan_stop = self._probe_context['prior_consider_scan_stop']
            self._probe_near_channel = None
            self.probe_events['probe_near_clear_success'] += 1

    def diagnostics(self):
        result = super().diagnostics()
        result['route_probe'] = dict(enabled=self.probe_enabled,
                                     max_probes_per_channel=self.max_probes_per_channel,
                                     triggers=dict(self.probe_events), history=self.probe_history,
                                     added_cost=self.probe_added_cost,
                                     added_cost_scope='detection and total round-trip channel-switch increment versus the held scan action; optical success is counted by the driver',
                                     target_samples_are_only_ranking=True)
        result['triggers'].update({'route_probe_' + key: value for key, value in self.probe_events.items()})
        return result
\end{lstlisting}
%</paper-block>


\clearpage
\linespread{1.12}\selectfont
%<paper-block id="appendix.h006" type="heading">
\section*{附录C\quad 补充图件}
%</paper-block>
%<paper-block id="appendix.p032" type="paragraph">
本附录放置正文已引用的局部细节、数值夹逼与辅助结果图。几何示例不表示实际运行轨迹；演练源数为结束后公开信息，正式入口记录只表示已清除数及已记录时间。
%</paper-block>

%<paper-block id="appendix.f001" type="figure">
\begin{figure}[!htbp]
\centering
\paperfigure{0.96}{fig_app_q1_two_station_detail.png}{12}
%<paper-block id="appendix.c001" type="caption">
\caption{仅两次有效测向也可能使直径圆不能覆盖交会区域：右图放大第三顶点的微小径向超出量。}
%</paper-block>
\label{fig:appq1two}
\end{figure}
%</paper-block>

%<paper-block id="appendix.f002" type="figure">
\begin{figure}[!htbp]
\centering
\paperfigure{0.85}{fig_app_q1_error_sensitivity.png}{13}
%<paper-block id="appendix.c002" type="caption">
\caption{固定配置下，允许误差半宽增大时纯交会区域直径的变化；不是对任意中心读数扰动的一般导数上界。}
%</paper-block>
\label{fig:appq1sensitivity}
\end{figure}
%</paper-block>
\clearpage

%<paper-block id="appendix.f003" type="figure">
\begin{figure}[!htbp]
\centering
\paperfigure{0.97}{fig_app_q1_disk_bracket.png}{14}
%<paper-block id="appendix.c003" type="caption">
\caption{圆盘内接、外切多边形与直径夹逼示例。径向误差与最终交集的直径误差分别解释。}
%</paper-block>
\label{fig:appq1disk}
\end{figure}
%</paper-block>

%<paper-block id="appendix.f004" type="figure">
\begin{figure}[!htbp]
\centering
\paperfigure{0.97}{fig_app_q2_continuous_certificate.png}{15}
%<paper-block id="appendix.c004" type="caption">
\caption{第二问对非活动读数部分的63个闭区间核验，活动区间另由解析论证处理；区间拼接覆盖全部读数方向。}
%</paper-block>
\label{fig:appq2cert}
\end{figure}
%</paper-block>
\clearpage

%<paper-block id="appendix.f005" type="figure">
\begin{figure}[!htbp]
\centering
\paperfigure{0.82}{q3_02_practice.pdf}{16}
%<paper-block id="appendix.c005" type="caption">
\caption{第三问8场同版本演练的源数与平均清除时间。P1至P8按正文演练表顺序对应，不同场景不作算法提速配对。}
%</paper-block>
\label{fig:appq3practice}
\end{figure}
%</paper-block>

%<paper-block id="appendix.f006" type="figure">
\begin{figure}[!htbp]
\centering
\paperfigure{0.82}{q4_04_formal_entry.pdf}{17}
%<paper-block id="appendix.c006" type="caption">
\caption{第四问三条正常正式入口记录的每个已清除源平均时间；另两条连接失败记录仍保留在正文表中。}
%</paper-block>
\label{fig:appq4formal}
\end{figure}
%</paper-block>
